\documentclass[10pt,a4paper]{article}

% ----------------- Paquetes -------------------%
\usepackage[T1]{fontenc}
\usepackage[utf8]{inputenc}
\usepackage[a4paper,margin=2.5cm]{geometry}
\usepackage{mathptmx}             
\usepackage{changepage} 
\usepackage{setspace}
\usepackage[spanish]{babel}
\usepackage[labelfont=bf,labelsep=space]{caption}
\usepackage{amssymb}
\usepackage{amsmath}
\usepackage{ebgaramond}
\usepackage{placeins}
\usepackage{etoolbox}
\usepackage{setspace}
\usepackage{parskip} 
\usepackage{tcolorbox}
\usepackage{needspace}
\usepackage{xparse}
\usepackage{ocgx2}
\usepackage{hyperref}
\usepackage{hypcap}


%%%%%%%%%%%%%%%%%%%%%%%%%%%%%%%%%%%%%%%%%%%%%%%%%%%%%%%%%%%%%%%%
% --- Fija primero tus valores globales "buenos" ---
\setstretch{1.2}                 % Interlineado global
\setlength{\parskip}{0.7\baselineskip}
\setlength{\parindent}{0pt}
\newlength{\GlobalParskip}
\setlength{\GlobalParskip}{\parskip}

\newcounter{eqnum}[subsection]
\renewcommand{\theeqnum}{\arabic{eqnum}}
\newcounter{alphaitem}
\newcounter{numitem}

%% ------------------- Recuadros----------------------%%
\newcommand{\puntosclave}[1]{%
  \begin{tcolorbox}[
    colframe=black,
    title={Puntos clave},
    before upper={\setlength{\parindent}{0.5cm}\setlength{\parskip}{0.5\baselineskip}}
  ]
  #1
  \end{tcolorbox}
}

\newcommand{\modificaciones}[1]{
  \begin{tcolorbox}[
    colback=white,
    colframe=red,
    title={Modificaciones a realizar},
    before upper={\setlength{\parindent}{0.5cm}\setlength{\parskip}{0.5\baselineskip}}
  ]
  #1
  \end{tcolorbox}
}

%% ------------------- Listas----------------------%%
    % --- Lista alfabética ---
    \newenvironment{la}
      {\begin{list}{\alph{alphaitem})}{%
          \usecounter{alphaitem}%
          \setlength{\leftmargin}{1cm}%
          \setlength{\parskip}{3pt}% 
          \setlength{\itemsep}{0pt}% ← espacio entre ítems
          \setlength{\parsep}{0pt}%  ← párrafos dentro de un ítem
      }}
      {\end{list}}
      
    % --- Lista numérica ---
    \newenvironment{lnum}
      {\begin{list}{\arabic{numitem}.}{%
          \usecounter{numitem}%
          \setlength{\leftmargin}{1cm}%
          \setlength{\parskip}{3pt}% 
          \setlength{\itemsep}{0pt}% ← espacio entre ítems
          \setlength{\parsep}{0pt}%  ← párrafos dentro de un ítem
      }}
      {\end{list}}
      
% ------------------- Imágenes----------------------%
    \graphicspath{{Img/}}% Rutas por defecto 
    % --- Imagen adaptativa ---
    % Registros de dimensión definidos una sola vez
    \newdimen\imgcap
    \newdimen\imgmin
    \newdimen\imgmax
    \newdimen\imgremain
    \newdimen\imgh
    \newdimen\imgneed
    
    \newcommand{\imagenfit}[3]{%
      \addvspace{10pt}% separación antes
      \par\begingroup
        % Parámetros de composición (asignaciones, no \newdimen)
        \imgcap=1\baselineskip            % alto estimado de título+fuente
        \imgmin=0.15\textheight
        \imgmax=0.15\textheight
        \imgremain=\pagegoal
        \ifdim\pagegoal=\maxdimen \imgremain=0pt\fi
        \advance\imgremain by -\pagetotal
        \advance\imgremain by -\imgcap
        % Decisión de altura destino
        \ifdim\imgremain<\imgmin
          \newpage
          \imgh=0.20\textheight
        \else
          \imgh=\imgremain
          \ifdim\imgh>\imgmax \imgh=\imgmax \fi
        \fi
        % Reservar espacio para evitar cortes del bloque completo
        \imgneed=\imgh \advance\imgneed by \imgcap
        \Needspace{\imgneed}
        % Cuerpo
        \noindent\begin{minipage}{\textwidth}
          \centering
          \captionof{figure}{\textemdash\ #2}\par\vspace{0.3em}% título arriba
          \includegraphics[width=\linewidth,height=\imgh,keepaspectratio]{\detokenize{#1}}\par
          \vspace{0.3em}\small\textit{Fuente: }#3% fuente abajo
        \end{minipage}%
      \endgroup
      \par\addvspace{10pt}%
    }

%------------ Bloque de RECUADRO ESPECIAL -------------%
    \newenvironment{specialblock}[1]{%
      \begin{quote} % crea sangría a izquierda y derecha
      \small % texto más pequeño
      \noindent\textbf{#1}\\[-0.3em] % título en negrita
      \rule{\linewidth}{0.4pt}\\[0.6em]}{
      \end{quote}}
      

%-------------------------- Author Font -----------------------%
    \newcommand{\authorfont}[1]{{\fontfamily{EBGaramond-TLF}\selectfont #1}}

%-------------------------- Anexos -----------------------%
    \makeatletter
    \newcounter{anexo}
    \renewcommand{\theanexo}{\Roman{anexo}}
    \newcommand{\anexo}[1]{%
      \clearpage
      \phantomsection
      \refstepcounter{anexo}%
      \noindent\textbf{Anexo \theanexo.- }#1\par\medskip
      \addcontentsline{stc}{section}{Anexo \theanexo.- #1}%
    }
    \makeatother

    %-------------------------- Lista de figuras (personal) -----------------------%
\makeatletter
\newcommand{\MyListOfFigures}{%
  \clearpage
  \phantomsection
  {\centering\bfseries\Large Índice de figuras\par}%
  \medskip
  % Entrada en nuestro índice de contenido (stc)
  \addcontentsline{stc}{section}{Índice de figuras}%
  % Recuperación del listado (sin cabecera propia)
  \begingroup
    \parindent=0pt\parskip=2pt
    \@starttoc{lof}%
  \endgroup
}
\makeatother

%-------------------------- Bibliografía -----------------------%
\makeatletter
% Contador para las referencias
\newcounter{myref}
% Cabecera de la bibliografía, entra en el índice 'stc'
\newcommand{\MyBibliography}{%
  \clearpage
  \phantomsection
  {\centering\bfseries\Large Bibliografía y referencias\par}%
  \medskip
  \addcontentsline{stc}{section}{Bibliografía y referencias}%
  \setcounter{myref}{0}%
}
\newcommand{\MyBibItem}[4]{%
  \refstepcounter{myref}%
  \noindent[\themyref]\ #1.\ \emph{#2}. #3. #4\par
}
\makeatother

%--------- Establecer cuatro niveles de índice --------%
    \setcounter{tocdepth}{4}   
    \setcounter{secnumdepth}{4}% numera hasta \paragraph

%%%%%%%%%%%%%%%%%%%%%%%%%%%%%%%%

\makeatletter

    %--------- Interlineado Global --------%
    \edef\GlobalBaseLineStretch{\baselinestretch}
    
    %--------- Espaciado de seccinones --------%
    \renewcommand\section{\@startsection{section}
      {1}{\z@}
      {2.5ex \@plus 1ex \@minus .2ex} % espacio antes
      {1.5ex \@plus .2ex}             % espacio después
      {\normalfont\Large\bfseries}    % formato del título
    }
    \renewcommand\subsection{\@startsection{subsection}{2}{\z@}
      {3ex \@plus 0.8ex \@minus .2ex}
      {1.5ex \@plus .2ex}
      {\normalfont\large\bfseries}
    }
    \renewcommand\subsubsection{\@startsection{subsubsection}{3}{\z@}
      {3ex \@plus 0.5ex \@minus .2ex}
      {1ex \@plus .2ex}
      {\normalfont\normalsize\bfseries}
    }
    \renewcommand\paragraph{\@startsection{paragraph}{4}{\z@}%
    {-2ex \@plus -0.5ex \@minus -.2ex}% espacio antes
    {0.8ex \@plus .2ex}% espacio después
    {\normalfont\normalsize\itshape}}% ← cursiva en lugar de negrita

    %--------- Ecuaciones --------%   
    % Variables globales para almacenar títulos y notas
    \gdef\leq@current@section{}%
    \gdef\leq@current@subsection{}%
    \gdef\leq@last@printed@section{}%
    \gdef\leq@last@printed@subsection{}%
    
    % Comando para añadir nota (invisible en el cuerpo)
    \newcommand{\eqnote}[1]{%
      \addcontentsline{leq}{leqnote}{#1}%
    }
    
    % Comando para ecuaciones
    \newcommand{\eqblock}[2]{%
      % Verificar si necesitamos imprimir el título de sección
      \ifx\leq@current@section\leq@last@printed@section\else
        \addcontentsline{leq}{leqsection}{\leq@current@section}%
        \global\let\leq@last@printed@section\leq@current@section
        \gdef\leq@last@printed@subsection{}% Reset subsección al cambiar sección
      \fi
      % Verificar si necesitamos imprimir el título de subsección
      \ifx\leq@current@subsection\leq@last@printed@subsection\else
        \ifx\leq@current@subsection\@empty\else
          \addcontentsline{leq}{leqsubsection}{\leq@current@subsection}%
          \global\let\leq@last@printed@subsection\leq@current@subsection
        \fi
      \fi
      % Ecuación
      \refstepcounter{eqnum}%
      \begin{equation*}
        #1\tag{E.\theeqnum}
      \end{equation*}%
      \addcontentsline{leq}{leqt}{[E.\theeqnum] -- #2}%
      \addcontentsline{leq}{leqm}{#1}%
    }
    
    %%% ----- Índice de ecuaciones ----------%%%
    % Tamaño para []
    \newcommand{\leqIndexFont}{\normalsize}
    \newcommand{\leqTitleFont}{\tiny}
    \newcommand{\leqNoteFont}{\tiny}
    % Cabecera de sección 
    \newcommand*\l@leqsection[2]{%
      \addvspace{25pt}% Mayor separación antes de nueva sección
      \noindent\leqIndexFont
      \makebox[\linewidth]{\rule{.38\linewidth}{0.4pt}\ \ \textbf{  \MakeUppercase{#1}}\ \ \rule{.38\linewidth}{0.4pt}}%
      \par\vspace{-10pt}
    }
    % Cabecera de subsección 
    \newcommand*\l@leqsubsection[2]{%
      \par\addvspace{15pt}%
      \noindent\leqIndexFont #1
      \par\vspace{-7pt}
    }
    % Título de ecuación 
    \newcommand*\l@leqt[2]{%
    \par\addvspace{10pt}%
      \par\noindent\leqTitleFont #1\par
    }
    % Miniatura de ecuación
    \newcommand*\l@leqm[2]{%
      \vspace{0.3\baselineskip}%
      \par\scriptsize\noindent\hspace*{1.5em}\(#1\)\par%
    }
    % Nota de ecuación 
    \newcommand*\l@leqnote[2]{%
      \vspace{0.2\baselineskip}%
      \par\noindent\leqNoteFont\hspace*{1.5em}\textbf{Nota.--} #1\par\vspace{0.5\baselineskip}%
    }
    % Generación del índice
    \newcommand{\mylistofequations}{%
      \clearpage
      \phantomsection
      % Título visual de la lista (ajústalo a tu gusto):
      {\centering\bfseries\Large Índice de ecuaciones\par}
      % Entrada en nuestro índice 'stc':
      \addcontentsline{stc}{section}{Índice de ecuaciones}%
      % Llamada a tu lista real (usa el nombre exacto de tu macro):
      \@starttoc{leq}%
    }
    
    %--------- Captura de títulos de sección --------%
    \let\leq@original@section\section
    \let\leq@original@subsection\subsection
    % Flag para saber si estamos en una sección asterisco
    \newif\ifLeqStarSection
    \LeqStarSectionfalse
    
    % Redefinir \section CON CONTROL DE ASTERISCO
    \renewcommand{\section}{%
      \@ifstar{\leq@section@star}{\@dblarg\leq@section@nostar}%
    }
    \def\leq@section@star#1{%
      \global\LeqStarSectiontrue
      \gdef\leq@current@section{#1}%
      \gdef\leq@current@subsection{}%
      \leq@original@section*{#1}%
      \global\LeqStarSectionfalse
    }
    \def\leq@section@nostar[#1]#2{%
      \global\LeqStarSectionfalse
      \gdef\leq@current@section{#2}%
      \gdef\leq@current@subsection{}%
      \leq@original@section[#1]{#2}%
    }
    % Redefinir \subsection
    \renewcommand{\subsection}{%
      \@ifstar{\leq@subsection@star}{\@dblarg\leq@subsection@nostar}%
    }
    \def\leq@subsection@star#1{%
      \global\LeqStarSectiontrue
      \gdef\leq@current@subsection{#1}%
      \leq@original@subsection*{#1}%
      \global\LeqStarSectionfalse
    }
    \def\leq@subsection@nostar[#1]#2{%
      \global\LeqStarSectionfalse
      \gdef\leq@current@subsection{#2}%
      \leq@original@subsection[#1]{#2}%
    }

    % ----- Restauración de valores Globales de estilo ----------%
    % Flags
    \newif\ifInFootnote
    \InFootnotefalse
    \pretocmd{\@footnotetext}{\InFootnotetrue}{}{}
    \apptocmd{\@footnotetext}{\InFootnotefalse}{}{}
    % Profundidad de listas (soporta anidadas)
    \newcount\MyListDepth \MyListDepth=0
    \AtBeginEnvironment{la}{\advance\MyListDepth by 1}
    \AfterEndEnvironment{la}{\advance\MyListDepth by -1}
    \AtBeginEnvironment{lnum}{\advance\MyListDepth by 1}
    \AfterEndEnvironment{lnum}{\advance\MyListDepth by -1}
    \AtBeginEnvironment{itemize}{\advance\MyListDepth by 1}
    \AfterEndEnvironment{itemize}{\advance\MyListDepth by -1}
    \AtBeginEnvironment{enumerate}{\advance\MyListDepth by 1}
    \AfterEndEnvironment{enumerate}{\advance\MyListDepth by -1}
    % Restauración diferida: hazlo al INICIO del próximo párrafo real
    \newif\if@pendingRestore
    \@pendingRestorefalse
    \newcommand{\RequestRestore}{\global\@pendingRestoretrue}
    \everypar{%
      \if@pendingRestore
        \global\@pendingRestorefalse
        \ifInFootnote\else
          \ifnum\MyListDepth=0
            \setlength{\parskip}{\GlobalParskip}%
            \setstretch{\GlobalBaseLineStretch}%
          \fi
        \fi
      \fi
    }
    % Al final de bloques:
    \AfterEndEnvironment{equation}{\RequestRestore}
    \AfterEndEnvironment{equation*}{\RequestRestore}
    \AfterEndEnvironment{la}{\RequestRestore}
    \AfterEndEnvironment{lnum}{\RequestRestore}
    % Al final de macros:
    \apptocmd{\eqblock}{\RequestRestore}{}{}
    \apptocmd{\imagenfit}{\RequestRestore}{}{}
    
\makeatother

% ===================== ToC propio con 4 niveles (único bloque) =====================
\makeatletter

% --- Escritores: añadimos una línea al .stc en cada cabecera numerada ---
% Guardamos las versiones actuales para llamar al comportamiento normal
\let\MyTOC@orig@section\section
\let\MyTOC@orig@subsection\subsection
\let\MyTOC@orig@subsubsection\subsubsection
\let\MyTOC@orig@paragraph\paragraph

% SECTION
\renewcommand{\section}{%
  \@ifstar{\MyTOC@section@star}{\@dblarg\MyTOC@section@nostar}%
}
\def\MyTOC@section@star#1{%
  \MyTOC@orig@section*{#1}% sin entrada al .stc
}
\def\MyTOC@section@nostar[#1]#2{%
  \MyTOC@orig@section[{#1}]{#2}% comportamiento normal (escribe al .toc)
  % Escribimos también nuestra línea al .stc (texto con \numberline)
  \addcontentsline{stc}{section}{\protect\numberline{\thesection}#2}%
}

% SUBSECTION
\renewcommand{\subsection}{%
  \@ifstar{\MyTOC@subsection@star}{\@dblarg\MyTOC@subsection@nostar}%
}
\def\MyTOC@subsection@star#1{%
  \MyTOC@orig@subsection*{#1}%
}
\def\MyTOC@subsection@nostar[#1]#2{%
  \MyTOC@orig@subsection[{#1}]{#2}%
  \addcontentsline{stc}{subsection}{\protect\numberline{\thesubsection}#2}%
}

% SUBSUBSECTION
\renewcommand{\subsubsection}{%
  \@ifstar{\MyTOC@subsubsection@star}{\@dblarg\MyTOC@subsubsection@nostar}%
}
\def\MyTOC@subsubsection@star#1{%
  \MyTOC@orig@subsubsection*{#1}%
}
\def\MyTOC@subsubsection@nostar[#1]#2{%
  \MyTOC@orig@subsubsection[{#1}]{#2}%
  \addcontentsline{stc}{subsubsection}{\protect\numberline{\thesubsubsection}#2}%
}

% PARAGRAPH
\renewcommand{\paragraph}{%
  \@ifstar{\MyTOC@paragraph@star}{\@dblarg\MyTOC@paragraph@nostar}%
}
\def\MyTOC@paragraph@star#1{%
  \MyTOC@orig@paragraph*{#1}%
}
\def\MyTOC@paragraph@nostar[#1]#2{%
  \MyTOC@orig@paragraph[{#1}]{#2}%
  % Si no numeras los \paragraph, cambia \theparagraph por texto plano sin \numberline
  \addcontentsline{stc}{paragraph}{\protect\numberline{\theparagraph}#2}%
}

% --- Impresor del índice propio (lee el .stc) ---
\newcommand{\MyTOC}{%
  \par\bigskip
  {\centering\bfseries\Large Índice de contenido\par}%
  \medskip
  \begingroup
    \parindent=0pt\parskip=2pt
    \@starttoc{stc}%
  \endgroup
}

\makeatother
% ===================== fin ToC propio =====================


%%%%%%%%%%%%%%


% --- Datos del tema ---
\title{Teoría impositiva (VIII): Teoría del federalismo fiscal\\
\large Asignación óptima de las funciones entre los distintos niveles de Gobierno para alcanzar el bienestar máximo. La financiación de las haciendas territoriales: impuestos, transferencias y deuda}
\author{Marcelo Ortega}
\date{Spetiembre 2026}

\begin{document}
\maketitle
\newpage

% --- Página de resumen y modificaciones ---
\puntosclave{ %% Escribir puntos clave Apuntes CECO
     Para la profundización en la materia, se recomienda centrarse en la investigación empírica, parte de la cual se puede encontrar en los Anexos del Tema de ICEX-CECO.
    }
\vspace{1cm}

\modificaciones{
    
    }

\newpage
\MyTOC
\newpage

\mylistofequations
\clearpage

\MyListOfFigures
\clearpage


\pagenumbering{arabic}
\setcounter{page}{1}


\section{Introducción}



\subsection{Contextualización}


\subsubsection{Relevancia}




\subsubsection{Problemática}



\subsection{Estructura de la exposición}










\clearpage
\section{Dimensión política: Marco coneptual}
\puntosclave{
    
}

\subsection{ICEX-CECO}
El federalismo fiscal es la teoría que estudia la deseabilidad de la provisión de bienes públicos (y de su financiación) a un determinado nivel regional. Típicamente se refiere a crear subunidades administrativas dentro de un mismo país, pero también puede aplicarse a la creación de federaciones internacionales.\footnote{Por ejemplo, puede servir como marco de análisis tanto para estudiar si es deseable que servicios como la sanidad pública o la educación se provean por el Estado central o por las Comunidades Autónomas, como también para estudiar si son deseables la formación de un ejército o de un sistema de infraestructuras a nivel de la Unión.

Los principales casos de estudio en la literatura han sido los Estados Unidos, la Unión Soviética, y más recientemente la Unión Europea.
}

De lo anterior resultan dos cuestiones fundamentales a tener en cuenta cuando se trata el Tema. En primer lugar, la problemática es similar y complementaria a la de la Teoría de la integración monetaria. Sin embargo. mientras que ésta última se enmarca dentro de la macroeconomía, el federalismo fiscal pertenece a la Economia Pública. En ambos casos. el análisis se enriquece fuertemente complementándolo con una perspectiva de Economía Política. En segundo lugar. se trata de una cuestión marcadamente política, dado que tiene que ver con cuestiones estudiadas típicamente por la Ciencia Política como el papel del Estado, su organización territorial u otros elementos definitorios, y la soberanía. Como consecuencia de esto, no es sorprendente comprobar que existen numerosas teorías y opiniones sobre el federalismo fiscal.\footnote{Dado que la mayoría de la literatura sobre la materia procede de EEUU, donde las posiciones conservadoras han defendido históricamente una mayor descentralización, mayor atrihución de competencias para los estados frente al Gobierno federal, no sorprende que en este campo abunden trabajos de economistas ideológicamente conservadores que defiendan la descentralización (en especial, los de la Public Choice School)
} Si bien algunos trabajos sí son centrales en la literatura, se trnta en general de un campo que se caracteriza por la atomicidad de aportaciones por economistas con visiones y métodos muy diferentes.

El presente tema busca recopilar esta heterogeneidad a través de una estructura clara, en la que se opta por no adentrarse en profundidad en modelos concretos sino en enumerar argumentos y tratar de sintetizarlos en fomrn de resultados amplios.

Se busca que el opositor comprenda algunos aspectos. Primero, cuál es la motivación o por qué se justifican sistemas descentralizados. Segundo, debe comprenderse bien cuáles son las distintas formas de financiación a las jurisdicciones, y qué implicaciones tienen en la vehiculización del gasto. Tercero, es fw1damental que el opositor comprenda los trade-off.s· que conlleva un sistema descentralizado y por tanto la ausencia de un resultado úico y generalizable a todos los casos: en Economía, no se puede concluir que la descentralización sea mejor o peor que un sistema descentralizado para cualquier jurisdicción (probablemente. ni siquiera para una jurisdicción concreta).

Finalmente, y como apunte muy importante, debemos dividir el estudio en dos líneas de investigación diferentes: las aportaciones de la primera generación de la Teoría del Federalismo Fiscal, y las de la segunda. Mientras que la primera (hasta los años 70) adopta un enfoque teórico-normativo, la segunda (desde ese momento hasta hoy) se fundamenta en modelos de Economía Política. Se remite al opositor a la nota previa al tema 4B2 para una explicación de conjunto de esa rama de la literatura y su tratamiento en el temario de nuestra oposición.




\tetxcolor{yellos}{\textbf{Descentralización por el lado del ingreso}}
III.II.I. Implicaciones de la descentralización de los ingresos en los objetivos del gobierno central

En este sentido, NISKANEN indicaba que el Sector Público se caracteriza por tratar de maximizar sus presupuestos para así aumentar su poder e influencia. aumentar el número de burócratas. y sus salarios. BUCHANAN en 1980 extendió esta idea al considerar que el Sector Público puede considerarse como un ·'Leviatán'' que se extiende maximizando las rentas que extrae de la economía. aumentando también su poder monopolístico.

Como se ·mostró en el análisis de la descentralización del gasto. BRENNAN y BUCHANAN concluían en 1 980, que el sistema descentralizado podía aumentar la rendición de cuentas de cara a satisfacer las preferencias de los agentes. Esta conclusión tiene también importantes implicaciones por el lado de los ingresos. En este sentido, estos autores afirman que el sistema descentralizado sirve de mecanismo para reducir el control del gobierno central sobre los recursos de la economía. En esta línea, resaltan el importante papel que tiene la competencia entre distintos gobiernos en un sistema descentralizado. Así. cuando se tienen agentes y empresas móviles. el sistema actúa corno un sustituto parcial o total de la imposición de restricciones fiscales e:-.plícitas al poder impositi,·o de los gobiernos.

Esta idea del Leviatán ha sido analizada por muchos autores de la segunda generación de la Teoría del Federalismo Fiscal, existiendo una gran diversidad de opiniones. Sin embargo. como Roddcn en 2003 <;eñaló. el sistema descentralizado efectivamente puede ser una solución a problemas del Sector Publico. aunque lo importante es la fomrn que toma esta de::.centralización.

De esta manera. Roddcn afinna que en los sistemas descentralizados en los que las jurisdicciones únicamente dependen de sus propios tributos sí se observa un menor tamaño del Sector Público. Sin embargo, señala que si la fuente de financiación principal de los gobiernos locales se debe alas transferencias, pueden surgir incentivos perversos que acaben dando lugar a un Sector Público de mayor tamaño.

Existen otras teorías que analizan la forma que debe tomar el sistema descentralizado en cuanto a la financiación de las haciendas territoriales. Podemos mencionar aquí el trabajo de Levaggi de 2002, basado en el problema principal-agente.

LEVAGGI considera que el gobierno central es el principal, y que los gobiernos locales son agentes. Supone que el principal actúa de tal manera que trata de influir en las juiisdicciones para que cumplan los objetivos centrales. Así. el gobierno central solo puede controlar aspectos fiscales. y actúa con información asimétrica.

De esta manera. el autor seíiala dos importantes conclusiones para el gobierno local. En primer lugar, de cara a l ograr sus objetivos, el gobierno central debe limitar la capacidad de decisión de las j urisdicciones en cuanto al total de presupuesto y su distribución en distintos programas.

Además. Levaggi concluye que no se debería establecer un sistema de transferencias. Vemos, por tanto. cómo existen distintas ópticas al introducir aspectos relevantes de la Teoría de la Elección Pública a la Teoría del Federalismo Fiscal. ya que existen autores que consideran que la descentralización puede suponer una solución al problema del Leviatán. pero otros analizan este problema que se da en los gobiernos locales.

En definitiva, la forma que toma el sistema descentralizado tiene un papel fundamental en el funcionamiento de éste y en el equilibiio resultante.

\textbf{IV. Conclusión}

Dentro del análisis realizado sobre el sistema descentralizado se han resaltado las diferencias fundamentales existentes en cuanto a los enfoques de las dos generaciones de la Teoría del Federalismo Fiscal.

Las aportaciones de la primera generación justificaban la descentralización de las competencias para lograr una mejor provisión de bienes públicos locales y con esto satisfacer en mayor medida las preferencias de los agentes, alcanzando así un m ayor nivel de bienestar social. Analizaban para esto cuál era el tamaño óptimo de la jurisdicción. y las distintas fuentes de financiación y sus implicaciones de cara a lograr estos objetivos. 

Sin embargo. los autores de la segunda generación tienen por objeto anal izar qué factores afectan al buen funcionamiento del sistema, afectados originalmente por la Escuela de la Elección Pública e incorporando al análisis cuestiones de Teoría de Juegos (comportamientos estratégicos. asimetrías de información. y fallos endógenos del sistema). Cabe señalar que las conclusiones de estas aportaciones son muy dispares, existiendo un amplio abanico en cuanto a la modelización para el análisis de distintas líneas de investigación.

A modo conclusión es importante destacar algunos aprendizajes. En primer lugar, el hecho de que el sistema multinivel tiene tanto ventajas como desventajas. Es importante. por tanto, ponderar los beneficios que supone una mejora en la eficiencia en la asignación. en comparación a la pérdida que puede suponer no internalizar externalidades o no aprovechar economías de escala. En especial, no hay resultados generales extrapolables a toda jurisdicción.

Pero otra de las principales conclusiones del análisis. es que la forma o estructura que tome el sistema descentralizado es un aspecto que tiene una gran influencia en su funcionamiento. El policy maker debe tener en cuenta las potenciales interacciones entre unidades de gobierno, y las reglas que las definen de cara a diseñar el sistema, ya que afectará al resultado. 

El estudio del federalismo fiscal cobra un gran interés en el contexto actual. ya que un gran número de economías han establecido el sistema multinivel, ya sea a nivel nacional, o supranacional. Podemos decir que uno ele los principales retos a los que se enfrentan estos sistemas descentralizados en el contexto actual, se refiere a la financiación y al endeudamiento de las diferentes unidades económicas. También es importante mencionar como otro de los retos principales la propia sostenibilidad y estabilidad del sistema. que, en muchos casos a nivel mundial, se ve amenazada.

En la actualidad, la investigación está centrada principalmente en los fallos del sistema, y se avanza en el enfoque de la segunda generación de la Teoría del Federalismo Fiscal en un m arco tanto teórico como empírico.

No obstante. en la práctica se observa que muchas estructuras descentralizadas se han construido en base a criterios históricos y políticos. ya que ésta5 o cumplen con los postulados analizados en el cuerpo del tema.







\subsection{GRUBER (2017)}
\textcolor{blue}{\textbf{OJO -> Origen:} GRUBER (2017)}

Conclusión

En todos los países, el gobierno central recauda solo una parte de los ingresos fiscales nacionales totales y realiza solo una parte del gasto público nacional. El resto de la imposición y del gasto es realizado por los gobiernos subnacionales, como los gobiernos estatales y locales en Estados Unidos. En relación con otros países desarrollados, Estados Unidos asigna una gran proporción de las responsabilidades gubernamentales a sus gobiernos subnacionales. Este capítulo presentó una teoría para explicar por qué el gasto podría dividirse entre los gobiernos nacionales y subnacionales.

Cuando el gasto se destina a bienes para los cuales las preferencias locales son relativamente similares, y cuando la mayoría de los residentes pueden beneficiarse de esos bienes, el modelo de Tiebout sugiere que el gasto debería realizarse localmente. Cuando el gasto se destina a bienes que benefician únicamente a una minoría de la población, como la redistribución de la renta, el modelo de Tiebout sugiere que podría ser difícil realizar este gasto localmente porque la mayoría de las personas que no se benefician «votarán con los pies» y se trasladarán a otro lugar. Estos resultados son coherentes con la división de responsabilidades respecto al gasto en educación y seguridad pública (local) y redistribución (nacional). Además, si el gasto tiene efectos externos sobre otras comunidades, la provisión local también puede ser ineficiente, lo cual es coherente con el hecho de que la financiación de la educación en Estados Unidos sea compartida entre los gobiernos locales y estatales, aunque plantea la cuestión de si el gobierno federal debería desempeñar un papel mayor.

Los niveles superiores de gobierno pueden no creer en las conclusiones del modelo idealizado de Tiebout, en cuyo caso querrán redistribuir entre los niveles inferiores de gobierno. Si el gobierno de nivel superior decide que quiere redistribuir entre los niveles inferiores, puede hacerlo mediante varios tipos diferentes de subvenciones. La elección apropiada depende del objetivo (redistribuir para compensar los fallos de Tiebout o redistribuir para compensar las externalidades).




\subsection{ROSEN y GARBER}
\textcolor{red}{\textbf{OJO -> Origen:} ROSEN y GARBER}

\textcolor{red}{\textbf{OJO -> Origen:} ROSEN y GARBER}
\textbf{Federalismo óptimo}

Ahora que tenemos una idea de cómo caracterizar los gobiernos locales, volvemos a nuestra pregunta anterior. ¿Cuál es la asignación óptima de responsabilidades económicas entre los niveles de gobierno en un sistema federal? Consideremos primero brevemente las funciones macroeconómicas. La mayoría de los economistas están de acuerdo en que las decisiones de gasto e imposición destinadas a afectar a los niveles de desempleo e inflación deberían ser tomadas por el gobierno central. Ningún gobierno estatal o local es lo suficientemente grande como para afectar al nivel general de actividad económica. No tendría sentido, por ejemplo, que cada localidad emitiera su propia oferta monetaria y siguiera una política monetaria independiente.

Con respecto a las actividades microeconómicas de mejorar la eficiencia y la equidad, existe considerablemente más controversia. Planteada dentro del marco de la economía del bienestar, la cuestión es si un sistema centralizado o descentralizado tiene más probabilidades de maximizar el bienestar social. Para simplificar, la mayor parte de nuestro análisis supone únicamente dos niveles de gobierno, central y local. Con este supuesto no se pierde ninguna idea importante.

\textbf{Desventajas de un sistema descentralizado}

Consideremos un país compuesto por un grupo de pequeñas comunidades. El gobierno de cada comunidad toma decisiones para maximizar una función de bienestar social que depende únicamente de las utilidades de sus miembros; los de fuera no cuentan. ¿Cómo se comparan los resultados con los que surgirían de maximizar una función de bienestar social nacional que tuviera en cuenta las utilidades de todos los ciudadanos? Consideramos primero las cuestiones de eficiencia y después las de equidad.

\textbf{Cuestiones de eficiencia}
Un sistema de gobiernos descentralizados podría conducir a una asignación ineficiente de los recursos por varias razones.

\textbf{Externalidades}
Un bien público cuyos beneficios corresponden únicamente a los miembros de una comunidad concreta se denomina bien público local. Por ejemplo, la biblioteca pública de Austin, Texas, tiene poco efecto sobre el bienestar de las personas de Ann Arbor, Michigan. Sin embargo, las actividades emprendidas por una comunidad pueden afectar en ocasiones al bienestar de las personas de otras comunidades. Si un municipio proporciona una buena educación pública a sus niños y algunos de ellos finalmente emigran, otras comunidades pueden beneficiarse de tener una fuerza de trabajo mejor educada. Los municipios también pueden afectarse negativamente entre sí. Victoria, Columbia Británica, vierte sus aguas residuales sin tratar al mar; parte de los residuos llega hasta Seattle, Washington, a cuyos ciudadanos no les gusta en absoluto. En resumen, las comunidades se imponen externalidades (tanto positivas como negativas) unas a otras.

Si cada comunidad se preocupa únicamente por sus propios miembros, estas externalidades se pasan por alto. Por tanto, de acuerdo con el argumento estándar, los recursos se asignan de manera ineficiente.

\textbf{Economías de escala en la provisión de bienes públicos}
Para determinados servicios públicos, el coste por persona disminuye a medida que aumenta el número de usuarios. Por ejemplo, cuantas más personas utilizan una biblioteca pública, menor es el coste por usuario. Si cada comunidad establece su propia biblioteca, los costes por usuario son más altos de lo necesario. Una jurisdicción central, por otro lado, podría construir una sola biblioteca, permitiendo que las personas se beneficiaran de las economías de escala.

Por supuesto, distintas actividades están sujetas a diferentes economías de escala. La escala óptima para los servicios bibliotecarios podría diferir de la correspondiente a la protección contra incendios. Y ambas seguramente difieren de la escala óptima para la defensa nacional. Esta observación, por cierto, ayuda a racionalizar un sistema de jurisdicciones superpuestas: cada jurisdicción puede encargarse de aquellos servicios con economías de escala que sean apropiadas para el tamaño de la jurisdicción.

Por otro lado, la consolidación no es la única manera de que las comunidades aprovechen las economías de escala. Por ejemplo, la figura supramunicipal de la \textit{mancomunidad}, en España, permite a diferentes municipios gestionar conjuntamente determinados servicios (escueles, bibliotecas, recogida de basuras, etc.), aprovechando las economías de escala y conservando al mismo tiempo su independencia. Estos acuerdos debilitan el vínculo entre las decisiones de la jurisdicción sobre cuánto consumir de un bien proporcionado públicamente y cuánto producir.

\textbf{Sistemas fiscales ineficientes }
En términos generales, una imposición eficiente requiere que los bienes cuya demanda u oferta sean inelásticas sean gravados a tipos relativamente altos, y viceversa. Supongamos que la oferta de capital para todo el país es fija, pero que el capital es altamente móvil entre jurisdicciones subfederales. Cada jurisdicción se da cuenta de que, si impone un impuesto considerable sobre el capital, el capital simplemente se trasladará a otro lugar, empeorando así la situación de la jurisdicción. En tal situación, una jurisdicción racional grava el capital muy ligeramente, o incluso lo subvenciona. Un ejemplo es Connecticut, que desde 2006 ha ofrecido a los productores cinematográficos un crédito fiscal del 30\%. Durante un período de dos años, el crédito fiscal atrajo más de 66 largometrajes, programas de televisión y anuncios (FODERARO, 2008). De forma más general, CHIRINKO y WILSON (2006) encontraron que, durante los últimos 40 años, los incentivos fiscales estatales a la inversión se han vuelto cada vez mayores y cada vez más comunes.

En realidad, por supuesto, la oferta del stock total de capital no es fija. Tampoco se sabe exactamente en qué medida responden las decisiones de localización de las empresas a las diferencias en los tipos impositivos locales, aunque existen algunas pruebas estadísticas de que el crecimiento del empleo en una jurisdicción está inversamente correlacionado con sus tipos impositivos sobre las empresas [Mark et al., 2000]. Pero el punto básico permanece: es improbable que los impuestos recaudados por comunidades descentralizadas sean eficientes desde un punto de vista nacional. En cambio, es probable que las comunidades seleccionen los impuestos sobre la base de si pueden trasladarse a personas de fuera. Por ejemplo, si una comunidad tiene la única mina de carbón del país, esperamos que la incidencia de un impuesto sobre el carbón impuesto localmente recaiga en gran medida sobre los usuarios de carbón de fuera de la comunidad.⁷ Un impuesto sobre el carbón sería una buena idea desde el punto de vista de la comunidad, pero no necesariamente desde el de la nación.⁸

Una implicación importante del traslado de impuestos es que las comunidades pueden comprar demasiados bienes públicos locales. La eficiencia exige que los bienes públicos locales se compren hasta el punto en el que su beneficio social marginal sea igual al coste social marginal. Si las comunidades pueden trasladar parte de la carga a otras jurisdicciones, el coste marginal percibido por la comunidad es menor que el coste social marginal. Cuando las comunidades igualan el beneficio social marginal al coste marginal percibido, el resultado es una cantidad ineficientemente grande de bienes públicos locales.

\textbf{Economías de escala en la recaudación de impuestos}
Las comunidades individuales pueden no ser capaces de aprovechar las economías de escala en la recaudación de impuestos. Cada comunidad tiene que dedicar recursos a la administración tributaria, y pueden obtenerse ahorros teniendo una autoridad tributaria conjunta. ¿Por qué no dividir los costes de un único ordenador para llevar el seguimiento de las declaraciones de impuestos, en lugar de que cada comunidad compre el suyo propio? Por supuesto, algunas de estas economías podrían lograrse simplemente mediante la cooperación entre las jurisdicciones, sin que se produjera una consolidación efectiva. 

\textbf{Cuestiones de equidad} 
Maximizar el bienestar social puede requerir transferencias de renta a los pobres. Supongamos que el patrón de impuestos y gastos de una comunidad concreta es favorable a sus miembros de renta baja. Si no existen barreras al movimiento entre comunidades, esperamos una inmigración de pobres procedentes del resto del país. A medida que aumenta la población pobre, también lo hace el coste de la política fiscal redistributiva. Al mismo tiempo, las personas de renta alta del municipio pueden decidir marcharse. ¿Por qué deberían pagar impuestos elevados para los pobres cuando pueden trasladarse a otra comunidad con una estructura fiscal más ventajosa? Así, las demandas sobre la base impositiva de la comunidad aumentan mientras su tamaño disminuye. Finalmente, el programa redistributivo tiene que ser abandonado.

Este argumento se apoya en gran medida en la idea de que las decisiones de las personas de establecerse en una determinada jurisdicción están influidas por el paquete disponible de impuestos y prestaciones sociales. Existe cierto apoyo anecdótico para esta proposición. En la década de 1990, los legisladores de California estaban lo suficientemente preocupados por la migración hacia su estado inducida por la asistencia social como para restringir a los nuevos migrantes, durante su primer año en el estado, a las prestaciones de asistencia social de los estados desde los que se habían trasladado. Sin embargo, el Tribunal Supremo declaró inconstitucionales dichas leyes en 1999.

Feldstein y Wrobel [1998] proporcionan algunas pruebas en esta línea, señalando que, si los individuos de renta alta pueden evitar condiciones fiscales desfavorables migrando a estados con tipos impositivos más bajos, entonces los empleadores de los estados con impuestos elevados tendrán que pagar salarios antes de impuestos más altos para conservar a sus trabajadores. El efecto neto es que no se produce ningún cambio en la distribución de la renta. Feldstein y Wrobel encuentran que, de hecho, cuando los estados elevan sus tipos impositivos, los salarios antes de impuestos aumentan poco después. La interpretación de este hallazgo es un poco complicada; podría darse el caso de que la causalidad opere en la otra dirección: los estados cuyos ciudadanos han experimentado aumentos salariales votan a favor de sistemas fiscales más progresivos. En cualquier caso, el resultado sugiere que se requiere cautela cuando las jurisdicciones descentralizadas intentan llevar a cabo una redistribución de la renta.


\textcolor{red}{Implicaciones}

El análisis anterior deja claro que no puede esperarse que un sistema puramente descentralizado maximice el bienestar social. La eficiencia exige que los bienes con efectos de desbordamiento que afectan a todo el país —bienes públicos nacionales como la defensa— sean proporcionados a nivel nacional. Por otro lado, los bienes públicos locales deberían proporcionarse localmente. Como expresó Dave Cieslewicz, alcalde de Madison, Wisconsin, cuando la gente de su municipio debatía si adoptar una postura sobre el conflicto entre israelíes y palestinos: \textit{Fui elegido para que se recogiera la basura y se limpiaran las calles de nieve, }[no para]\textit{ actuar sobre cuestiones de política internacional} (NAPOLITANO, 2004).

Esto nos deja con el caso intermedio de las actividades comunitarias que crean efectos de desbordamiento que no tienen alcance nacional. Una posible solución es situar a todas las comunidades que se afectan entre sí bajo un único gobierno regional. En teoría, este gobierno tendría en cuenta el bienestar de todos sus ciudadanos y, por tanto, internalizaría las externalidades. Sin embargo, una jurisdicción gubernamental más grande puede responder menos a las diferencias locales en las preferencias.

Un método alternativo para tratar las externalidades es un sistema de impuestos y subvenciones pigouvianos. Sabemos que el gobierno puede aumentar la eficiencia gravando las actividades que crean externalidades negativas y subvencionando las actividades que crean externalidades positivas ([\authorfont{Ver Tema 3.A.24}]). Podemos imaginar al gobierno central utilizando mecanismos similares para influir en las decisiones de los gobiernos locales. Por ejemplo, si la educación primaria y secundaria crea beneficios que van más allá de los límites de una jurisdicción, el gobierno central puede proporcionar subvenciones educativas a las comunidades. Se mantiene la autonomía local, pero se corrige la externalidad. 

Nuestra teoría sugiere una división bastante clara de la responsabilidad de la provisión de bienes públicos: los bienes públicos locales, por las localidades, y los bienes públicos nacionales, por el gobierno central. En la práctica, existe una interacción considerable entre los niveles de gobierno. Por ejemplo, la mayoría de los agentes encargados de hacer cumplir la ley son funcionarios estatales y locales. Sin embargo, muchas de sus actuaciones están regidas por el derecho penal federal, que \textit{ha crecido explosivamente a medida que el Congreso ha adoptado posturas contra delitos tales como el robo de automóviles con violencia y la quema de iglesias, la interrupción de un rodeo y los daños a una instalación ganadera} (DERTHICK, 2000). Dado que las localidades podrían actuar de manera inapropiada en ausencia de tales regulaciones, su presencia puede mejorar el bienestar. Sin embargo, algunos creen que el sistema de regulación federal sobre las unidades gubernamentales subfederales se ha vuelto tan complicado que resulta difícil determinar qué nivel de gobierno tiene responsabilidad sobre qué. Esto podría ayudar a explicar la respuesta inadecuada al huracán Katrina en 2005: la confusión sobre las funciones que debía desempeñar cada nivel de gobierno condujo a una falta de coordinación que retrasó servicios esenciales.

Si una división de responsabilidades es apropiada desde el punto de vista de la eficiencia, ¿ocurre lo mismo con la distribución de la renta? La mayoría de los economistas creen que las consideraciones de movilidad analizadas anteriormente descartan depender en gran medida de los gobiernos locales para alcanzar objetivos distributivos. Es probable que una jurisdicción individual que intente hacerlo se encuentre con problemas financieros. 

\textbf{La educación pública en un sistema federal}

Una forma útil de aplicar la teoría del federalismo óptimo es utilizarla para analizar la educación, una de las partidas más importantes de los presupuestos de los gobiernos estatales y locales. 

¿Se ajusta este patrón de gasto en educación por parte de los diferentes niveles de gobierno a nuestras ideas sobre el federalismo óptimo? Un argumento a favor de la provisión descentralizada de un bien es que puede adaptarse a las preferencias locales. Debido a que muchos padres tienen opiniones firmes sobre la educación de sus hijos y estas opiniones difieren entre comunidades, el papel principal desempeñado por los gobiernos locales en la provisión de educación tiene sentido. Por supuesto, se podría permitir discreción local sobre la política escolar mientras se proporciona financiación desde los niveles estatal o federal de gobierno. Políticamente, sin embargo, puede ser difícil mantener el control de las escuelas si la financiación procede de algún otro nivel de gobierno: quien paga al músico elige la melodía.\footnote{En California, por ejemplo, una cantidad sustancial de financiación pública procede del gobierno estatal. Las escuelas públicas están sujetas a un código educativo estatal de 9.000 páginas, que les dice qué libros de texto comprar, cómo enseñar fonética y que sus cafeterías deben tener cocinas de servicio completo, entre otras cosas (KRONHOLZ, 2000).
}

Los gobiernos locales recaudan dinero para la educación principalmente mediante la imposición sobre la propiedad; existen amplias variaciones en la cantidad de riqueza inmobiliaria disponible para los distritos escolares. Las variaciones en la base del impuesto sobre la propiedad pueden estar asociadas con enormes diferencias en la financiación de los distritos escolares. En 2006, por ejemplo, entre los distritos escolares de California con al menos 10.000 estudiantes, el gasto por alumno era 11 veces mayor en el distrito más rico que en el distrito más pobre [US Bureau of the Census, 2008d]. Una visión igualitaria del gasto educativo exigiría financiación procedente de un nivel de gobierno que pudiera redistribuir recursos a través de las fronteras locales, independientemente de sus posibles efectos sobre la autonomía local. 

La financiación federal de la educación se concentra en dos áreas: en los niveles de educación primaria y secundaria, la financiación del Departamento de Educación se destina principalmente a programas que atienden a niños desfavorecidos (14.700 millones de dólares en 2006) y de educación especial (11.800 millones de dólares en 2006) [US Bureau of the Census, 2009, p. 141]. Esto es coherente con la observación de que la redistribución es difícil de llevar a cabo a nivel local. En la educación superior, una gran cantidad del gasto federal se dirige hacia la investigación. La información procedente de la investigación es un bien público, y hemos visto que la provisión centralizada o la subvención de bienes públicos puede evitar el problema del free rider que podría surgir a nivel local.

Debe señalarse, sin embargo, que el papel federal en la educación no termina con la financiación. Un vasto cuerpo de leyes y regulaciones federales rige la educación pública. La legislación federal abarca temas tan diversos como la formación de profesores, las bibliotecas, los estándares para estudiantes con discapacidad y la educación sexual. Los estados cuyas prácticas no sigan las normas pueden perder fondos federales. Así, aunque el sistema de financiación de la educación estadounidense parece ampliamente coherente con los principios básicos del federalismo óptimo, la división de la toma de decisiones no es tan clara como sugeriría la teoría.




Visión general

Al comienzo de este capítulo planteamos algunas preguntas relativas a los sistemas federales: ¿Es deseable la toma de decisiones descentralizada? ¿Cómo deberían asignarse las responsabilidades? ¿Cómo deberían financiarse los gobiernos locales? Nuestras respuestas sugieren que el federalismo es un sistema sensato. Permitir que las comunidades locales tomen sus propias decisiones muy probablemente aumenta la eficiencia en la provisión de bienes públicos locales. Sin embargo, también es probable que la eficiencia y la equidad requieran un papel económico significativo para un gobierno central. En particular, un sistema en el que únicamente se utilizan recursos locales para financiar bienes públicos locales es considerado por muchos como inequitativo.

Aunque nuestro enfoque se ha centrado naturalmente en cuestiones económicas, las cuestiones de poder y política nunca están muy lejos de la superficie en los debates sobre el federalismo. La dispersión del poder económico está generalmente asociada con la dispersión del poder político. ¿Cómo debería asignarse el poder? ¿Es su imagen del gobierno subfederal la de un gobernador racista que mantiene a los estudiantes negros fuera de la universidad estatal, o la de una reunión municipal en la que los ciudadanos toman democráticamente decisiones colectivas? Cuando piensa en el gobierno central, ¿se imagina a un burócrata indiferente y distante imponiendo regulaciones molestas, o a un abogado del Departamento de Justicia trabajando para garantizar los derechos civiles de todos los ciudadanos? Las diferentes imágenes coexisten en nuestras mentes, creando sentimientos contradictorios acerca de la distribución apropiada del poder gubernamental.

Resumen



• Los impuestos sobre la propiedad son una fuente importante de ingresos para los gobiernos estatales y locales. La «visión tradicional» del impuesto sobre la propiedad es que se trata de un impuesto especial sobre el suelo y las estructuras. La «nueva visión» sostiene que el impuesto sobre la propiedad es un impuesto general sobre todo el capital, con tipos que varían entre jurisdicciones y diferentes tipos de capital. La «visión de la tasa de usuario» considera los impuestos sobre la propiedad como un pago por los servicios públicos locales.

• El impuesto sobre la propiedad es muy impopular. Quizá su principal ventaja en el contexto de un sistema federal sea que puede administrarse localmente.

• Las subvenciones pueden ser condicionadas (categóricas) o incondicionadas (de suma fija). Cada tipo de subvención incorpora diferentes incentivos para los gobiernos locales. La combinación final entre un aumento del gasto y unos impuestos locales más bajos depende de las preferencias que determinan las elecciones locales.

• Los estudios empíricos de las subvenciones intergubernamentales indican un efecto papel matamoscas: un aumento del dinero de las subvenciones induce un mayor gasto en bienes públicos que un aumento equivalente de la renta local. Una posible explicación es que los burócratas explotan la información incompleta de los ciudadanos acerca de la restricción presupuestaria de la comunidad.

He mantenido también la terminología de las secciones anteriores para que el capítulo quede coherente.





\clearpage
\section{Federalismo Fiscal: Marco conceptual}
\puntosclave{
    En un sistema federal, diferentes gobiernos proporcionan diferentes servicios a jurisdicciones superpuestas.
    El modelo de club de formación de comunidades indica que el tamaño de la comunidad y la cantidad de bienes públicos dependen de las preferencias por los bienes públicos, de los costes de proporcionar servicios públicos y de los costes de la congestión.

• El modelo de Tiebout pone de relieve las funciones de la movilidad, los impuestos sobre la propiedad y las normas de zonificación en las finanzas públicas locales. Bajo determinadas condiciones, «votar con los pies» —trasladarse a la comunidad preferida— da lugar a una asignación de bienes públicos eficiente en el sentido de Pareto.

• Las desventajas de la descentralización son las externalidades entre comunidades, las economías de escala desaprovechadas en la provisión de bienes públicos, la imposición ineficiente y la falta de capacidad para redistribuir la renta.

• Las ventajas de la descentralización son la capacidad de modificar la combinación de servicios públicos para adaptarla a las preferencias locales, los efectos beneficiosos de la competencia entre gobiernos locales y el potencial para la experimentación de bajo coste a nivel subfederal.

• La responsabilidad local sobre la educación puede justificarse sobre la base de las diferentes preferencias entre comunidades. Sin embargo, puede ser apropiada cierta participación federal en la distribución de los recursos disponibles para la educación.
}

A veces es útil pensar que las decisiones de finanzas públicas son tomadas por un único gobierno. Sin embargo, un número asombroso de entidades tiene el poder de recaudar impuestos y gastar.\footnote{En Estados Unidos, por ejemplo, hay más de 89.000 jurisdicciones gubernamentales: 1 federal, 50 estatales, 3.033 de condado, 19.492 municipales, 16.519 de municipio, 13.051 distritos escolares y 37.381 distritos especiales (US Bureau of the Census, 2009, p. 259).
} La interacción entre los distintos niveles de gobierno desempeña un papel crucial en las finanzas públicas. Examinemos las cuestiones fiscales que surgen en los sistemas federales.




\tectolor{red}{\textbf{Origen:} ICEX-CECO} l. Introducción
La Teoría de la Hacienda Pública se encarga de analizar la intervención pública a través de distintos programas de ingresos y gastos. Sin embargo, esta intervención no tiene por qué llevarse a cabo enteramente desde el gobierno central. En la práctica se observa que la intervención puede llevarse a cabo a través de varios niveles, existiendo en muchos países varios niveles de gobierno. La literatura económica se refiere a esta estructura multinivel como Federalismo Fiscal. 

La Teoría del Federalismo Fiscal analiza la división de las responsabilidades políticas a través de distintos niveles de gobierno y la interacción tanto vertical como horizontal entre éstos. Concretamente, estudia los \textbf{efectos que tiene la descentralización} de competencias sobre el sector público y el privado, y los \textbf{efectos de la financiación} intergubernamental y las interacciones fiscales.

Como señaló OATES (2005), se puede observar una gran evolución dentro de esta teoría que permite clasificar las aportaciones en este ámbito en dos generaciones. La primera generación se desarrolló durante los años 50 y 60. Se integraba dentro de la Teoría de las Finanzas Públicas, y tenía como objeto definir las competencias de los distintos niveles de gobierno. y disei'iar un sistema de financiación adecuado. De esta manera, las aportaciones de la primera generación afirmaban que la descentralización de gastos e ingresos públicos estaba justificada para alcanzar el objetivo de un mayor nivel de bienestar social. Con este sistema. se podía obtener una asignación eficiente, aunque los recursos móviles podrían interferir en esa asignación. Consideraban así un contexto sin restTicciones ni problemas de información. Más recientemente han surgido nuevas aportaciones que incluyen características fundamentales de la Teoría de la Elección Colectiva, como, por ejemplo: problemas de información, maximización de los intereses del gobierno, etc. Estas aportaciones suponen la creación de la segunda generación de la Teoría del Federalismo Fiscal, y analizan cómo interfieren los fallos propios del Sector Público al sistema descentralizado y los problemas endógenos incluidos en éste.

El interés por el análisis de las estructuras multinivel ha aumentado a lo largo de los años. Comenzó siendo de gran interés para algunos de los países más desarrollados. Pero este interés se amplió a economías en desarrollo desde los años 90, cuando muchas comenzaron a experimentar procesos de descentralización.

Cabe preguntarse entonces, cuáles son los beneficios y las ventajas que puede suponer un sistema descentralizado, y cómo debe ser su estructura. Para esto, se va a analizar por un lado la descentralización de los gastos. y. por otro. la de los ingresos. En ambos casos se hará referencia a las aportaciones de las dos generaciones de la Teoría del Federalismo Fiscal.

En comparación con la mayoría de las demás naciones desarrolladas, los gobiernos subnacionales (estatales y locales) de Estados Unidos recaudan una proporción mucho mayor de los ingresos públicos totales (nacionales más subnacionales) y gastan una proporción algo mayor del gasto público total. Un estudio reciente de las naciones de la OCDE mostró que los gobiernos subnacionales de la nación promedio recaudan solo el 27,2\% de los ingresos públicos totales, mientras que en Estados Unidos los gobiernos subnacionales recaudan el 36,8\% de los ingresos totales. Las diferencias entre países por el lado del gasto son aún más significativas: el gobierno subnacional de la nación promedio de la OCDE representa el 30,4\% del gasto, frente al 46,7\% en Estados Unidos.

El mayor nivel de centralización en otras naciones existe porque, en muchos países, como México, Austria y Noruega, los gobiernos subnacionales casi no tienen poder legal para gravar a los ciudadanos: este poder está reservado al gobierno central. Además, en la mayoría de los países, los gobiernos centrales redistribuyen una proporción mayor de sus ingresos a los gobiernos subnacionales. Muchos países practican la igualación fiscal, mediante la cual el gobierno nacional distribuye subvenciones a los gobiernos subnacionales en un esfuerzo por igualar las diferencias de riqueza. Esto puede lograrse proporcionando mayores subvenciones nacionales per cápita a las áreas subnacionales más pobres. En Austria, por ejemplo, el gobierno federal compensa más de la mitad de la diferencia entre áreas subnacionales en los ingresos que estas son capaces de recaudar mediante impuestos. El gobierno federal de Estados Unidos es notable porque no utiliza subvenciones para la igualación; el único programa de este tipo, iniciado por el presidente Richard Nixon a principios de la década de 1970, fue eliminado en 1986.

Otras naciones también tienen una distribución muy diferente del gasto entre los gobiernos nacionales y subnacionales. En Estados Unidos, por ejemplo, entre el 30 y el 40\% del gasto estatal y local se dedica a la educación, mientras que el promedio en las naciones de la OCDE es de alrededor del 20\%, lo que pone de relieve el mayor papel que desempeña el gobierno central en la educación en otros países.

Los últimos años han visto un movimiento hacia la descentralización fiscal en todo el mundo. En Estados Unidos, se han incrementado los esfuerzos por trasladar el control y la financiación de los programas públicos a los estados, como demuestra el ejemplo de la reforma de la asistencia social. En países tan diversos como Hungría, Italia, Corea del Sur, México y España, ha habido esfuerzos por trasladar la responsabilidad de la atención sanitaria, la educación y la asistencia social de los gobiernos nacionales a los subnacionales. Así, en la mayoría de los países, el gasto de los gobiernos subnacionales ha aumentado durante las últimas dos décadas, financiado a menudo mediante subvenciones del gobierno nacional. Este aumento de la financiación y del control ha ido acompañado normalmente de una creciente imposición de normas nacionales y estándares de calidad sobre los bienes proporcionados localmente (como planes de estudio nacionales cada vez más rígidos en educación).

• ¿Es deseable un gobierno descentralizado?
• En caso afirmativo, ¿qué niveles de gobierno deberían decidir sobre las diferentes políticas?
• ¿Son los impuestos recaudados localmente una buena manera de pagar los servicios proporcionados por los gobiernos estatales y locales? ¿O debería proceder el dinero del gobierno federal?

Estas son cuestiones importantes en Estados Unidos, donde la división apropiada del poder entre los distintos niveles de gobierno ha sido objeto de controversia desde la fundación de la nación. Las cuestiones tienen igual importancia para China, que está considerando si transferir o no poder a los gobiernos provinciales, y para las naciones europeas, que actualmente están decidiendo qué funciones de formulación de políticas económicas serán cedidas a la Unión Europea. Este capítulo examina los aspectos normativos y positivos de las finanzas públicas en un sistema federal.

\textcolor{blue}{\textbf{OJO -> Origen:} GRUBER (2017)}
Este debate plantea la importante cuestión del federalismo fiscal óptimo, la cuestión de qué actividades deberían tener lugar en qué nivel de gobierno. Quienes se oponen a la centralización tienen razón al afirmar que la provisión local de servicios gubernamentales permite a las comunidades elegir el paquete de servicios que mejor se ajusta a los gustos de sus residentes, mejorando potencialmente la eficiencia de la provisión de bienes públicos. Por otro lado, los defensores de la centralización también tienen razón al afirmar que, en algunos casos, los programas que mejor sirven a los intereses locales pueden no ser de interés nacional.

Como podemos intuir, vamos a presentar el conjunto de cuestiones que rodean el gasto de los gobiernos estatales y locales, o subnacionales, y la división de responsabilidades entre los diferentes niveles de gobierno. Comenzamos con un análisis de la actual división de responsabilidades en diferentes naciones desarrolladas. A continuación, pasamos a analizar si la provisión de bienes públicos por parte de los gobiernos locales resuelve los problemas de la provisión gubernamental de bienes públicos. En particular, al permitir que los individuos elijan la jurisdicción que mejor se ajuste a sus gustos, la provisión de bienes públicos por parte de los gobiernos locales puede permitir a los gobiernos locales proporcionar la cantidad óptima de bienes públicos, superando los problemas de revelación de preferencias y agregación de preferencias que obstaculizan las decisiones sobre la provisión nacional de bienes públicos.

Más adelante abordaremos si el gobierno debería redistribuir recursos entre comunidades y cómo debería hacerlo. Existen enormes diferencias entre comunidades en cuanto a la capacidad de financiar bienes públicos locales, en gran medida debido a diferencias en la capacidad recaudatoria (o recaudación per cápita). ¿Deberían preocuparse los gobiernos estatales y federal por estas diferencias? En caso afirmativo, ¿qué instrumentos pueden utilizar estos niveles superiores de gobierno para redistribuir recursos entre las comunidades?


\subsection{Comunidades heterogéneas: ¿Por qué se forman comunidades diferenciadas?}
Un sistema federal consta de diferentes niveles de gobierno que proporcionan bienes y servicios públicos y tienen cierto margen para tomar decisiones. La materia del federalismo fiscal examina las funciones llevadas a cabo por los diferentes niveles de gobierno y cómo interactúan entre sí los diferentes niveles de gobierno. Un sistema federal está más centralizado que otro cuando una mayor parte de sus poderes de toma de decisiones está en manos de autoridades con una jurisdicción más amplia. La medida más común del grado en que un sistema está centralizado es la ratio de centralización, la proporción de los gastos públicos directos totales realizados por el gobierno central.\footnote{El gasto público directo comprende todo el gasto excepto las transferencias realizadas a otras unidades gubernamentales.
} Las ratios de centralización varían ampliamente entre naciones. En Francia, es del 81\%; en Canadá, del 43\%; y en Estados Unidos, del 46\%.

Sin embargo, la ratio de centralización no es en absoluto un indicador infalible. Por ejemplo, los estados y las localidades realizan gastos en ordenadores para bibliotecas públicas, pero parte del dinero procede en forma de subvenciones del gobierno federal. La Ley de Protección Infantil en Internet exige que las bibliotecas instalen programas informáticos para filtrar materiales obscenos; las bibliotecas que no cumplen pierden sus subvenciones. La mayoría de las bibliotecas cumplen. ¿Quién está realmente al mando? La cuestión es que, si el comportamiento de gasto de los gobiernos locales y estatales está restringido por el gobierno central, la ratio de centralización subestima el verdadero grado de centralización del sistema.

De hecho, una cantidad sustancial del gasto estatal y local viene dictada por el gobierno federal. El gobierno federal simplemente ordena que el gobierno subfederal proporcione determinados servicios, pero sin un aumento correspondiente del apoyo financiero.

Una serie de actividades importantes están principalmente en manos de los gobiernos estatales y locales, entre ellas la educación y la seguridad pública. Por otro lado, el gobierno federal tiene toda la responsabilidad de la defensa y la Seguridad Social. Y los tres niveles de gobierno gastan cantidades sustanciales de dinero en asistencia social pública. ¿Es sensata esta división de poderes en el sistema fiscal de Estados Unidos? Antes de proporcionar una respuesta, necesitamos analizar las características especiales asociadas con el gobierno local.

Existe, por tanto. un trade-o.fl ya que corno se sefialó anteriormente, cuanto más homogéneas sean las preferencias de los agentes de una jurisdicción y más pequeña sea la jurisdicción. más probabilidades hay de que la provisión del bien público local satisfaga las demandas de estos agentes, y que sea eficiente.

Sin embargo, si existen economías de escala en la producción del bien. tiene sentido tener jurisdicciones de mayor tamafio. Pero pueden surgir entonces problemas de congestión. de forma que la calidad del bien se vea reducida según éste es consumido por más agentes.

Habrá que ponderar los beneficios que supone la división del coste de provisión entre un mayor número de agentes. y los costes en el bienestar de los agentes que esto supone. 


\subsubsection{Fuerzas centrípetas: Modelo de BUCHANAN (1965)}
Para comprender las funciones fiscales apropiadas de las jurisdicciones locales, examinamos por qué se forman las comunidades. En este contexto, es útil pensar en una comunidad como un club (una asociación voluntaria de personas que se agrupan para compartir algún tipo de beneficio). La Teoría de Clubs fue formulada incialmente por BUCHANAN (1965) y constituye, en la actualidad, una de las herramientas centrales del análisis normativo de la agrupación y descentralización en Economía Pública. Para este autor, un club se concibe como una asociación voluntaria de individuos que desean compartir un mismo beneficio. Veámoslo con un ejemplo.

Consideremos un grupo de personas que desean agruparse para comprar terrenos para un uso comunitario. Para simplificar, supongamos que todos los miembros del grupo tienen preferencias idénticas y que pretenden compartir por igual el uso del parque y sus costes. La comunidad puede excluir sin coste alguno a todos los que no son miembros, y funciona sin costes de transacción. Dado el supuesto de preferencias idénticas, solo necesitamos considerar los deseos de un miembro representativo. Deben tomarse dos decisiones: \textbf{qué tamaño de parque adquirir ($G$) y cuántos miembros tener en la comunidad ($n$)}. Suponiendo que quiere maximizar el bienestar de sus ciudadanos, ¿cómo decide la comunidad?\footnote{Analíticamente, estaríamos resolviendo el siguiente problema. Sea la función de utilidad del agente representativo: $U(G,X)$; donde $G$ es la cantidad de bien público (o, de forma equivalente, el tamaño óptimo de los bienes comunitarios); y, $X$ el conjunto de alternativas de consumo factibles, agrupadas bajo un bien numerario, que actúa como sustituto del bien comunitario y, por tanto, cuyo valoración será el coste de oportunidad del bien común. El individuo está sujeto a la restricción presupuestaria $Y=X+\left\[\frac{P_G}{n}\right]G$. Además, hay que tener en cuenta que el bien público está sujeto a congestión: $C(n)$, creciente en $n$. En este caso, el óptimo se alcanzaría en la siguiente igualdad: $n\left[\frac{\partial U/\partial G}{\partial U/\partial X}\right]=\frac{\mathrm dC(n)}{\mathrm dn}$; es decir, el beneficio marginal social sea igual a los coste marginal social. Se conoce como la \textbf{condición de SAMUELSON}.
}

\paragraph{Número de miembros: ¿cuántos miembros debería tener?}
Consideremos primero la relación entre el coste total por miembro y el número de miembros, dado que se selecciona un terreno de un determinado tamaño. Claramente, cuanto mayor es la comunidad, más personas hay para soportar el gasto y menor es la contribución requerida por miembro. Si el coste per cápita disminuye continuamente con el tamaño del número de miembros, ¿por qué no invitar simplemente a tantas personas como sea posible a unirse? El problema es que, a medida que más personas se unen a la comunidad, el parque se congestiona ($C(n)$). El coste marginal de congestión mide el coste en dólares de la congestión incremental creada por cada nuevo miembro. Suponemos que el coste marginal de congestión aumenta con el número de miembros. La comunidad debería ampliar su número de miembros hasta que la disminución marginal de la cuota de afiliación sea exactamente igual al aumento marginal por persona de los costes de congestión.

Es decir, al ampliar el tamaño de la jurisdicción surgen dos efectos que actúan en sentido contrario: (i) beneficio marginal, ya que la financiación del bien se reparte entre más individuos, por lo que el coste de financiación per cápita se reduce; y, (ii) coste marginal al aumentar la congestión del bien público.\footnote{\textbf{NOTA.-} En el caso de un bien público puro, para el que no existen costes de congestión, el tamaño del club es ilimitado.
}

\paragraph{Provisión: ¿cuál es el tamaño óptimo de los bienes comunitarios?}
Pasemos ahora al reverso del problema: para cualquier número dado de miembros de la comunidad, ¿qué tamaño debería tener el parque? Un parque más grande proporciona mayores beneficios, aunque, como la mayoría de los bienes, suponemos que está sujeto a una utilidad marginal decreciente. El coste marginal por miembro del aumento de la superficie del parque es simplemente el precio del terreno adicional dividido por el número de miembros que comparten su coste. La superficie debería aumentarse hasta el punto en el que el beneficio marginal de cada miembro sea exactamente igual al coste marginal por miembro.


\paragraph{Club óptimo: ¿cómo se define?}
Ahora podemos juntar estas dos partes de la imagen para describir una comunidad o club óptimo. La comunidad óptima es aquella en la que el número de miembros y el nivel de servicios satisfacen simultáneamente la condición de que el coste marginal sea igual al correspondiente beneficio marginal. Aunque este modelo de club es muy sencillo, \textbf{pone de relieve los aspectos cruciales del proceso de formación de comunidades}. En concreto, sugiere cómo el tamaño de la comunidad depende del tipo de bienes públicos que las personas quieren consumir, del grado en que estos bienes están sujetos a congestión y de los costes de obtenerlos, entre otras cosas.


\begin{specialblock}{\textbf{Extensiones:} Adhesión a las nóminas}
    La separación de los individuos en distintas jurisdicciones en la práctica es mucho más compl��ia, ya que en muchas ocasiones la adhesión es no-anónima. de manera que las características de los nuevos miembros, además de sus rentas, son importantes para los miembros existentes del club.

    Como señala SCOTHMER (1997). en muchas ocasiones puede no existir un equilibrio en la asignación de residentes a distintas jurisdicciones.
    
    Además, existen análisis que consideran que las adhesiones no-anónimas afectan también al coste de provisión de bienes públicos y las demandas de los bienes públicos locales. Esto hace que los individuos de una jurisdicción tengan en cuenta las características del resto de individuos (o potenciales residentes).
\end{specialblock}


\subsubsection{Fuerzas centrífugas}
Como vemos, resulta muy intrigante ver como podemos utilizar este sencillo modelo para verificar las fuerzas centrítepas que actúan en favor de la integración. Es evidente que las comunidades políticas no son sujetos ajenos a la transformación. Sin embargo, también es fácil intuir que existen elementos que actúan de pegamento político, social o económico, y a veces resulta difícil identificar algo como tal. De esta forma, el modelo antes visto nos proporciona una herramienta de simplificación para analizar y contrastar nuestras intuiciones personales. No obstante, estas no son estables y constantes, sino que actúan frente a otras cuya dirección es opuesta. Es decir, existen fuerzas centrífugas (políticas, sociales u económicas) que tienden a desestabilizar estos sistemas políticos. 

Recuperemos el ejemplo anterior ....

Por ejemplo, BUCHANAN y FAITH (1987) plantearon un modelo donde determinadas haciendas territoriales se aprovechaban de otras a través de un sistema de nivelación por transferencias. Así, los autores señalan que una jurisdicción tendrá incentivos para separarse cuando obtenga más pérdidas que ganancias por formar parte del sistema descentralizado. De hecho., esos incentivos pueden establecer límites a la explotación entre jurisdicciones. BOLTON y ROLAND (1997) mostraron que un aspecto importante para la desintegración del sistema se debe a las diferencias que se dan entre distintas jurisdicciones en cuanto al nivel de preferencias satisfechas. En linea con los resultados anteriores, señalaban que la amenaza de secesión puede poner límites a las políticas nacionales. 

DE FIGUEIRDO y WEINGAST (2002), plantean la posibilidad de que existan fuerzas endógenas que actúen como desestabilizantes. Por ejemplo, el solapamiento del gobierno central o la progresiva erosión de las funciones locales autónomas, el deterioro de las jurisdicciones mediante acciones oportunistas para beneficiarse fiscalmente a expensas de otras jurisdicciones.

No obstante. un aspecto crucial de las aportaciones de la segunda generación, es que consideran que existen estas características endógenas del sistema que pueden suponer una amenaza. pero que también existen otras que pueden suponer un refuerzo propio del sistema descentralizado. Para el refuerzo, una de las características más importantes es contar con un gobierno central que sea suficientemente fuerte para evitar capturas por parte de los gobiernos locales, y que éstos no lleven a cabo políticas de empobrecimiento del vecino. Pero es fundamental también que cuente con las restricciones adecuadas para reducir la intrusión central.

DOUGHERTY y WEINGAST destacan otras características importantes basadas en el análisis de algunos sistemas descentralizados fallidos, como los Estados Confederados de América. Señalan que el gobierno central fue incapaz de movilizar suficientes recursos económicos, de tener una moneda nacional estable y de evitar el establecimiento de barreras comerciales internas. En esta línea, una economía de mercado que funcione bien es clave para disciplinar a los gobiernos locales, como se mencionó anteriormente.

QIAN y WEINGAST (1997) enfatizan el papel que tiene la economía de mercado en reforzar el sistema descentralizado, y viceversa. Consideran un Sector Público relativamente descentralizado donde los gobiernos locales son los principales responsables de proveer bienes públicos y de ejercer funciones regulatorias en un mercado común sin barreras al comercio. Y, en relación con lo analizado antes, en este sistema la restricción presupuestaria de las jurisdicciones será dura. Argumentan que en este sistema habrá un compromiso creíble de proteger y fomentar los mercados privados y que, al contrario, los mercados contribuirán a la viabilidad de la estructura descentralizada. Estos autores, igual que BUCHANAN, hacen alusión a la competencia fiscal entre jurisdicciones, que permite castigar a los políticos locales y refuerza el sistema.

Pero queda pendiente la cuestión que se planteó de contener al gobierno central para evitar una intrusión en el sistema. Se trata de un asunto complejo donde la literatura más reciente destaca el papel del control por parte de instituciones. Por ejemplo, INMAN y RUBINFELD (1997) destacan el papel del poder judicial. RIKER y WEINGAST destacan el papel del sistema de partidos políticos con elementos locales. Además, también es importante limitar la participación de la política nacional con autoridades descentralizadas.

En definitiva, la segunda generación de la Teoría del Federalismo Fiscal resalta la importancia de la forma que toma un sistema descentralizado, y el papel que tiene la competencia entre jurisdicciones, para garantizar un sistema descentralizado estable. que satisfaga las preferencias de los agentes, y promueva la estabilidad económica.

Completa las aportaciones previas ya que el tipo de financiación elegida para las haciendas territoriales puede afectar al funcionamiento del sistema.

\textit{Ámalo o déjalo}. Cuando a las personas que se oponen a la política del gobierno federal de Estados Unidos se les da este consejo, es aproximadamente tan constructivo como decirles \textit{muérete}. Solo en casos extremos esperamos que las personas abandonen su país debido a la política gubernamental. Los grandes costes pecuniarios y psíquicos de emigrar llevan a que la opción más realista sea quedarse en casa e intentar cambiar la política. Por otro lado, la mayoría de los ciudadanos no están tan fuertemente vinculados a sus comunidades locales. Si no te gustan las políticas que se siguen en Barcelona, lo más fácil puede ser trasladarte a Badalona. Analicemos la relación entre la movilidad entre comunidades, la formación voluntaria de comunidades y la provisión eficiente de bienes públicos.


\subsection{Argumentos a favor: ¿por qué es deseable el autogobierno?}
¿Hasta qué punto es estrecha la analogía entre un club y una comunidad del mundo real? En muchos casos, es más estrecha de lo que podría pensarse. Por ejemplo, en Estados Unidos más de $50$ millones de personas viven en áreas gobernadas por asociaciones vecinales (TIERNEY, 2005). Estas comunidades cerradas deciden cuántos miembros tendrán, cuántos guardias de seguridad contratarán, si construirán campos de golf y piscinas comunitarias, etcétera. 

No obstante, en la mayoría de los casos, considerar las comunidades como clubes deja sin respuesta varias cuestiones importantes que son relevantes para comprender las finanzas públicas locales:
• ¿Cómo deben financiarse los servicios públicos? Un club de campo puede cobrar una cuota de afiliación, pero un municipio normalmente recauda impuestos para pagar los bienes públicos.
• Un club (o una comunidad cerrada) puede excluir a quienes no son miembros y eliminar así el problema del free rider. ¿Cómo pueden los municipios lograr este fin?
• Cuando las personas de todo el país se organizan en muchos clubes (comunidades) diferentes, ¿es equitativa y eficiente la asignación global de bienes públicos?

Ahora que conocemos por qué surgen las agrupaciones sociales, las razones por las que estás pueden tener diferentes tamaños y hemos visto cómo diferentes naciones han adoptado diferentes enfoques del federalismo fiscal, se plantea una pregunta natural: ¿Cuál es la división óptima de responsabilidades entre los diferentes niveles de gobierno? ¿Por qué debería hacerse algo por parte de los gobiernos locales? Alternativamente, ¿por qué existe algún papel para un gobierno central? ¿Y qué tipos concretos de programas son administrados de manera más apropiada en qué nivel de gobierno? 

Existen diversas justificaciones que se han ido desarrollando para responder a estas preguntas, veamos algunas.


\subsubsection{Eficiencia: Modelo de TIEBOUT}
Sin duda, la más importante/conocida tiene que ver con el concepto de eficiencia económico, con el que tan familiarizados debemos de estar llegados a este punto. 

Sabemos que la intervención del Sector Público se puede justificar por razones de estabilización, equidad y eficiencia. Así, una primera aproximación al asunto podría hacernos pensar que las mismas razones nos llevarían a la descentralización. Sin embargo, esto no es del todo cierto. Las dos primeras, estabilización y equidad, conducen a situaciones de díficil escapatoria, donde la deseabilidad de la fragmentación en diferentes niveles de gobierno no es concluyente (o incluso proporciona argumentos en contra).\footnote{Puede consultarse el [\authorfont{Ver Anexo \ref{estabilizacion_equidad}}] para más información.
} Pero el tercero si nos arroja luz sobre el tema: en un importante artículo, TIEBOUT (1956) argumentó que la capacidad de los individuos para trasladarse entre jurisdicciones produce una solución similar a la del mercado para el problema de los bienes públicos locales; las razones que dió no son difíciles de comprender. 

Como sabemos, dos problemas principales de la provisión gubernamental de bienes públicos son: la revelación de preferencias y la agregación de preferencias; es difícil diseñar instituciones democráticas que hagan que los individuos revelen honestamente sus preferencias por los bienes públicos, y también es difícil agregar las preferencias individuales en una decisión social. Como resultado, los gobiernos son a menudo incapaces de proporcionar en la práctica el nivel óptimo de bienes públicos. Partiendo de este presupuesto, TIEBOUT se preguntó: \textbf{¿qué hay en el mercado privado que garantiza la provisión óptima de bienes privados que está ausente en el caso de los bienes públicos?} 

Su idea fue que los factores ausentes del mercado de bienes públicos eran la posibilidad de elegir y la competencia. La posibilidad de elegir es la fuerza fundamental que induce eficiencia en los mercados de bienes privados. Si una empresa está vendiendo un bien inferior en relación con sus competidores, los consumidores comprarán a los competidores, no a la empresa. Esta competencia lleva a las empresas a producir eficientemente en el mercado perfectamente competitivo de bienes privados. Sin embargo, con muchos bienes públicos no existe posibilidad de elegir. Los individuos no debaten si vivir en Estados Unidos o en Canadá basándose en si el misil marginal es producido por el gobierno federal. Los votantes pueden elegir entre partidos basándose en sus promesas, pero este es solo uno de un gran número de factores que determinan los votos para cargos federales, y el proceso de cambiar la toma de decisiones federal es lento. Así, debido a que existe poca competencia real a la que se enfrenta un gobierno nacional, cuando toma sus decisiones de proporcionar bienes públicos las decisiones pueden dar lugar a una provisión ineficiente de bienes públicos.

No obstante, señaló TIEBOUT, la situación es diferente cuando los bienes públicos son proporcionados a nivel local por ciudades y municipios. En este caso, argumentó, la competencia surgirá naturalmente porque los individuos pueden votar con los pies: \textbf{si no les gusta el nivel de provisión de bienes públicos de un municipio, pueden trasladarse al municipio de al lado, sin una alteración de sus vidas ni de lejos tan grande como la que supondría trasladarse a otro país}.\footnote{Supongamos que descubrimos que tu biblioteca pública local estaba gastando $37.000€$ al año para proporcionar al director de la biblioteca un coche deportivo. Este despilfarro eleva claramente los impuestos locales que pagamos para financiar el gobierno municipal, o como mínimo la sensación de indignación. En este caso, tienes una opción realista: puedes trasladarte al municipio de al lado, que puede ser similar en la mayoría de las dimensiones pero mejor en términos de disciplina fiscal. De esta forma, con los bienes públicos locales tenemos un nuevo mecanismo de revelación de preferencias: la movilidad. Por el contrario, supongamos que leemos que el Ministerio de Transportes estaba gastando $110€$ en una bombilla, o $300€$ en un solo martillo. ¿Qué podríamos hacer al respecto? Es poco probable que nos trasladásemos a otra nación. Prodríamos expulsar del poder mediante  voto al partido gobernante, pero éste se basa en un gran número de factores, de los cuales es solo uno. Así que realmente hay poco que hacer para poner fin a tal ineficiencia.
} ¿Por qué? TIEBOUT argumentó que esta amenaza de salida puede inducir eficiencia en la producción local de bienes públicos. De hecho, fue un paso más allá y argumentó que, bajo ciertas condiciones, la provisión de bienes públicos será plenamente eficiente a nivel local. ¿Cuáles? Analicemos el modelo formal que subyace a su intuición.

El modelo de TIEBOUT supone que hay muchas personas que se distribuyen entre municipios que proporcionan diferentes niveles de bienes públicos, y que hay suficientes comunidades diferentes para que cada individuo pueda encontrar una con servicios públicos que satisfagan sus demandas. Dentro de este mundo, la provocadora afirmación depende crucialmente de las siguiente condiciones:
\begin{la}
    \item \textbf{Ausencia de externalidades}. Las actividades gubernamentales no generan externalidades. Como se señala más adelante, los efectos de desbordamiento entre comunidades pueden conducir a ineficiencias;
    \item \textbf{Movilidad perfecta}. Los individuos son completamente móviles. Cada persona puede trasladarse sin coste a una jurisdicción cuyos servicios públicos sean los mejores para ella. La ubicación de su lugar de trabajo no impone ninguna restricción sobre dónde reside y no afecta a sus ingresos;\footnote{De esta forma, el modelo resuelve los problemas de revelación y agregación de preferencias que causan dificultades con la provisión pública de bienes públicos. ¿Por qué? Porque, al no haber incentivo para que las personas mientan con un impuesto uniforme que financia los bienes públicos, sencillamente no existe problema de revelación. Por otro lado el problema de agregación de preferencias también se resuelve porque todos los habitantes del municipio quieren el mismo nivel de bienes públicos, $G_i$, y el gobierno municipal puede simplemente dividir esa cantidad por la población para obtener la financiación apropiada.
    
    Para verlo claro supongamos que Ramón se une a un municipio de individuos idénticos a él, que cuenta con cuatro bibliotecas que costean comúnmente a un precio de cien euros por habitante. Supongamos que Ramón miente diciendo que tiene las mismas preferencias que Ana, otra ciudadana. En el modelo de TIEBOUT, para llevar a cabo esa mentira, tendría que trasladarse realmente a un municipio de individuos como Ana. En el municipio de Ana, solo desean dos bibliotecas, a un coste de cincuenta euros por habitante. Al trasladarse al municipio de Ana, Ramón paga la mitad de lo que pagaba pero obtiene solo la mitad de lo que obtenía. Ramón no tiene ningún incentivo para mentir porque debe actuar conforme a su mentira trasladándose a un municipio diferente que coincida con sus preferencias declaradas. Es decir, Ramón no puede actuar como free rider cuando los individuos de cada municipio son idénticos y comparten por igual la financiación del bien público. 
    }
    \item \textbf{Información perfecta}. Las personas tienen información perfecta con respecto a los servicios públicos y los impuestos de cada comunidad;
    \item \textbf{Rendimientos constantes a escala}. El coste por unidad de servicios públicos es constante, de modo que si la cantidad de servicios públicos se duplica, el coste total también se duplica. Además, si el número de residentes se duplica, la cantidad del servicio público proporcionado debe duplicarse;\footnote{Para ver por qué estas condiciones son necesarias para que un equilibrio de TIEBOUT sea eficiente, imaginemos, por el contrario, que el coste por unidad de servicios públicos disminuyera a medida que aumentara la escala de provisión. En ese caso, habría economías de escala que las comunidades que operan de forma independiente podrían no aprovechar. Este supuesto hace que el servicio público sea esencialmente un bien privado proporcionado públicamente. Los bienes públicos puros (como la defensa nacional) no satisfacen este supuesto. Sin embargo, muchos servicios públicos locales, como la educación y la recogida de basura, se ajustan razonablemente bien a esta descripción.
    }
    \item \textbf{Impuestos de suma fija}. Los servicios públicos se financian mediante un impuesto proporcional sobre la propiedad y el tipo impositivo puede variar entre comunidades.
\end{la}

Con los problemas de revelación y agregación de preferencias resueltos, la fijación de precios de LINDAHL funciona en el modelo de TIEBOUT. Cada individuo declara su verdadera valoración del bien público, las valoraciones se suman y, a continuación, a cada individuo se le cobra el coste total del bien público dividido por el tamaño de la población ([\authorfont{Ver Tema 3.A.24}]). Esto constituye un equilibrio porque cada persona está dispuesta a pagar su parte del impuesto para obtener el bien público, y se cumple la condición para la provisión óptima de bienes públicos porque el nivel de bienes públicos proporcionado viene determinado por la suma de los beneficios individuales


Según la misma lógica por la que el equilibrio competitivo proporciona el nivel eficiente de bienes privados, la competencia entre localidades en la provisión de bienes públicos proporcionará el nivel eficiente de bienes públicos. Los municipios que no proporcionen niveles eficientes de bienes públicos perderán ciudadanos en favor de los municipios que sí logren la eficiencia, y finalmente dejarán de funcionar.


\footnote{Este modelo hace una serie de supuestos que no son realistas, como veremos después. Sin embargo, el mensaje principal del modelo, que la competencia entre jurisdicciones locales ejerce presiones competitivas sobre la provisión de bienes públicos locales, es un mensaje importante que es coherente con la evidencia que examinamos más adelante.
} 









\textbf{Tiebout y el mundo real}

El modelo de Tiebout claramente no es una descripción exacta del mundo real. Las personas no son perfectamente móviles; no hay suficientes comunidades para proporcionar a cada familia un conjunto de servicios que se adapte perfectamente a ella; etcétera. Además, contrariamente a lo que implica el modelo, observamos muchas comunidades con enormes diferencias de renta y, por tanto, presumiblemente con diferentes niveles deseados de provisión de servicios públicos. Basta considerar cualquier gran ciudad.

Sin embargo, no deberíamos descartar demasiado precipitadamente el mecanismo de TIEBOUT. Existe mucha movilidad en la economía estadounidense. Un patrón persistente es que, en cualquier año dado, alrededor del 17\% de los estadounidenses tienen residencias diferentes de las que tenían el año anterior (US Bureau of the Census, 2009). Además, la mayoría de las áreas metropolitanas permiten una amplia variedad de elección con respecto al tipo de comunidad. Dentro de un radio de 20 millas de una gran ciudad estadounidense, a menudo se puede elegir establecerse entre varios cientos de suburbios. Ciertamente, la observación casual sugiere que entre los suburbios existe una considerable segregación residencial por renta, que la zonificación excluyente se practica ampliamente y que los niveles de servicios difieren (incluso cuando las rentas son similares).

Ha habido varias pruebas empíricas formales de la Hipótesis de TIEBOUT. Un tipo de estudio examina si los valores de los servicios públicos locales y los impuestos se capitalizan en los valores de las propiedades locales. La idea es que, si las personas se trasladan en respuesta a los conjuntos locales de impuestos y servicios públicos, las diferencias entre estos conjuntos deberían reflejarse en los valores de las propiedades. Una comunidad con mejores servicios públicos debería tener valores de las propiedades más elevados, manteniéndose iguales las demás cosas (incluidos los impuestos). Estos estudios de capitalización se analizan más adelante en este capítulo en el contexto de la imposición sobre la propiedad. Como se señala allí, la capitalización parece ser, efectivamente, un fenómeno generalizado.

Otro tipo de estudio examina si los cambios en los niveles de bienes públicos locales conducen a la migración entre jurisdicciones. Por ejemplo, BANZHAF y WALSH (2008) encuentran pruebas de que las personas migran entre comunidades a medida que cambia la calidad del aire local. Concluyen que \textit{los hogares sí parecen votar con los pies en respuesta a cambios en los bienes públicos}. Estos resultados sugieren que, al menos en algunos contextos, el modelo de TIEBOUT es una buena representación de la realidad.


\textcolor{red}{\textbf{Origen:} ICEX-CECO}
\textcolor{red}{IJ. Descentralización por el lado del gasto : primera y segunda generación}



ll.l.J. Motivos de actuación del sector público y descentralización

i. Estabilización

Con relación a la función de estabilización, según Wildasin (20 18). existe un amplio consenso de que debe ser una decisión cemralizada. Existen varios argumentos que justifican esta afirmación.

Un pnmcr argumento considera que, dado que las distintas Jurisdicciones que forman una unión fiscal son economías relativamente más abiertas que la economía nacional. el impacto de políticas macroeconómicas de estabilización se ve diluido a través de los flujos de capital y trabajo cuando se llevan a cabo a nivel local. Y es que, tal y como seíiala el modelo ue MUNDELL y FLEMING desarrollado en los afios 60, una política fiscal local se filtraría a otras jurisdicciones sin ser efectiva. Por tanto. se cumplen los requisitos ser'íalados por MUNDELL en 1961 para formar un Area Monetaria Óptima, de manera que los gobiernos locales no tendrán autonomía sobre la política monetaria, pues; esta debe estar centralizada (o pertenecer incluso a una entidad supranacional, como es el caso de la política monetaria en la zona euro).

Por otro lado. también es importante señalar que la centralización de la función de estabilización permite un desarrollo coordinado de esta entre las distintas regiones, en caso de existir asimetría en los shocks.

ii. Redistribución

En cuanto a la función de redistribución, la primera generación de la Teoría del Federalismo Fiscal. también consideraba que el rol principal debe pertenecer al gobierno central. Esto se debe a que, en caso de ser descentralizada. podrían surgir problemas de equidad interte1Titorial. Por ejemplo, si la función de redistribución fuera descentralizada. podría haber jurisdicciones que tuvieran una menor capacidad para llevar a cabo esta función con relación a otras. Persistiría. por lo tanto, la desigualdad entre distintas unidades.

Además. como señaló STIGLER en 1957, con alta movilidad de agentes, un gobierno local que fuera más redistributivo que los demás vería emigrar a los individuos más ricos. e inmigrar a los más pobres de otras jurisdicciones. Debido a esto. un gobierno local podría tener incentivos a no realizar gastos destinados a una redistribución y una mayor equidad, con el objetivo de evitar la emigración de individuos de rentas bajas, algo inconsistente con la concepción global de justicia redistributiva.

En definitiva, para tener una redistribución uniforme y efectiva, lo más deseable es que sea centralizada.

\textcolor{red}{\textbf{iii. Eficiencia}}
Por tanto, la descentralización se ha justificado tradicionalmente por motivos de eficiencia en la asignación. Esto está relacionado con los bienes públicos locales. que son aquellos bienes públicos cuyos beneficios están limitados espacialmente, de manera que sólo podrán ser consumidos por aquellos agentes que residen en un área determinada. 

De esta manera, las decisiones solo deberían tomarse a nivel nacional cuando afectan al total de la economía. Un ejemplo típico de esta situación. es el g asto en defensa, cuyos beneficios no afectan solo a un grupo de agentes de la economía. Estos bienes pueden proveerse a nivel nacional, aprovechando así las economías de escala debido a la mutualización de costes fijos, pero existen importantes argumentos que sugieren que una provisión a niveles inferiores es preferible en otros casos.

En primer lugar, es necesario que el Sector Público conozca qué bienes desean los individuos, y cuánto los valoran. Sin embargo. las preferencias individuales no son observables; y no existen incentivos para revelarlas, de manera que el Sector Público debe basarse en estimaciones. De esta forma, la descentralización presenta la ventaja de que los gobiernos territoriales conocerán mejor las preferencias y necesidades de los agentes que residan en sus jurisdicciones. y realizarán estimaciones más precisas. dado que están más próximos a los ciudadanos que el gobierno central.

Otro importante argumento a favor de la descentralización se refiere a la diversificación en la oferta de los bienes públicos. Una provisión centralizada puede traducirse en una provisión uniforme e idéntica para los distintos territorios. ya sea por: los problemas de información comentados en cuanto a las estimaciones que se dehen realizar; restricciones legales que garanticen la no discriminación: o por presiones políticas que traten de evitar una diferenciación, o de favorecer a una mayoría.

Una provisión uniforme de bienes públicos en todas las jurisdicciones solo cubrirá las necesidades de todos los agentes cuando sus preferencias sean homogéneas. En el caso contrario, siempre supondrá una pérdida de eficiencia en comparación con la provisión diferenciada.

Este argumento se puede ilustrar considerando el ejemplo de HINDRICKS y MYLES de una economía donde existen dos grupos de consumidores (A y B). Suponen así que existen n,, y ns miembros en cada grupo. Los gustos y preferencias en cuanto al bien público que se provee difieren entre cada grupo. Supondremos que el grupo B tiene una mayor preferencia por el bien público que el grupo A.

\textbf{GRÁFICO: CURVAS DE INDIFERENCIA GRUPOS A Y B. NIVEL DE UTILIDAD-PROVISIÓN DE BIEN PÚBLICO}

Se puede representar en el Gráfico I las curvas de indiferencia y los niveles de utilidad de cada
uno de los grupos.
Los niveles de provisión óptimos de los grupos A y B son GA y GB respectivamente, siendo
GA<G13. Podemos suponer que existe una provisión uniforme igual a Go, intcnnedia entre GA y
Gs. Con esta provisión se obtiene una pérdida de bienestar social en comparación a la que cada
grupo alcanzaría con una provisión diferenciada como la siguiente:

$L=n_A[u_A(G_A)-u_A(G_0)]+n_B[u_B(G_B)-u_B(G_0)]$

Esta pérdida puede minimizarse estableciendo un nivel de provi-sión G0 para el cual el beneficio marginal del grupo B compense la pérdida marginal del grupo A. Sin embargo, siempre habrá una pérdida de bienestar social positiva que será mayor cuanto mayor sea la dispersión de las preferencias de ambos grupos.

En definitiva, la provisión uniforme conlleva costes de bienestar social. y solo será preferible si los costes de una provisión diferenciada exceden sus beneficios.

Esta es la idea central del conocido \textbf{Teorema de la Descentralización} formulado por OATES en 1 972. El teorema señala que el nivel de bienestar alcanzado con una provisión descentralizada y diferenciada de acuerdo con las preferencias de las jurisdicciones será siempre PARETO superior al que resultaría de una provisión uniforme.

Para esto. es necesario que las preferencias sean heterogéneas de forma que existan demandas diferenciadas entre regiones, y que no existan efectos externos interregionales. Además. como sefiala el propio Oates. es necesario que existan impedimentos para que un gobierno central ofrezca esta provisión diferenciada, tales como la existencia de problemas de información. o las restricciones legales y políticas que se mencionaron previamente.

En relación a esto, es fundamental hacer mención al modelo desarrollado por TIEBOUT en el año 1956, que ofrece una solución a estos problemas. La intención inicial de su desarrollo era desafiar la afirmación de SAMUELSON de que la elección descentralizada no resulta en una provisión eficiente de bienes públicos. TIEBOUT señala así que existe un tipo de estos bienes para los que sí se tiene un equilibrio eficiente en sentido PARETO, los bienes públicos locales. Crea así la conocida como Hipótesis de TIEBOUT, considerada por muchos autores una pieza central de la Teoría del Federalismo Fiscal.

Tiebout supone que existen agentes con preferencias heterogéneas. Considera que existe movilidad perfecta y costes de transacción nulos, y que los gobiernos locales pueden proveer bienes públicos en el mínimo de los costes medios y financian esta provisión con impuestos de suma fija de manera que no generan distorsiones.

El autor se basa en la analogía de agentes con preferencias heterogéneas que adquieren bienes en el mercado, y agemes que eligen una residencia entre un grupo de diferentes localidades. Indica que los agentes revelan sus preferencias sobre los bienes públicos locales con la decisión de residencia.

La idea es que los agentes se irán agrupando según preferencias homogéneas o parecidas, y se desplazarán para instalarse en la jurisdicción que les ofrezca una combinación de gastos e ingresos públicos más acordes a sus preferencias. Es decir, votan con los pies, y muestran su acuerdo o desacuerdo con las ofertas públicas quedándose o marchándose de una jurisdicción.

Este emparejamiento entre demanda y oferta lleva a una provisión eficiente de bienes públicos locales. La descentralización, por tanto, permite crear un mercado para la provisión de bienes públicos locales en el que cada gobierno local es el oferente, y los ciudadanos revelan al final sus preferencias mediante su decisión de donde residir.

Como vemos, se trata de un análisis simplificado en el que estamos asumiendo decisiones en una sola dimensión. Otra importante limitación es el hecho de que en la práctica no existe esa movilidad perfecta con costes de transacción nulos, y los individuos no solo deciden dónde residir en función de los bienes públicos. También hay que señalar la dificultad que se da en la práctica de emplear impuestos de suma fija.

Otra limitación es que el autor considera que las distintas jurisdicciones son unidades económicas diferenciadas que pueden llevar a cabo la función de asignación de una manera que no es posible para el gobierno central. No explica, sin embargo, el proceso por el que estas jurisdicciones determinan cuál es el nivel de provisión de los bienes públicos locales.

Además, en la práctica existen efectos de desbordamiento entre distintas jurisdicciones, de forma que las decisiones de una unidad territorial afectan a otras, y para obtener el resultado que predice TIEBOUT no debería haber externalidades.

También. de acuerdo con la hipótesis de Tiebout, deberían existir un número altísimo de jurisdicciones para cubrir el abanico de preferencias. En relación con esto. se analiza más adelante la Teoría de los Clubes que analiza cuál es el tamaño óptimo de cada jurisdicción. En definitiva, las aportaciones de Tiebout destacan los beneficios potenciales de un sistema descentralizado en cuanto a la provisión eficiente de bienes públicos. No obstante. existe un trade-o.fr ya que la descentralización de ciertos servicios puede reducir el aprovechamiento de las economías de escala y producir pérdidas importantes de eficiencia productiva.

\begin{specialblock}{\textbf{Extensiones:} Modelo de TIEBOUT}
    Existen contribuciones y extensiones posteriores a la hipótesis de Tiebout, que incluyen análisis tanto teóricos como empíricos. Anal izan las implicaciones de la organización de 1 icbout. y del sistema estratificado.
    
    En cuanto al contenido empírico destaca el trabajo de Oates ( 1969), que indicaba que la votación con los pies resultaría en una «capitalización» al afectar al precio de las viviendas.
    
    Otras extensiones que podemos destacar en el marco teórico son : Hirschman ( 1970), Bergstrom y Goodman ( 1973). o Goldstein y Pauly que profundizan en la definición del nivel de provisión de las jurisdicciones ; Ellickson ( 1971 ), Westhoff ( 1977), Nechyba ( 1997), deBartolome ( 1990), Benabou ( 1996) o Durlauf ( 1996) que analizan el equilibrio en cuanto a los residentes de las jurisdicciones ; y Hamilton (1975), Milis y Oates ( 1 975), Rogcrson (1997), Henderson ( 1985), o Glacser y Gyourko ('.!002 ) que anal izan la segregación fiscal en distintas comunidades dentro de las jurisdicciones.
\end{specialblock}

Con todo esto, se puede concluir que la primera generación de la Teoría del Federalismo Fiscal abogaba por un sistema descentralizado en el que la función de estabilización y la de redistribución pertenecieran a un nivel centralizado de gobierno. mientras que la función de asignación pe1teneciera a los gobiernos locales.

Pero aquí es importante señalar, que en muchos sistemas descentralizados surge la dificultad de que muchas políticas públicas. debido a su propia naturaleza. tienen un impacto tanto en la asignación como en la redistribución. Por ello, una separación de estas funciones entre los distintos niveles de gobierno no siempre es posible.

Algunos ejemplos de estas políticas son las de sanidad. educación. desarrollo económico. o transporte, ya que se encuentran dentro de la función de asignación. pero también son fundamentales para lograr los objetivos de redistribución.

Además. de cumplirse la hipótesis de Tiebout y formarse jurisdi<;ciones de agentes con preferencias homogéneas, las ganancias en eficiencia que se producen por la descentralización, pueden suponer importantes aumentos en cuanto a la disparidad de la provisión de estos bienes públicos en el caso de que existan muchos grupos de preferencias.

\textcolor{blue}{\textbf{OJO -> Origen:} GRUBER (2017)}
\textcolor{red}{\textbf{Federalismo fiscal óptimo}}







\textcolor{blue}{\textbf{OJO -> Origen:} GRUBER (2017)}
\textcolor{red}{\textbf{Evidencia empírica: Modelo de TIEBOUT}}

El modelo de TIEBOUT impone claramente un conjunto muy restrictivo de supuestos si se toma literalmente, pero la intuición básica de que los individuos votan con los pies sigue siendo sólida. De hecho, dos tipos de pruebas revelan que la provisión de bienes públicos locales es generalmente coherente con la descripción de TIEBOUT.

\textbf{Similitud de los residentes entre áreas }
Una predicción clara del modelo de TIEBOUT es que las personas que viven en una determinada comunidad local deberían tener preferencias similares por los bienes públicos locales. Cuantas más comunidades locales haya entre las que elegir, según la lógica de este modelo, más podrán los residentes clasificarse en agrupaciones similares. Si una ciudad tiene solo un suburbio a una distancia que permita desplazarse al trabajo, será difícil para los residentes que trabajan en la ciudad votar con los pies si no les gusta el nivel de provisión de bienes públicos de ese único suburbio. Por tanto, una implicación contrastable es que, cuando las personas tienen más posibilidades de elección de comunidad local, los gustos por los bienes públicos serán más similares entre los residentes del municipio que cuando las personas no tienen muchas opciones (y, por tanto, no pueden clasificarse en comunidades de TIEBOUT de personas con ideas afines).

Evidencia favorable sobre este punto procede de GRAMLICH y RUBINFELD (1982), para un conjunto de datos de Estados Unidos (encuestas a hogares de Michigan sobre su demanda de bienes públicos. Los autores encontraron que en áreas metropolitanas más grandes (es decir, suburbios cercanos a ciudades), donde las personas tienen una mayor posibilidad de elegir en qué comunidad pueden vivir, las preferencias por los bienes públicos eran más similares dentro de los municipios que en áreas más pequeñas con menos municipios independientes entre los que elegir. Además, en las áreas urbanas/suburbanas, los residentes estaban mucho más satisfechos con el nivel de gasto en bienes públicos que en las áreas no urbanas, donde hay menos formas de votar con los pies porque hay menos municipios a los que mudarse. BERGSTROM et al. (1988) utilizaron los mismos datos para estimar las demandas individuales de bienes públicos y mostraron que la provisión de bienes públicos locales parecía satisfacer la condición de eficiencia de que el coste marginal es igual a la suma de las tasas marginales de sustitución de los residentes.

\textbf{Capitalización de las diferencias fiscales en los precios de las viviendas }
Para muchos individuos, la decisión sobre dónde vivir no está determinada principalmente por el nivel de bienes públicos locales. De hecho, muchos residentes ni siquiera demuestran el conocimiento básico de los impuestos y el gasto locales que se requiere para que opere el mecanismo de TIEBOUT. Al mismo tiempo, el mecanismo de TIEBOUT no requiere que todos los residentes estén dispuestos a votar con los pies, sino que suficientes residentes estén dispuestos a hacerlo para imponer la provisión óptima de bienes públicos.

Un municipio no tiene que vaciarse por completo antes de que los funcionarios locales capten el mensaje de que los residentes están descontentos con la provisión de bienes públicos; todo lo que se requiere es que exista suficiente movilidad entre una minoría informada en respuesta a las decisiones sobre bienes públicos. De hecho, se requiere muy poca movilidad real para que opere el mecanismo de TIEBOUT porque las personas no solo votan con los pies, sino que también votan con su bolsillo, en forma de precios de las viviendas. 

El modelo de TIEBOUT predice que cualquier diferencia en el atractivo fiscal de un municipio será capitalizada en los precios de las viviendas. El precio de cualquier vivienda refleja el coste (incluidos los impuestos locales) y los beneficios (incluidos los bienes públicos locales) de vivir en esa vivienda. Por tanto, los municipios que tengan un nivel relativamente alto de bienes públicos, dados los impuestos pagados, tendrán viviendas más caras; a la inversa, los municipios que tengan una carga fiscal relativamente más alta, dados los bienes públicos proporcionados, tendrán viviendas menos caras. La fijación de precios de las viviendas, por tanto, representa votar con el bolsillo: las personas pagarán más por una vivienda en un municipio que proporcione bienes públicos locales de manera más eficiente.

Existe una sólida evidencia a favor de votar con el bolsillo (\textcolor{red}{, como se examina en el recuadro de Evidencia Empírica}). Por tanto, incluso si algunos residentes no eligen su ubicación basándose en factores de TIEBOUT, suficientes residentes sí toman decisiones de esa manera como para que ello determine la fijación de precios de las viviendas entre las comunidades locales.


\textcolor{red}{\textbf{OJO -> Origen:} ROSEN y GARBER}

\textcolor{red}{Ventajas de un sistema descentralizado}

\paragraph{Preferencias heterogeneas}
Algunas personas quieren que los institutos de sus hijos tengan amplios programas deportivos; otras creen que esto es innecesario. Algunas personas disfrutan de los parques; otras no. Un gobierno centralizado tiende a proporcionar el mismo nivel de servicios públicos en todo el país, independientemente del hecho de que las preferencias de las personas difieran. Como observó de TOCQUEVILLE: \textit{En las grandes naciones centralizadas, el legislador está obligado a dar un carácter de uniformidad a las leyes, que no siempre se adapta a la diversidad de costumbres y de distritos}. Claramente, es ineficiente proporcionar a los individuos una cantidad mayor o menor de un bien público de la que desean si la cantidad que reciben puede adaptarse más estrechamente a sus preferencias. Bajo un sistema descentralizado, los individuos con preferencias similares por los bienes públicos se agrupan, de modo que las comunidades proporcionan los tipos y cantidades de bienes públicos deseados por sus habitantes.

Una idea estrechamente relacionada es que la mayor proximidad de un gobierno local a las personas hace que responda mejor a las preferencias de los ciudadanos que el gobierno central. Es especialmente probable que este sea el caso en un país grande donde los costes de obtener y procesar información sobre las preferencias de todo el mundo son sustanciales. El director ejecutivo de McDonald’s dijo una vez: \textit{No se pueden gestionar 25.000 restaurantes de manera centralizada. Muchas decisiones necesitan decidirse más cerca del mercado} (BARBOZA, 1999). Un sistema federal aplica el mismo principio al gobierno.

Esta lógica sugiere que cuanto más varíen las preferencias dentro de un área, mayores serán los beneficios de la toma de decisiones descentralizada dentro de esa área. Para examinar si esta idea tiene algún poder predictivo, STRUMPF y OBERHOLZER-GEE (2002) examinaron cómo difieren los estados con respecto a qué nivel de gobierno regula la venta de bebidas alcohólicas. Las personas de diferentes orígenes religiosos difieren acerca de si las bebidas alcohólicas deberían estar prohibidas. Por tanto, la teoría del federalismo sugiere que los estados con mayor diversidad religiosa deberían tener más probabilidades de descentralizar el control sobre la política regulatoria relativa al alcohol, manteniéndose iguales las demás cosas. Encontraron apoyo para esta hipótesis: el control local aumenta con la variación de las preferencias dentro del estado.

La lógica del federalismo también sugiere que las regulaciones económicas promulgadas a nivel nacional pueden no tener sentido en todas las comunidades. Por ejemplo, no tiene sentido que las regulaciones medioambientales sean uniformes en todo el país. Los costes y beneficios marginales de la reducción de la contaminación dependen de la densidad de población, de los patrones meteorológicos, etcétera. En la medida en que los funcionarios de una determinada jurisdicción tengan mejor información sobre cuestiones específicas relacionadas con su área que el gobierno federal, tiene sentido darles cierto margen para determinar la política regulatoria. En Estados Unidos, los estados pueden optar por asumir la responsabilidad de aplicar y hacer cumplir algunas políticas medioambientales federales. Existen algunas pruebas de que los estados que aprovechan esta opción son más estrictos que el gobierno federal a la hora de hacer cumplir las regulaciones (SIGMAN, 2003).

\paragraph{Competencia intergubernamental: BRENAN y BUCHANAN (1980)}
En muchos contextos, los gestores gubernamentales carecen de incentivos para producir al coste mínimo factible. Los gestores de empresas privadas que no logran minimizar los costes finalmente son expulsados del mercado. En cambio, los gestores gubernamentales pueden continuar funcionando de manera deficiente. Sin embargo, si los ciudadanos pueden elegir entre comunidades, una mala gestión considerable puede hacer que los ciudadanos simplemente se marchen. Esta amenaza puede crear incentivos para que los gestores gubernamentales produzcan de manera más eficiente y respondan mejor a sus ciudadanos. En este contexto, es interesante señalar que algunas pruebas sugieren que cuanto más descentralizado es el sistema fiscal de un país, menos probable es que su gobierno sea corrupto, manteniéndose iguales las demás cosas (FISMAN y GATTI, 2002).

\subsubsection{Rendición de cuentas: ¿Se puede incrementar la fiscalización de la acción pública mediante la descentralización?}
Otra vía para justificar la descentralización parte de los problemas inherentes a la toma de decisiones cuando existe información incompleta y, sobre todo, cuando quien administra los recursos públicos no soporta íntegramente las consecuencias de sus decisiones. Entramos, por tanto, en el terreno de la Teoría de la Agencia.

Hasta ahora hemos supuesto implícitamente que el planificador público actuaba de manera benevolente: conocía las preferencias de los ciudadanos y utilizaba la información disponible para maximizar su bienestar. Sin embargo, ¿qué ocurre si abandonamos este supuesto? Los responsables políticos son también agentes con objetivos propios. Además del bienestar de sus electores, pueden valorar la permanencia en el cargo, el prestigio, el poder o, más sencillamente, la posibilidad de apropiarse de determinadas rentas. El problema deja entonces de consistir exclusivamente en proporcionar al Gobierno información acerca de las preferencias de los ciudadanos. Aparece una segunda cuestión: ¿cómo pueden los ciudadanos saber si el Gobierno está utilizando correctamente la información que posee?

Esta cuestión proporciona , como ahora veremos, una nueva justificación para la descentralización. El trabajo seminal en esta rama de la literatura fue BRENNAN y BUCHANAN (1980) que, partiendo de su concepción del Gobierno como un Leviatán, muestran que la existencia de gobiernos competidores puede limitar la capacidad de un poder público monopolístico para apropiarse de recursos. No obstante, esta intuición puede complementarse, como ahora presentamos, con los modelos de competencia por comparación: cuando existen varias jurisdicciones comparables, la actuación de unas proporciona información con la que los ciudadanos pueden evaluar la actuación de las demás.

\paragraph{Información asimétrica: ¿Se comporta bien el Gobierno?}
Representemos esta idea mediante un juego secuencial muy sencillo. Supongamos que el Gobierno debe proporcionar un determinado servicio público y que existen dos posibles estados de la naturaleza, ($a$) y ($b$). En el estado favorable ($a$), proporcionar eficientemente el servicio cuesta una unidad, mientras que en el estado desfavorable ($b$) cuesta tres: $[c(a)=1,\, c(b)=3]$.

La naturaleza determina inicialmente cuál de los dos estados se produce. El Gobierno observa este estado, pero los ciudadanos no; existiendo, por tanto, información asimétrica. Una vez conocido el verdadero coste del servicio, el Gobierno determina los recursos que exigirá a los ciudadanos para financiarlo. Podemos representar esta cantidad como $T=c(\theta)+r$, donde ($c(\theta)$) es el coste efectivo de proporcionar el servicio y $r\geq0$ representa las rentas apropiadas por el Gobierno. Un Gobierno responsable escogerá ($r=0$); un Gobierno oportunista puede aprovechar su ventaja informativa para escoger ($r>0$).

¿Qué está en juego para cada jugador? En primer lugar, el Gobierno obtiene una utilidad derivada de las rentas que consigue apropiarse. Sin embargo, permanecer en el poder también tiene valor. Denotemos por ($V$) el valor esperado de la reelección y su función de pagos puede expresarse como $U_G=r+V\mathbb{1}(\text{reelección})$. Supondremos que ($V>r$): apropiarse de rentas resulta atractivo, pero conservar el cargo es todavía más valioso. En segundo lugar, los ciudadanos obtienen el beneficio ($B$) asociado al servicio público, pero soportan los recursos exigidos para financiarlo. Su bienestar será, por tanto, $U_C=B-T$.

La dificultad aparece porque los ciudadanos observan cuánto han tenido que pagar, pero no conocen directamente las circunstancias a las que se enfrentaba el Gobierno. Supongamos, por ejemplo, que observan ($T=3$). Esta observación es compatible con dos historias completamente diferentes. Puede haberse producido el estado desfavorable (b), en cuyo caso un Gobierno responsable necesitaba efectivamente tres unidades para proporcionar el servicio: $(b,r=0)\longrightarrow T=3$. Pero también puede haberse producido el estado favorable (a) y haber aprovechado el Gobierno su ventaja informativa para apropiarse de dos unidades: $(a,r=2)\longrightarrow T=3$. El resultado observado por los ciudadanos es exactamente el mismo. Lo que cambia es aquello que no pueden observar: el estado de la naturaleza y, con él, la existencia o no de rentas políticas.

Aquí reside el problema fundamental de un Gobierno monopolístico (centralizado). Los ciudadanos pueden observar sus resultados, pero carecen de un término de comparación que les permita determinar qué parte de esos resultados es atribuible a circunstancias objetivas y qué parte responde al comportamiento del gobernante. En términos de teoría de juegos, ambos nodos pertenecen al mismo conjunto de información del electorado. 

Esto limita considerablemente la eficacia de las elecciones como mecanismo disciplinario. Si los votantes establecen un criterio muy exigente de reelección, pueden castigar a un Gobierno responsable simplemente por haber tenido que actuar bajo circunstancias adversas. Si, por el contrario, relajan ese criterio para evitar dicho error, proporcionan margen al gobernante para apropiarse de rentas cuando las circunstancias son favorables.

\paragraph{Competencia por comparación: ¿Qué información proporcionan otros gobiernos del mismo nivel?}
La descentralización puede alterar precisamente este problema informativo. Supongamos ahora que existen dos jurisdicciones, ($i$) y ($j$), sometidas a condiciones económicas iguales o suficientemente correlacionadas, pero gobernadas de manera independiente. Los ciudadanos de cada jurisdicción pueden observar no sólo el coste que soportan en su propia jurisdicción, sino también el resultado obtenido en la otra. Consideremos nuevamente que los ciudadanos de ($i$) observan ($T_i=3$). Aislada, esta información continúa siendo ambigua. Sin embargo, supongamos que simultáneamente observan $T_j=1$

La actuación de la segunda jurisdicción proporciona ahora una información que antes no existía. Si ambas administraciones se enfrentan a circunstancias comparables y el Gobierno de ($j$) ha podido proporcionar el mismo servicio exigiendo únicamente una unidad, resulta mucho menos verosímil que las tres unidades exigidas por el Gobierno de ($i$) respondan exclusivamente a unas condiciones económicas desfavorables. La jurisdicción ($j$) se convierte, por tanto, en un benchmark con el que evaluar a ($i$). 

Obsérvese que la descentralización no tiene por qué hacer que los políticos sean más benevolentes. Tampoco elimina necesariamente su capacidad para apropiarse de rentas. Lo que modifica es el conjunto de información de los ciudadanos. La coexistencia de varios gobiernos genera observaciones adicionales que permiten distinguir mejor entre un mal resultado provocado por circunstancias adversas y un mal resultado provocado por una mala actuación del gobernante.

Esta posibilidad modifica, a su vez, los incentivos del político. Bajo un sistema centralizado, la asimetría informativa puede permitirle apropiarse de una renta ($r$) sin que los ciudadanos puedan identificar con precisión su comportamiento. Bajo un sistema descentralizado, la comparación con otras jurisdicciones eleva la probabilidad de que esa conducta sea detectada y castigada electoralmente. Si el valor de permanecer en el cargo ($V$) es suficientemente elevado en relación con las rentas que puede obtener, el gobernante tendrá incentivos para renunciar a estas últimas y actuar de acuerdo con el interés de sus electores.

La descentralización introduce así una forma de competencia que no depende necesariamente de que los ciudadanos se desplacen de una jurisdicción a otra; \textbf{es la propia información la que se desplaza}: los resultados de otras administraciones proporcionan a los votantes una vara de medir con la que evaluar a la suya. Por esta razón, la existencia de múltiples gobiernos puede aumentar la rendición de cuentas. Frente a un Gobierno central monopolístico, cuya actuación debe evaluarse de forma aislada, una estructura descentralizada genera términos de comparación. La descentralización puede justificarse entonces no sólo porque acerca las decisiones públicas a unas preferencias territorialmente heterogéneas, sino porque reduce la ventaja informativa del agente frente al principal y, con ello, limita su capacidad para apropiarse de rentas.

\textcolor{red}{Sin embargo, dentro de las aportaciones de la segunda generación de la Teoría del Federalismo fiscal, existen otros enfoques relacionados con la rendición de cuentas considerando también el problema principal-agente. Aquí por ejemplo podemos citar las aportaciones de TOMASI en el año 2003, donde los votantes tienen menos información y actúan como principal. En el caso del agente, el autor considera dos opciones: el sistema centralizado, donde el gobierno central es un solo agente que sirve a la población; y el sistema descentralizado, donde existe un agente por jurisdicción, cada gobierno local. TOMASI (2003) desarrolló un modelo en el que el sistema centralizado es un solo agente que sirve a la población y el descentralizado considera un agente por jurisdicción. Analizaba cuál es el contrato óptimo que incentive a los agentes a satisfacer las preferencias del electorado. Mostró así que cuantas mayores sean las externalidades asociadas a los bienes públicos locales más preferible es establecer un sistema centralizado. Como vemos, esta conclusión es similar a la del enfoque que siguen las aportaciones anteriores. Sin embargo, es importante destacar una diferencia clave. En el análisis de TOMASI, la descentralización puede ser preferible incluso cuando las preferencias de los agentes son idénticas para todas las jurisdicciones. En definitiva, elegir la descentralización no solo dependería de las diferencias en las preferencias como tradicionalmente se había considerado, sino también de un potencial mayor control o una mayor rendición de cuentas por parte de los gobiernos locales. Además, las teorías de la primera generación consideraban que los problemas de información asimétrica también eran una justificación tradicional para la descentralización, dado que las haciendas territoriales tienen mayor acceso a la información sobre las preferencias de los agentes. Se asumía así que el gobierno centralizado sí conoce las preferencias de los agentes para los bienes públicos nacionales. Pero autores como CREMER, ESTACHE y SEABRIGHT (1996) argumentaron que en esa afirmación existe una dicotomía ya que la adquisición de información es endógena. De esta forma, el gobierno central tiene la posibilidad de establecer mecanismos para conocer las condiciones en todas las regiones, aunque sean costoso.}

\subsubsection{Riesgo}





\paragraph{Experimentación e innovación: Provisión fragmentada}
\textbf{Experimentación e innovación en bienes y servicios proporcionados localmente} 
Para muchas cuestiones de política, nadie sabe cuál es la respuesta correcta, ni siquiera si una única solución es la mejor en todas las situaciones. Una forma de averiguarlo es permitir que cada comunidad elija su propio camino y después comparar los resultados. Un sistema de gobiernos diversos aumenta las posibilidades de que se busquen nuevas soluciones a los problemas. Como observó una vez el juez del Tribunal Supremo Louis Brandeis: \textit{Es uno de los felices accidentes del sistema federal que un único estado valiente pueda, si sus ciudadanos así lo eligen, servir como laboratorio y probar experimentos morales, sociales y económicos sin riesgo para el resto del país}.

A juzgar por todas las apariencias, los laboratorios de Brandeis están trabajando intensamente. Históricamente, algunos programas que comenzaron como experimentos a nivel estatal finalmente se convirtieron en política federal. Durante la Gran Depresión, por ejemplo, los diseñadores de la Seguridad Social aprovecharon la experiencia de varios estados que anteriormente habían establecido programas de seguro social.

\subsection{Argumentos en contra: }
Las dos vías propuestas para justificar la descentralización resultan muy interesantes y, hasta cierto punto, simplifican adecuadamente algunos aspectos de la realidad. Sin embargo, en muchos casos estas simplificaciones exigen imponer unos supuestos demasiado restrictivos que, de relajarlos, pueden llegar a contradecir las implicaciones anteriores. 

Veamos una serie de problemas que presentan las justificaciones anteriores y, en cierto modo, justifican la centralización de la acción pública.

\subsubsection{Competencia entre regiones: ¿Se observa realmente en la práctica?}
El modelo de TIEBOUT requiere una serie de supuestos que pueden no cumplirse en la realidad. 

\paragraph{Movilidad perfecta: }
El primer supuesto es la \textbf{movilidad perfecta}: los individuos no solo deben querer votar con los pies, sino que deben ser capaces de llevar realmente a cabo ese voto. Esto es difícil en la práctica. Por ejemplo, puede estar muy asentado en una localidad concreta, con muchos amigos y otras comodidades. Haría falta mucho más que la financiación inapropiada del coche deportivo de un director de una biblioteca para conseguir que alguien se mudara.

\paragraph{Información perfecta}
Quizá aún más inverosímil sea el supuesto de que los individuos tienen \textbf{información perfecta} sobre los beneficios que reciben del municipio y los impuestos que pagan. Incluso si la biblioteca pública local estuviera comprando coches deportivos a todos sus bibliotecarios, yo nunca me enteraría a menos que de algún modo fuera revelado por los medios de comunicación locales (y yo estuviera prestando atención).

\paragraph{Libertad de elección: }
Además, para que se cumpla el modelo de TIEBOUT, debo poder \textbf{elegir libremente} entre una variedad de municipios que puedan ajustarse a mi gusto por los bienes públicos. Esta variedad puede existir alrededor de un municipio grande. Pero podría no ser así en otras áreas, donde los municipios están más dispersos y votar con los pies significaría mudarme considerablemente más lejos. Tales restricciones sobre los sustitutos adecuados para el municipio de uno podrían limitar la utilidad del mecanismo de TIEBOUT para áreas metropolitanas más pequeñas o en declive.

Al mismo tiempo, el modelo de TIEBOUT requiere que haya \textbf{suficientes municipios} para que los individuos puedan clasificarse en grupos con preferencias similares por los bienes públicos. Esto plantea una clara tensión: ¿Podemos dividir a la población en grupos de personas que tengan todas preferencias similares por los bienes públicos y, al mismo tiempo, garantizar que estos grupos sean lo suficientemente grandes como para sostener las economías de escala requeridas por los bienes públicos?


\subsubsection{Eficiencia: ¿No existen externalidades ni fallos de mercado?}

Finalmente, la provisión de algunos bienes públicos requiere una \textbf{escala o tamaño suficientes}. No es eficiente gestionar una escuela con solo unos pocos estudiantes ni construir un parque que vaya a ser utilizado únicamente por unos pocos residentes debido a los grandes costes fijos de construir la escuela o el parque. Estos costes fijos dan lugar a econocomías de escala, por las cuales la eficiencia de un bien público es mucho mayor si es utilizado por muchos en lugar de por pocos. Una escuela que es utilizada por mil estudiantes puede financiarse mediante impuestos menores que una escuela utilizada por diez estudiantes, porque los grandes costes fijos de la escolarización (por ejemplo, el edificio, el director) pueden repartirse entre el conjunto más amplio de hogares.

\textbf{Ausencia de externalidades/efectos de desbordamiento} 
Un tercer problema con el modelo de TIEBOUT es que supone que los bienes públicos tienen efectos únicamente en un municipio determinado y que los efectos no se desbordan hacia los municipios vecinos. Si existen tales efectos de desbordamiento, hay argumentos a favor de la provisión de bienes públicos en un nivel superior de gobierno, o de subvenciones que subsidien las compras locales.

Imaginemos que Barcelona está considerando construir un gran parque público nuevo. Este parque será disfrutado principalmente por sus habitantes, pero muchas personas de municipios vecinos también visitarán sus hermosos terrenos. Bajo el mecanismo de TIEBOUT, cuando Barcelona decida si construir el parque, considerará únicamente las preferencias de sus residentes, no las preferencias de los residentes de otros municipios que podrían disfrutar del parque. Por tanto, nos enfrentamos al problema estándar de la provisión de bienes públicos: dado que las personas de otros municipios actúan como free riders sobre el parque, se suministrará una cantidad insuficiente de servicios. Si los beneficios sociales (para Barcelona y su área metropolitana) superan el coste de construir el parque, este debería construirse, pero si los beneficios privados son menores que los costes de construir el parque, entonces no se construirá, lo cual es socialmente ineficiente.

Muchos bienes públicos locales presentan características similares de externalidad o efecto de desbordamiento: la policía (si no es lo suficientemente grande, la actividad delictiva podría desbordarse hacia otros municipios); las obras públicas (si las calles están cubiertas de baches, los conductores de municipios vecinos podrían sufrir las consecuencias al atravesar mi municipio); la educación (toda la nación se beneficia de una ciudadanía más educada); etc. Por tanto, existe una disyuntiva fundamental en el enfoque de TIEBOUT. Existen ventajas en los bienes públicos proporcionados localmente para municipios pequeños de individuos similares, pero puede ser óptimo proporcionar los bienes públicos que tienen efectos externos o efectos de desbordamiento sobre otros municipios en un nivel superior de gobierno que pueda internalizar las externalidades.

\subsubsection{Conflicto interdimensional: ¿Es justo el resultado final?}
Un segundo gran problema con el funcionamiento del modelo de TIEBOUT es que requiere una \textbf{financiación igualitaria} del bien público entre todos los residentes.\footnote{Este tipo de financiación se denomina impuesto de suma fija, una suma fija que una persona paga en impuestos independientemente de los ingresos de esa persona, su consumo de bienes y servicios o su riqueza.
} Como sabemos, esta forma de imposición es considerada altamente inequitativa por el público porque tanto ricos como pobres pagan la misma cantidad de impuestos. Como resultado, la imposición de suma fija se utiliza muy raramente para financiar los gastos gubernamentales. De hecho, el intento de mayor repercusión de imponer impuestos de suma fija, por parte del gobierno británico de Margaret Thatcher en 1990, dio lugar a grandes disturbios que condujeron a la dimisión de la que previamente había sido increíblemente popular.

Los municipios suelen financiar sus bienes públicos, en cambio, mediante impuestos (normalmente) progresivos. Esto genera otro problema: los pobres persiguen a los ricos. Las personas más ricas pagan una parte mayor de la factura de los bienes públicos que las personas más pobres, por lo que las personas que valoran esos bienes querrían vivir en una comunidad con personas más ricas que ellas. De ese modo, las personas más pobres pueden beneficiarse de los mayores impuestos pagados por sus vecinos más ricos. En otras palabras, todo el mundo quiere vivir en municipios con personas que sean más ricas que ellos para poder actuar como \textit{free riders} sobre los mayores pagos de impuestos de sus vecinos.

Una forma en que los municipios han tratado de resolver este problema es mediante el uso de la zonificación. Las regulaciones de zonificación son restricciones que los municipios imponen sobre cómo puede utilizarse la propiedad inmobiliaria, aparentemente con el objetivo de preservar el carácter de la comunidad local. Por ejemplo, una regulación de zonificación común exige que las casas se construyan a una determinada distancia de la calle para preservar cierto espacio de jardín y, por tanto, el carácter estético del vecindario. Otros ejemplos de regulaciones de zonificación incluyen prohibiciones contra el uso de la propia vivienda para gestionar un negocio en un vecindario residencial, restricciones sobre el número máximo de ocupantes que puede albergar un terreno o edificio, requisitos de tamaños mínimos de parcela, limitaciones sobre el tamaño máximo de los edificios y prohibiciones de viviendas multifamiliares. Las regulaciones de zonificación protegen la base impositiva de los municipios ricos al excluir del mercado de la vivienda, mediante los precios, a las personas de menores ingresos. Por ejemplo, un municipio que prohíbe las viviendas multifamiliares (como las casas para dos familias y los edificios de apartamentos) reduce la cantidad disponible de viviendas y, por tanto, infla el valor de las viviendas existentes de modo que las personas pobres no puedan permitirse mudarse allí y actuar como \textit{free riders} sobre los pagos de impuestos de los vecinos con mayores ingresos. De hecho, GLAESER y GYOURKO (2002) compararon áreas con diferentes leyes de zonificación y encontraron que los precios del suelo en las áreas zonificadas son diez veces más altos que los precios en los mercados no zonificados.

Las comunidades pueden promulgar leyes de zonificación excluyentes, leyes que prohíben determinados usos del suelo. En concreto, pueden exigir que todas las viviendas tengan un determinado tamaño mínimo. Para ver por qué este supuesto es crucial, recordemos que, en el equilibrio de TIEBOUT, las comunidades están segregadas sobre la base de las demandas de bienes públicos de sus miembros. Si la renta está positivamente correlacionada con la demanda de servicios públicos, el resultado es la segregación de las comunidades por renta. En las comunidades de renta alta, el nivel de los valores de las propiedades tiende a ser elevado y, por tanto, la comunidad puede financiar una determinada cantidad de gasto público con un tipo del impuesto sobre la propiedad relativamente bajo. Las familias de renta baja tienen un incentivo para trasladarse a comunidades ricas y construir viviendas relativamente pequeñas. Debido al bajo tipo impositivo, las familias de renta baja tienen obligaciones fiscales relativamente pequeñas, pero, no obstante, disfrutan del elevado nivel de provisión de servicios públicos. A medida que más familias de renta baja tienen la misma idea y se trasladan allí, la base imponible por familia de la comunidad disminuye. Los tipos impositivos deben aumentarse para financiar el nivel ampliado de servicios públicos necesario para atender al aumento de la población. Puesto que suponemos movilidad perfecta, los ricos no tienen ninguna razón para soportar esto. Simplemente se trasladan a otra comunidad. Pero ¿qué impide que los pobres los sigan? En ausencia de restricciones a la movilidad, nada. Claramente, en un modelo de Tiebout puede desarrollarse un juego de las sillas de suburbios. La zonificación excluyente evita este fenómeno y mantiene así un equilibrio estable y eficiente en el sentido de Pareto.










\clearpage
\section{Gasto: ¿Cómo asignar las partidas de gasto?}
\puntosclave{
    
}



La clave está en cambiar ligeramente la pregunta. No sería tanto, ¿qué partidas de gasto debe realizar cada nivel de gobierno?, como, dada la existencia de distintos niveles de gobierno, ¿a qué nivel debe asignarse cada responsabilidad de gasto y bajo qué forma debe ejercerse?\footnote{Eso es exactamente lo que la literatura denomina asignación de responsabilidades de gasto. Se puede profundizar en el capítulo corresopndiente de AHMAD, HEWITT y RUGGIERO en \textit{Fiscal Federalism in Theory and Practice}.
}

La pregunta fundamental: ¿cuál es el tamaño óptimo de la jurisdicción? Una vez que aceptamos que pueden coexistir distintos niveles de gobierno, la cuestión ya no es decidir simplemente entre centralización o descentralización sino en asignar el \textbf{nivel óptimo de gasto asignado para cada nivel de gobierno}. Para empezar, sabemos que la literatura proporciona un principio de partida muy potente que OLSON (1969) denominó principio de equivalencia fiscal: idealmente, debe existir correspondencia entre el territorio sobre el que recaen los costes de una política y el territorio sobre el que recaen sus beneficios.\footnote{Esta idea sigue apareciendo en formulaciones modernas. La OCDE, por ejemplo, plantea que una función debería asignarse al nivel territorial cuya extensión geográfica abarque, en la medida de lo posible, tanto a quienes se benefician del servicio como a quienes soportan su coste.  
}

Lo primero que debemos tener en cuenta es el principio de equivalencia fiscal y correspondencia territorial de costes y beneficios. De esta forma, el problema entero queda formulada de una manera muy sencilla: consiste encontrar la asignación óptima de gasto dada la heterogeneidad de preferencias, las externalidades presentas, las economías de escala y la capacidad administrativa.

Estas cuatro fuerzas empujan en direcciones distintas:

Característica de la función	Favorece
Preferencias territorialmente heterogéneas	Descentralización
Beneficios concentrados territorialmente	Descentralización
Externalidades interjurisdiccionales elevadas	Centralización/coordinación
Economías de escala elevadas	Jurisdicciones mayores
Necesidad de información local	Descentralización
Elevada complejidad/capacidad técnica necesaria	Puede favorecer niveles superiores
Economías de alcance entre servicios	Gestión conjunta


Referencias que utilizaría como columna vertebral

No necesitas cincuenta artículos. Para construir este apartado con rigor, empezaría por estas pocas referencias: Oates (1972), Fiscal Federalism como punto de partida de la asignación eficiente; Olson (1969), “The Principle of Fiscal Equivalence” para la correspondencia entre jurisdicción, beneficiarios y costes; Oates (1999), “An Essay on Fiscal Federalism” como revisión canónica de la literatura; Ahmad, Hewitt y Ruggiero (1997), “Assigning Expenditure Responsibilities”, porque trata exactamente el apartado que estás construyendo; Alesina y Spolaore (1997), “On the Number and Size of Nations”, para formalizar escala frente a heterogeneidad; Alesina, Baqir y Hoxby (2004) para evidencia empírica sobre jurisdicciones, escala y heterogeneidad; Besley y Coate (2003) como revisión político-económica del problema centralización/descentralización con heterogeneidad y spillovers; Allain-Dupré (2018, OECD) para el paso de la teoría a criterios modernos de asignación institucional; y, si quieres específicamente el componente Hotelling, Hsu (2005), “A Hotelling Model of Fiscal Competition”.  

Y entre todas ellas, para lo que estás haciendo empezaría leyendo Ahmad et al. (1997), Olson (1969), Oates (1999) y Alesina-Spolaore (1997). Son las que más directamente te van a permitir construir el discurso.

Creo, además, que aquí tienes material para hacer algo bastante más interesante que una exposición puramente descriptiva: podemos construir todo el apartado 2 alrededor de un pequeño modelo gráfico-matemático común, de modo que heterogeneidad, externalidades y economías de escala no aparezcan como tres argumentos inconexos, sino como fuerzas que conjuntamente determinan g_j^*. Ese sería el siguiente desarrollo que yo haría: definir formalmente H(g),E(g),S(g), obtener la jurisdicción óptima y después aplicarlo sucesivamente a educación, sanidad, transporte, policía/bomberos y defensa. Ahí podemos conseguir un apartado muy redondo sin meter matemáticas artificiales.


\subsection{Principio de correspondencia: Benficiarios y contribuyentes}
Recordemos cuál era la pregunta fundamental de TIEBOUT: \textit{¿Por qué podría ser deseable descentralizar?} Formulemos ahora una alternativa ¿Qué funciones concretas deberían descentralizarse? La respuesta inicial, será que cuanto más territorialmente concentrados estén los beneficios de un bien y mayor sea la heterogeneidad de preferencias entre territorios, menor debería ser, ceteris paribus, la jurisdicción encargada de decidir su provisión.

Ese es el corazón del Teorema de Descentralización de OATES. Bajo sus supuestos, cuando las preferencias difieren entre jurisdicciones y no existen economías de escala ni externalidades relevantes, permitir diferentes niveles de provisión local produce al menos tanto bienestar como imponer una provisión uniforme centralizada.  

Supongamos que existen dos municipios, $A$ y $B$, con preferencias diferentes. Si un gobierno central proporciona una cantidad uniforme en ambos municipios, aparece una pérdida por falta de adaptación, $L_i=L(q_i^*-\bar q)$. Cuanto mayor sea la pérdida, mayor es el coste de imponer una política común.



\subsubsection{Diseño óptimo de políticas: }
Veamos las aportaciones de ALESINA y SAPOLAORE.\footnote{Por ejemplo, HSU (2005) construye un modelo en el que dos jurisdicciones compiten mediante tipos impositivos y mediante el «tipo» de bien público local que ofrecen, estudiando hasta qué punto terminan diferenciándose.  
}

Su modelo enfrenta directamente:

\underbrace{\text{beneficios de jurisdicciones grandes}}_{\text{economías de escala}}

con

\underbrace{\text{costes de jurisdicciones grandes}}_{\text{heterogeneidad de preferencias}}.

La lógica puede reducirse pedagógicamente a:

\text{Coste total por habitante}=\text{coste de provisión}+\text{coste de heterogeneidad}.

Si aumenta el tamaño de la jurisdicción, disminuye el coste por habitante porque repartes los costes fijos entre más personas. Pero simultáneamente aumenta, normalmente, H(N)\uparrow porque resulta más difícil ofrecer una política que satisfaga las preferencias de una población cada vez más heterogénea. Por tanto: AC(N)=\frac{F}{N}+H(N).

Y existe potencialmente:

N^*=\arg\min_N
\left\{
\frac{F}{N}+H(N)
\right\}.

Esto me parece muchísimo más apropiado para tu exposición que Hotelling. Alesina y Spolaore desarrollan precisamente el trade-off entre las economías derivadas del tamaño y los costes ocasionados por la heterogeneidad.  

Además tienes una prolongación empírica muy bonita: Alesina, Baqir y Hoxby (2004) estudian municipios, distritos escolares y otros tipos de jurisdicciones estadounidenses y encuentran evidencia de ese intercambio entre economías de escala y determinadas formas de heterogeneidad de la población.  

Yo usaría Hotelling después, si quieres introducir diferenciación horizontal entre bienes públicos, pero no para fundamentar todo el problema.












⸻

11. Esto te permite construir un modelo conceptual propio muy sencillo

Yo haría esto.

Para cada función pública j y nivel de gobierno g, defines:

SC_{jg}
=
H_{jg}
+
E_{jg}
+
S_{jg}
+
A_{jg}.

Donde:

H_{jg}
=
\text{coste por falta de adaptación a preferencias};

E_{jg}
=
\text{coste por externalidades no internalizadas};

S_{jg}
=
\text{pérdida de economías de escala};

A_{jg}
=
\text{costes administrativos y de coordinación}.

Entonces:

\boxed{
g_j^*
=
\arg\min_g SC_{jg}
}

No presentaría esto como «un modelo de Oates» ni se lo atribuiría a ningún autor concreto. Lo presentaría transparentemente como una síntesis analítica de los criterios de asignación que acabamos de estudiar.

Y puede funcionar maravillosamente en clase.

Porque puedes coger cualquier partida.

Parques urbanos

H_L\downarrow,\quad
E_L\approx0,\quad
S_L\approx0

\Rightarrow\text{local}.

Transporte metropolitano

Municipio:

E_L\uparrow,\quad S_L\uparrow

Región metropolitana:

E_M\downarrow,\quad S_M\downarrow

\Rightarrow\text{metropolitano/regional}.

Defensa nacional

E_L\rightarrow\text{enorme},
\qquad
S_L\rightarrow\text{enorme}

\Rightarrow\text{central}.

De repente el alumno puede aplicar el federalismo fiscal, no solamente memorizarlo.

⸻

12. Y ahora una cuestión importante: ¿qué haría con la equidad?

Aquí hay que tener cuidado porque tú ya has decidido tratarla en el apartado 1 como límite de la descentralización, y creo que está bien.

No volvería a desarrollar toda la teoría redistributiva aquí.

Pero sí introduciría una advertencia.

Hay ciertas funciones —educación, sanidad, servicios sociales— para las que la decisión de asignación no depende exclusivamente de eficiencia espacial.

Puede ocurrir que:

g_j^*=\text{local}

desde el punto de vista puramente asignativo, pero que el Estado quiera garantizar:

q_i\geq\underline q
\qquad\forall i

por razones de equidad.

Entonces aparece una solución multinivel:

\boxed{
\text{estándar central}
+
\text{prestación descentralizada}
}

y posiblemente:

+
\text{financiación/transferencias centrales}.

Esto es especialmente útil porque anticipa el apartado 3 sin repetirlo.

⸻

13. Tu apartado 2 podría quedar estructurado así

Yo propondría una estructura bastante concreta:

1. El problema de la asignación del gasto público. Introduces expenditure assignment, subsidiariedad y equivalencia fiscal de Olson: cada función debería corresponder, en principio, a la jurisdicción mínima capaz de internalizar sus costes y beneficios.  
2. Heterogeneidad de preferencias y adaptación de la provisión. Recuperas brevemente Oates y formalizas el coste de uniformidad. Aquí puedes introducir, como extensión, diferenciación espacial/Hotelling, aunque Alesina-Spolaore me parece una base mejor.  
3. Externalidades y área óptima de provisión. Presentas el problema de spillovers, la discrepancia entre beneficio privado/local y beneficio social y las posibles soluciones: centralización, coordinación o transferencias.
4. Economías de escala y tamaño óptimo de la jurisdicción. Introduces costes fijos y el trade-off escala-heterogeneidad de Alesina-Spolaore.  
5. Producción, provisión, financiación y regulación. Rompes la idea de que asignar una función significa entregársela íntegramente a un solo nivel. Ésta sería para mí una sección fundamental.  
6. Coordinación intergubernamental y soluciones intermedias. Consorcios, áreas metropolitanas, cooperación horizontal, estándares mínimos, etc. La literatura y la práctica contemporáneas insisten mucho en estas soluciones.  
7. Una regla de síntesis. Presentas g_j^*=\arg\min_g\{H+E+S+A\} y aplicas el criterio a 4–5 políticas públicas.

Creo que así queda mucho más fuerte que dividir simplemente en externalidades/economías de escala/Hotelling.


-------------------------------


Ahora que conocemos todas las razones que permiten argumentar acerca de las bondades y defectos que caracterizan a la descentralización fiscal, es necesario conocer las herramientas que nos permiten analizar su optimalidad material. Dicho de otra forma, suponiendo que trabajamos sobre un nivel de descentralización dado, ¿cómo podemos saber si este es óptimo?

Para resolver esta cuestión de forma clara y didáctica, dividiremos lo que queda de exposición en dos partes. A continuación presentamos qué factores influyen desde la óptica del gasto y, de esta forma, una vez tengamos claro cómo analizar la provisión de bienes públicos para un determinado nivel de Gobierno, podremos luego entender las diferentes figuras que permitan financiarlo.

Como se comentó, las aportaciones realizadas por la primera generación de la Teoría del Federalismo Fiscal respecto al análisis de la descentralización del gasto, se centraban principalmente en la asignación de funciones a los distintos niveles de gobierno. En este sentido cabe destacar las aportaciones de MUSGRAVE en el ai10 1 959 donde considera un rol activo por parte del Sector Público por razones de:
• Estabilización: de cara a conseguir una senda de crecimiento estable y reducir las fluctuaciones cíclicas de la economía.
• Equidad: para conseguir una redistribución más equitativa de la renta.
• Eficiencia: con el objetivo <le corregir fallos de mercado y alcanzar una asignación eficiente. Se procede, por tanto, a analizar la descentralización desde cada una de estas ópticas.

Aunque el modelo de TIEBOUT es una descripción imperfecta de la realidad, los cambios en la imposición y el gasto locales sí afectan a la movilidad y a los precios de las viviendas. Dados estos resultados positivos, ¿cuáles son las implicaciones normativas del modelo de TIEBOUT para el diseño óptimo del federalismo fiscal? Es decir, ¿qué implica el modelo de TIEBOUT que deberían ser los principios que guíen la provisión de bienes públicos en diferentes niveles de gobierno? El modelo de TIEBOUT implica que el grado en que los bienes públicos deberían proporcionarse a nivel local está determinado por tres factores. 


\subsection{Percepción: Vínculo entre impuestos y beneficios}
El primero son los \textbf{vínculos entre impuestos y beneficios}, el grado en que los residentes consideran que sus pagos de impuestos están directamente vinculados a los bienes y servicios que reciben. Los bienes con fuertes vínculos entre impuestos y beneficios, como las carreteras locales, deberían proporcionarse localmente. Existe un vínculo directo entre impuestos y beneficios en el gasto en carreteras locales: unos impuestos sobre la propiedad más altos financian carreteras de mejor calidad que benefician a la mayoría de los residentes de un municipio. Los bienes con vínculos más débiles entre impuestos y beneficios, como los pagos de asistencia social a los residentes de menores ingresos de un municipio, deberían proporcionarse a nivel estatal o federal. Existe un vínculo muy limitado entre impuestos y beneficios en el gasto en asistencia social: la mayoría de los residentes de un municipio no se benefician de la redistribución hacia los grupos de bajos ingresos (a menos que tengan preferencias altruistas hacia los pobres locales).

Si los residentes pueden ver directamente los beneficios que están comprando con el dinero de sus impuestos, estarán dispuestos a pagar impuestos locales. Si no pueden ver un beneficio derivado de sus pagos de impuestos, votarán con los pies trasladándose a un municipio que tenga impuestos más bajos. Si un municipio instaurara un programa de asistencia social en efectivo, los residentes con mayores ingresos tendrían un incentivo para marcharse y trasladarse a un municipio que no tuviera tal programa y que, como resultado, tuviera impuestos locales sobre la propiedad más bajos. 

La capacidad de los individuos para votar con los pies, por tanto, constituye una limitación fundamental de la capacidad de un municipio para llevar a cabo programas que beneficien únicamente a una minoría de los residentes.



\subsection{Diseño óptimo: ¿cómo definir las competencias de cada nivel de Gobierno?}


\textcolor{red}{\textbf{Origen:} ICEX-CECO}
JJ.l.llJ. Estructu ra óp tima

Sin embargo, se debe considerar la posible existencia de efectos de desbordamiento sobre agentr::s de otras regiones cercanas (por ej emplo, en sanidad local, educación. etc.).

Estas externalidades se pueden internalizar a través de distintas maneras, como introduciendo impuestos pigouvianos, o mediante la coordinación de políticas entre las regiones (aunque esta coordinación es costosa4, y por lo tanto la descentralización pueda resultar ineficiente).

En la práctica, de cara a definir la estructura óptima de gobierno. se debe tener también en cuenta el trade-off existente entre las ganancias de eficiencia de crear múltiples niveles de gobierno local, y los costes de gestión que ello conlleva.

La definición de la estructura del sistema descentralizado depende de las magnitudes relativas de costes y beneficios asociados. No obstante, siempre suele resultar un sistema más centralizado que el que habría si no existiera coste alguno para establecer distintos niveles de gobierno y una mayor estratificación.

Siguiendo a Hindrícks y Myles, este argumento se puede ilustrar con un modelo espacial donde se consideran las economías de escala y la diversidad de las preferencias. Consideran un segmento [O, 1 ] donde se representa la nación. y donde cada punto representa una localización geográfica en la que se sitúan los individuos y un hien público. La distribución de los agentes sobre el plano es unifonne.

Modelo baseline

Suponen que es el Sector Público el que provee el bien público. con un coste de provisión que llamaremos C. Los agentes deben afrontar este coste de provisión entre todos. lo que les genera desutilidad neta.

El bien público puede situarse en cualquier punto del segmento, de manera que la utilidad de un individuo será máxima si el bien se sitúa en su localización, reduciéndose de acuerdo con la distancia que exista entre la localización del bien y la localización del individuo. En este sentido, llamaremos "a'' a la tasa a la que se reduce la utilidad de los agentes según esta distancia. Se puede representar la nación en el espacio (O, 1 ] como se obse1va en el Gráfico 2.

\textbf{GRÁFICO HOTELING}
o 1

Fuente: Intermediate Public Economics, Jean Hindricks y Gareth D.
Mylcs.

Se considera aquí un sistema centralizado en el que se tendrá una provisión uniforme que no refleja la heterogeneidad de las preferencias. De esta manera, el bien público se localizará.a m itad del segmento, en el punto ½- Puesto que la provisión supone un coste C, la utilidad neta cada uno de los individuos en el sistema centralizado viene dada en función de la distancia en valor absoluto entre el bien público y el individuo, la tasa ·'a .. y los costes de provisión:
e
1
1
· 1
U¡ = (1 - a z - l ) - C
i: punto en el que se localiza el individuo 

Pero en el caso de un sistema descentralizado, podemos considerar dos jurisdicciones idénticas en nuestro espacio: la región L y la región R como aparece en el Gráfico 3.

\textbf{GRÁFICO HOTELING: 2}

Para cada una de las regiones. la prov1s1on del bien se dará a la mitad de cada uno de los segmentos. en los puntos ¡ y ¡- Ahora el coste de provisión para cada jurisdicción es de C. con lo que el coste total para proveer el bien público alcanzará el valor de 2C. La utilidad en este caso de los individuos de las jurisdicciones es:

R 13 vi = (l - a 4 ·1
- 1 ) - 2 c

Por lo tanto. la utilidad dependerá de la distancia media entre un individuo y el bien público. siendo ¡ en el sistema centralizado, y ¼ en cada región del sistema descentralizado. 

Pero este efecto positivo que tiene la descentralización sobre la utilidad de los agentes, debe compararse con el coste extra que supone una doble provisión del bien público. De esta manera, sólo será preferible si el coste extrn es menor al beneficio que se obtiene al reducir la distancia media de �� en el sistema centralizado a .: en el descentralizado. Así. en función de la tasa '"a":
4 8
C < a [-1 - -1] �� C < -a
- 4 8 - 8

Por tanto, para este caso concreto con dos regiones y distribución uniforme. se puede afirmar que solo en el caso en el que se cumpla la condición anterior. la descentralización será óptima. Esto debe analizarse para otros casos.

En definitiva. este modelo muestra que el nivel óptimo de estratificación se alcanzará cuando los beneficios de una mayor descentralización ( que se dan por el hecho de satisfacer en mayor medida las preferencias heterogéneas) compensen el coste de una provisión diferenciada.

Pero en la práctica. el grado de descentralización se determina a través de procesos políticos. Es interesante extender el modelo para ilustrar cómo la votación por mayoría da lugar a una descentralización exce��iva. bajo los supuestos del modelo.

Introducción de votación

Los autores asumen en este caso, que la provisión descentralizada prevalece cuando la mayoría de votantes son favorables a ello al menos en una región. Asumen que existe simetría entre las regiones, de manera que. si existe esta mayoría en una región. existirá también en el resto. De esta forma. nos centraremos en la región L.

La mayoría en la jurisdicción L estará formada por individuos tanto a la izquierda como a la
derecha del punto donde se localiza el bien público local. 2:.. Se puede ver fácilmente que todos
4
los individuos a la izquierda de este punto preferirán la descentralización. ya que la distancia
media al punto donde se localiza el bien público es menor que la que habría en un sistema
centralizado. tal y como se explicó anteriormente.
Por ello, habrá una mayoría hacia la descentralización si el individuo que está situado en el punto donde se localiza el bien público, i-1, prefiere el sistema descentraliz11do. Es decir, si: 4

$U_f^c=[1-\alpha...]$

Por tanto, se concluye que la descentralización será un equilibrio en la votación por mayoría si y solo si se cumple:

\textbf{CONDCIÖN: $c\leq\frac{\alpha}{4}$}

Este resultado sugiere que habrá un excesivo grado de descentralización con la votación por mayoría, porque el coste crítico bajo este tipo de elección social será mayor que el coste óptimo que se calculó antes. De hecho. para cualquier coste C que se sitúe entre �� y ��. la votación por 8 4 mayoría resulta en una mayor descentralización ( C ::; .:!4 .) aunque no es socialmente óptima (C >
��). 8

En definitiva, con la votación por mayoría hay una descentralización excesiva: los votantes que se sitúan en los extremos tienen incentivos para votar por una provisión descentralizada ya que obtienen así un bien público más próximo a sus preferencias.

Sin embargo, el sistema democrático no i nternaliza las externalidades negativas impuestas en los votantes que se sitúan más alejados, en el centro, que deben sopo1tar un coste extra y un peor emparejamiento con sus preterencias.

\textcolor{blue}{\textbf{OJO -> Origen:} GRUBER (2017)}




\subsection{Externalidades: ¿Cómo tratar?}
El segundo factor que determina el nivel óptimo de descentralización es el \textbf{grado de externalidades positivas}, o efectos de desbordamiento, en la provisión de bienes públicos. Si los bienes públicos locales tienen grandes efectos de desbordamiento sobre otras comunidades, los bienes serán proporcionados en una cantidad insuficiente por cualquier localidad. En este caso, los niveles superiores de gobierno tienen un papel en promover la provisión de estos bienes públicos (por ejemplo, mediante subvenciones).


5. Segunda dimensión: externalidades y área de beneficio

Aquí aparece otro concepto central: benefit area.

Piensa en cuatro bienes:

\text{alumbrado}
\rightarrow
\text{metro}
\rightarrow
\text{universidad/hospital especializado}
\rightarrow
\text{defensa}.

El área geográfica sobre la que se extienden sus beneficios es progresivamente mayor.

Puedes formularlo como:

B_j(d),

donde B_j representa los beneficios del bien j a una distancia d del lugar donde se proporciona.

Para determinados bienes:

\frac{\partial B_j}{\partial d}\ll0.

Sus beneficios desaparecen rápidamente con la distancia.

Para otros:

\frac{\partial B_j}{\partial d}\approx0.

Sus beneficios tienen un ámbito territorial enorme.

Esto proporciona una regla extraordinariamente intuitiva:

La jurisdicción competente debería aproximarse al área geográfica sobre la que se extienden los beneficios y costes relevantes de la política.

Es prácticamente la equivalencia fiscal de Olson llevada a la asignación del gasto.  

Y ahora introduces externalidades.

Supón que el municipio i decide q_i, pero parte de sus beneficios recaen sobre j:

B_i(q_i)+\alpha B_j(q_i),

con

0\leq\alpha\leq1.

El gobierno local sólo internaliza:

B_i(q_i),

por lo que escogerá:

MB_i(q_i)=MC(q_i).

Pero socialmente debería cumplirse:

MB_i(q_i)+\alpha MB_j(q_i)=MC(q_i).

Así que:

q_i^{local}<q_i^{social}.

Perfecto. Acabas de mostrar matemáticamente por qué un bien con fuertes spillovers puede requerir una jurisdicción superior.

Y puedes convertir \alpha en un indicador del grado óptimo de centralización:

\alpha\uparrow
\quad\Longrightarrow\quad
\text{mayor conveniencia de coordinación/centralización}.

Esta relación está en el núcleo de la teoría tradicional de Oates y continúa siendo central en los modelos posteriores, incluyendo Besley y Coate.  

⸻

6. Pero aquí viene algo importante: externalidad ≠ centralización obligatoria

Esto te puede dar mucha profundidad.

Una de las cosas que evitaría en la exposición es:

«Si hay externalidades, centralizamos.»

Eso es demasiado fuerte.

La verdadera conclusión es:

si la jurisdicción que decide no coincide con el área afectada por la decisión, hace falta algún mecanismo que permita internalizar la externalidad.

Y hay varias soluciones institucionales.

Puedes: \text{centralizar}; pero también: \text{crear una autoridad supramunicipal}; o \text{establecer cooperación intermunicipal}; o \text{utilizar transferencias condicionadas}.

Y aquí haces por primera vez una pequeña puerta hacia tu apartado 3. Una subvención del nivel superior puede conseguir: MB_i+s=MSC sin necesidad de retirar completamente al municipio la competencia.

Esto es muy importante porque rompe una dicotomía demasiado simple: \text{centralización}\quad vs.\quad\text{descentralización}.

En la realidad podemos tener: \boxed{\text{decisión local}+\text{financiación central}+\text{producción compartida}}







\subsection{Economías de escala: ¿?}
El tercer factor que determina el nivel óptimo de descentralización es la \textbf{economía de escala} en la naturaleza de los bienes públicos. Los bienes públicos que tienen grandes economías de escala, como la defensa nacional, no son proporcionados eficientemente por muchas jurisdicciones locales que compiten entre sí; los bienes públicos sin grandes economías de escala, como la protección policial, pueden ser proporcionados de manera más eficaz mediante la competencia de TIEBOUT.


8. Tercera dimensión: economías de escala

Aquí tu planteamiento es totalmente correcto.

Supón:

C(q)=F+cq,

donde F es un coste fijo.

Con N usuarios:

AC(N)=\frac{F}{N}+c.

Entonces:

\frac{\partial AC}{\partial N}<0.

Cuanto mayor sea la población atendida, menor es el coste medio.

Tienes entonces el argumento clásico contra una descentralización excesiva: fragmentar la producción puede obligar a duplicar costes fijos.

Si tienes diez municipios y cada uno necesita:

* sistema informático,
* departamento de contratación,
* planta de tratamiento,
* laboratorio,
* personal altamente especializado,

la descentralización de la producción puede resultar absurdamente cara.

Pero aquí hay otra distinción conceptual que merece mucho la pena:

\boxed{\text{descentralización de la decisión}\neq\text{descentralización de la producción}}

Un municipio puede conservar capacidad para decidir qué quiere proporcionar y, sin embargo, producirlo conjuntamente con otros municipios.

Por ejemplo:

A+B+C+D
\longrightarrow
\text{consorcio supramunicipal}.

Esto permite simultáneamente:

\text{adaptación local}
+
\text{economías de escala}.

La OCDE menciona precisamente cooperación interjurisdiccional para capturar economías de escala en servicios como transporte, gestión de residuos, agua, bomberos u hospitales.  

Esta parte puede ser muy buena porque introduces soluciones intermedias.

No tienes:

\text{municipio}\quad \text{o}\quad \text{Estado}.

Puedes tener:

\text{mancomunidades},
\quad
\text{consorcios},
\quad
\text{áreas metropolitanas},
\quad
\text{contratación conjunta}.

⸻

9. Aquí es donde Alesina-Spolaore te permite unir dos dimensiones

Y creo que ésta podría ser una de las piezas más elegantes de tu exposición.

En un eje horizontal:

N=\text{tamaño de la jurisdicción}.

Trazas una curva decreciente:

C_{\text{escala}}(N)

y otra creciente:

C_{\text{heterogeneidad}}(N).

La suma:

SC(N)
=
C_{\text{escala}}(N)
+
C_{\text{heterogeneidad}}(N)

tiene forma de U.

Su mínimo determina:

N^*.

Gráficamente:

Coste
  │\
  │ \       Coste heterogeneidad
  │  \     /
  │   \   /
  │    \ /      Coste total
  │     U
  │      \
  │       \_____ Coste de escala
  └──────────────────── Tamaño jurisdicción
              N*

Conceptualmente es potentísimo:

las jurisdicciones pequeñas son buenas para representar preferencias; las grandes son buenas para repartir costes.

Y el federalismo fiscal consiste parcialmente en encontrar el tamaño territorial apropiado para cada función.

Alesina y Spolaore construyen precisamente su teoría del tamaño de las jurisdicciones alrededor de este trade-off.  

⸻














Como resultado de estos factores, el modelo de TIEBOUT predice que el \textbf{gasto local debería centrarse en programas de amplia base con pocas externalidades y economías de escala relativamente bajas}, como la reparación de carreteras, la recogida de basura y la limpieza de calles. De manera similar, las comunidades locales deberían desempeñar un papel más limitado en la provisión de bienes públicos que sean redistributivos (como la asistencia social en efectivo), tengan grandes efectos de desbordamiento (como la educación) y tengan economías de escala muy grandes (como la defensa nacional). \textcolor{red}{La naturaleza del federalismo fiscal en Estados Unidos es en gran medida coherente con esta predicción. Las obras públicas se financian principalmente a nivel local, los programas redistributivos se financian a nivel estatal y federal, y la defensa es un programa nacional. La educación se financia aproximadamente en una mitad por las localidades y en otra mitad por niveles superiores de gobierno (principalmente el gobierno estatal), lo cual es coherente con los efectos de desbordamiento asociados a la educación. La única cuestión aquí es si las externalidades de la educación son lo suficientemente grandes a escala nacional como para que el gobierno federal, que actualmente proporciona menos del 10\% del gasto educativo, deba desempeñar un papel mayor en la financiación de la educación, como ocurre en la mayoría de las demás naciones industrializadas.} 

\subsection{Perspectiva institucional: Economías de alcance en el ámbito administrativo}
10. Cuarta dimensión: capacidad administrativa, coordinación y economías de alcance

Esta es la que añadiría a tus tres.

Porque incluso en ausencia de grandes economías tecnológicas de escala, una jurisdicción pequeña puede sencillamente no disponer de capacidad institucional suficiente.

No es lo mismo administrar:

un parque municipal

que:

una red hospitalaria con tecnologías altamente especializadas.

Hay funciones que requieren:

\text{capital humano especializado},

\text{sistemas complejos de información},

\text{capacidad regulatoria},

\text{contratación sofisticada}.

Esto genera lo que podrías denominar economías administrativas de escala.

Y existen también economías de alcance:

C(A,B)<C(A)+C(B).

Por ejemplo, puede ser más eficiente que una misma administración coordine:

\text{urbanismo}
+
\text{transporte}
+
\text{vivienda}

porque las tres políticas están fuertemente interrelacionadas.

El Banco Mundial incluye precisamente economías de escala, economías de alcance, spillovers y proximidad a los beneficiarios entre los criterios relevantes para asignar funciones públicas.  


7. No existe una única «competencia»: hay que descomponer cada función

Esta probablemente sería mi principal aportación a tu esquema.

Decir:

«la sanidad corresponde a las regiones»

es económicamente muy impreciso.

Porque dentro de una misma función pública podemos distinguir:

\text{regulación}
\neq
\text{financiación}
\neq
\text{decisión}
\neq
\text{administración}
\neq
\text{producción}.

La literatura moderna sobre expenditure assignment insiste precisamente en que las responsabilidades compartidas deben descomponerse en funciones concretas. El Banco Mundial, por ejemplo, señala explícitamente que puede ser necesario separar regulación, financiación, administración y prestación del servicio entre distintos niveles de gobierno.  

Esto te abre ejemplos mucho más realistas.

En educación podrías tener:

\begin{aligned}
\text{currículo mínimo}&\rightarrow \text{Estado},\\
\text{financiación}&\rightarrow \text{Estado/región},\\
\text{gestión de centros}&\rightarrow \text{región/municipio},\\
\text{producción efectiva}&\rightarrow \text{centro educativo}.
\end{aligned}

En transporte:

\begin{aligned}
\text{normativa}&\rightarrow \text{Estado/región},\\
\text{planificación metropolitana}&\rightarrow \text{autoridad metropolitana},\\
\text{financiación}&\rightarrow \text{varios niveles},\\
\text{operación}&\rightarrow \text{empresa pública/privada}.
\end{aligned}

Eso es muchísimo más interesante que intentar decidir simplemente si «el transporte es local o regional».

⸻

















\subsection{Mutualización del riesgo: ¿Cómo diseñar un seguro óptimo?}
\textcolor{red}{\textbf{Origen:} ICEX-CECO}
IJ.II.JJ. Mutualizacián del riesgo

Otra importante línea de investigación de las aportaciones más recientes se refiere a la mutualización del riesgo. Se centran en el análisis de los problemas que pueden surgir al establecer un sistema de aseguramiento para las jurisdicciones y en aspectos que mejoran su funcionamiento. En este sentido, para establecer un aseguramiento interregional, es necesario que ciertas variables económicas relevantes tengan una correlación cercana a $0$ entre las distintas partes. Es decir, que cuando una región sufra una contingencia, otras no.

La mutualización del riesgo será también más efectiva cuanto mayor número de riesgos se contemplen. Además, será más ventajosa para las distintas partes, cuando exista una correlación negativa entre los riesgos que les afectan a cada una, ya que se refuerza así la capacidad de cobertura cruzada entre las regiones.

Es también importante que exista simetría entre las regiones de cara a evitar que siempre unas mismas regiones cubran las pérdidas de otras. Y, por último, es clave un comportamiento recíproco entre las regiones y que ese compromiso sea creíble, de fonna que una región que expe1imcnte un shock positivo, asegure a otras, esperando que estas le ayuden cuando se encuentren ante shocks negativos.

Aunque existe la posibilidad de que se dé aseguramiento voluntario. en la práctica en los sistemas descentralizados. el aseguramiento se realiza a través de impuestos y transferencias, como se analizará más adelante en la segunda parte del tema. 

De esta manera, con la mutualización de riesgos. el gobierno central puede ejercer su función de estabilización económica, ya que asegura a las jurisdicciones frente a shocks estocásticos que pueden reducir el bienestar social. Así. el riesgo de cada j urisdicción se reduciría si aumenta el número de partes que mutual izan el riesgo ("risk pooling'').

Supongamos que tenemos n jurisdicciones. cada una de las cuales con un nivel de renta yi, que es una variable aleatoria y sigue la misma distribución de probabilidad para todas las jurisdicciones y presentan correlación nula entre ellas. En este caso. E[yi] = yi y Var[yi] = cr2

La mutualiz.ación del riesgo en cambios en los niveles de renta para las jurisdicciones (es decir, si éstas compmten sus rentas y cada una recibe la n-ésima parte del total). la nueva variable aleatoria es (¿yi)/n , por lo que:

E[(¿yi)/n] = yi y Varl(Iyi)/n] = ( 1 /n2)*ncr2 = cr2/n

Es decir, la esperanza de renta se mantiene y la varianza disminuye conforme se incrementa el número de j urisdicciones.

Sin embargo, varios autores. como. por ejemplo. Bucovetsy en 1 997. o Perrson y Tabellini en 1 996. sei1alan que con el aseguramiento interregional ��urgen problemas de riesgo moral.

El problema es que, al existir este aseguramiento central. las jurisdicciones tienen incentivos a reducir la provisión de los bienes públicos locales y programas públicos que tengan como objetivo mejorar la capacidad de estas unidades económicas para adaptarse a shocks exógenos.

En este ��cntido. los autores Perrson y Tabellini se11alaban que el riesgo moral dependía de cómo estu,;era constituido el sistema descentralizado. En su modelo concluyen que, si el aseguramiento del gobierno central ante shocks se dirige directamente a los ciudadanos, en vez de dirigirse a los gobiernos locales, se reducen los problemas de riesgo moral. 

Pero el problema de riesgo moral que surge debe ponderarse con las distintas presiones que surgen para el sistema descentralizado en el caso de que no exista un sistem�� de aseguramiento para las  haciendas territoriales. Por ejemplo. en caso de recesión. las jurisdicciones reducirían su gasto, siguiendo un comportamiento procíclico, y aumentarían así los efectos negativos a nivel ahrrcgado. Hansen y Perloff en 1 944 denominaron a este problema la ''hipótesis de la corruptibilidad fiscar·.

En este contexto. ante estos shocks estocásticos, el gobierno central debería siempre asistir a las haciendas para que mantengan su gasto (por ejemplo. a través ele transferencias). Sin embargo, surge un problema de riesgo moral como el que se mencionó anteriormente.

Así, dado que en el caso de que esta asistencia o aseguramiento no sea pe1manente, suponiendo que se diera en caso de recesión. las jrnisdicciones tienen incentivos para crear un ··colchón·' que sirva para afrontar futuros shock\·. Aunque esto es importante. puede ser contraproducente si para crearlo se reduce la provisión de los bienes públicos locales. Además. otro problema del aseguramiento, es su limitación, dada la dificultad que conlleva definir las condiciones para las que éste se ofrecerá a las jurisdicciones ..

En definitiva, el establecimiento de un sistema de aseguramiento interregional en un sistema descentralizado está justificado para llevar a cabo la función estabilizadora, aunque es importante tener en mente que existen numerosos problemas asociados.

En conclusión. hasta ahora, se ha analizado la descentralización desde la óptica del gasto, justificada principalmente por motivos de eficiencia. Esta justificación debe completarse con las aportaciones de la segunda generación de la Teoría del Federalismo Fiscal ya que el sistema descentralizado puede mejorar la rendición de cuenws de los gobiernos, pero surge el problema del principal-agente.

Además. de cara al buen funcionamiento del sistema descentralizado, puede ser necesario incluir un sistema de aseguramiento que no está libre de problemas. Estos problemas dependen en cierta medida de la fomia que tome la estructura del sistema descentralizado. Para esto, debe completarse en análisis de la descentralización desde la óptica de los ingresos.











\clearpage
\section{Ingreso: ¿Cómo financiar los diferentes niveles de Gobierno?}
\puntosclave{
    
}

A continuación, se procede analizar la otra parte • ele los sistemas descentralizados, la descentralización por el lado de los ingresos.

Ill.I. Prim era generación 

Las aportaciones de la primera generación analizaban las distintas fuentes de financiación disponibles para las jurisdicciones. Concretamente, podemos hablar de tres fuentes: impuestos, endeudamiento, y transferencias. Estas fuentes pueden coexistir.

Las aportaciones ele la segunda generación de la Teoría del Federalismo Fiscal, como señalamos. intrnducen aspectos relevantes de la Teoría de la Elección Pública. Con esto, igual que en el caso de la descentralización del gasto, analizan importantes implicaciones que tiene el sistema descentralizado por el lado de los ingresos.



\subsection{Impuesto}

\textcolor{red}{\textbf{Origen:} ICEX-CECO}
JJI.J.l. Impuestos

Existen distintos tipos de tributos para los que las jurisdicciones tienen competencias. Dentro de los tributos, cabe destacar a los impuestos por su importancia cuantitativa, que son aquellos tributos exigidos sin contraprestación directa. Es decir, el sujeto se beneficia de la actividad de la administracíón de forma general.

Dentro ele los impuestos, los que mayor corresponsabilidad fiscal conllevan serán los impuestos propios en sentido estricto. ya que serán aquel los enteramente diseñados y gestionados (i.e. liquidación, recaudación, e inspección) por las haciendas territoriales.

E n segundo lugar, podemos hablar de los llamadps impuestos cedidos. Son aquellos creados por el Estado. pero en los que el diseño y la gestión se c.ede a las haciendas territoriales. Por otro lado. encontramos los impuestos compartidos, que son aquellos en los que el Estado diseña los elementos básicos del impuesto y los gestiona, pero deja cierta discrecionalidad de diseño a las haciendas territoriales.

Y. en último lugar. encontramos 1ambién los impuestos participados. que son aquellos diseñados y gestionados por el Estado. pero en los que éste transfiere parte de la recaudación a las haciendas tc1Tiloriales.


\subsubsection{Corresponsabilidad fiscal}

\textcolor{red}{\textbf{Origen:} ICEX-CECO}
IIJ .Jl. ll. Corresponsabilidad fiscal

La financil:lción de las haciendas territoriales a través de tributos tiene como objetivo principal la corresponsabilidad fiscal. Este término, como anticipamos. se refiere al hecho de que las jurisdicciones financien una parte sustantiva de su gasto mediante impuestos propios y v isibles para los contribuyentes.

La con-esponsabilidad fiscal tiene importantes implicaciones. En primer lugar. puede mejorar la descentralización del gasto. ya que potencialmente aumenta la responsabilidad de los gobierno�� locales en sus decisiones de gasto al tratar de ajustarlo sus ingresos, los cuales están c laramente definidos.
 
Además. los individuos que residan en dicha región perciben el coste verdadero de los bienes y servicios públicos locales que consumen y de los que se benefician, especialmente si se financian íntegramente con sus impuestos.

Por último, la corresponsabilidad fiscal rebaja las tensiones que pueden surgir entre el gobierno central y los locales debido a las transferencias necesarias. ya que con tributos propios las haciendas pueden disponer de herramientas suficientes para cubrir sus gastos.

En definitiva, aumentar la autonomía de las jurisdicciones en cuanto al diseño de los i.mpuestos pcnnite cumplir mejor con el principio de corresponsabilidad fiscal. No obstante, existen algunas desventajas. ya que surge competencia fiscal.



\subsubsection{Competencia fiscal}

\textcolor{red}{\textbf{Origen:} ICEX-CECO}
lll.JT. JJJ. Competencia fiscal

La competencia fiscal se puede definir como la interacción que surge entre distintos gobiernos debido a la posibilidad de movilidad de la base imposible entre distintas jurisdicciones. La pérdida de base imponible en una jurisdicción supone una ganancia para otras. de manera que esta movilidad supone una extemalidad entre las jurisdicciones.

De esta manera, las haciendas territoriales pueden acabar incurriendo en una carrera a la baja de impuestos para atraer recursos. El resultado final es ineficiente. ya que, si todas hacen lo mismo. la base imponible para cada hacienda será la misma pero la recaudación se reducirá por la bajada de tipos. Se da así un juego repetido con jugadores con estrategias gatillo.

Por esta razón. suelen descentralizarse aquellos impuestos que gravan bienes y servicios poco móviles. Musgrave en concreto afinnaba que los impuestos que se emplean para llevar a cabo la función rcdistributiva (por ejemplo, con el impuesto sobre la renta para personas físicas) y estabilizadora ( por ejemplo. con el impuesto sobre la renta para personas jurídicas), debían dejarse en manos de los gobiernos centrales, ya que debían ser funciones centrales. Consideraba que sólo debían descentralizarse aquellos impuestos que gravan bienes inmóviles. como. por ejemplo. la propiedad.

Además. en cuanto a otros tipos de tributos Musgrave sefialó, que las tasas y las contribuciones especiales debían descentralizarse. Éstos están relacionados con la aplicación del principio del  beneficio, y sí están ligados a actividades específicas de la administración, a diferencia de los impuestos.

Pero. además de existir competencia entre gobiernos locales, es decir. competencia fiscal horizontal. es importante resaltar que también existe competencia fiscal vertical. Surge cuando niveles altos y bajos de gobierno comparten una misma base imponible. 

En este caso, el impuesto establecido por un gobierno reducirá la base imponible disponible para gobiernos de otros niveles, dando lugar a extemalidades verticales negativas, y a un equilibrio donde existe sobreimposición.

Se puede reflejar esta idea con un ��iemplo simple. Supondremos un bien que inicialmente se grava tanto a nivel local como a nivel central. El bien tiene una curva de oferta completamente elástica. La demanda del consumidor tendrá pendiente negativa. Así. las variaciones en los impuestos se trasladan completamente al precio que percibe el consumidor.

Si cualquier gobierno, ya sea local o central, establece un aumento 6. al tipo impositivo, el precio final que percibe el consumidor aumentará también b.. Así, la cantidad consumida en el equilibrio se verá reducida. Por otro lado, los ingresos del gobierno que aumenta el gravamen aumentarán, ya que el efecto del aumento de ¡,r r:wamen contraITesta el de la reducción de la cantidad demandada.

Sin em bargo, el otro gobierno solo experimentará este último efecto. y verá sus ingresos reducidos. Para aumentarlos tiene la opción de aumentar su tipo impositivo. aunque podría comenzar una guerra de precios que dé lugar a tipos ineficientemente elevados.

En definitiva, la ventaja principal de los tributos como forma de financiación de las hacicnda:c. ten-itoriales se debe a una mayor co1Tesponsabilidad íiscal. No obstante. existe el problema de competencia fiscal, tanto a nivel horizontal corno vertical, que puede dar lugar a resultados ineficientes.

Sin embargo. algunos autores enmarcados en la segunda generación de la Teoría del Federalismo Fiscal resaltan algunas ventajas que puede ofrecer la competencia fiscal. Existen varios argumentos que justifican esto. Ya se comentó anteriormente. que, de acuerdo con las aportaciones de Brennan y Buchanan, la competencia fiscal puede reducir el poder impositivo de los gobiernos.

Otro argumento diferente se refiere a que la competencia fiscal actúa como mecanismo de compromiso. I--lindricks y Myles analizan esta idea mediante un juego donde participan los gobiernos locales y los propietarios de capital. En este juego, los gobiernos deciden qué tipos impositivos establecer, y los propietarios de capital dónde invertir. Las decisiones sobre los tipos son reversibles. pero las de inversión no.

Así, en un primer periodo los gobiernos anuncian los tipos impositivos de fonna independiente. Posteriormente, los inversores deciden Jonde invertirán. Como los gobiernos pueden modificar sus tipos una vez se han tomado las decisiones de inversión. aparece un problema de compromiso.

Si las decisiones de inversión son irreversibles, un gobierno local tendrá incentivos a gravar todos  los rendimientos de capital. Pero el propietario del capital anticipará esa situación. y decidirá no invertir en esa jurisdicción.

Con esto, si existe competencia fiscal. se producirá una relocalización del capital entre las distintas regiones. Así, la competencia fiscal actúa como un. mecanismo de compromiso, ya que se reducirán los incentivos perversos de los gobiernos locales. que no modificarán sus decisiones en aras de atraer o mantener el nivel de inversión actual.

Cabe mencionar, en último lugar. otro argumento que justifica una mejora en el bienestar debido a la competencia fiscal. Este caso es más fácil de aplicar a distintos países dentro de un sistema descentra lizado.

Hindricks y Myles analizan el caso en el que distintos países establecen subsidios para ayudar a las empresas de su país. Suponen que no existe la posibilidad de discriminar entre empresas nacionales y empresas extranjeras. De esta forma. sin movilidad de empresas. un país deseará establecer subsidios a las empresas de su territorio (nacionales) y tener ventaja comparativa respecto a otros países. Sin embargo, el resto de países responderá estableciendo también subsidios. Así. para evitar si tuaciones de ventaja, en equilibrio lodos los países terminarían estableciendo este subsidio y los efectos se cancelarían.

Si existe movilidad de empresas. cada país deberá otorgar estos subsidios tanto a las empresas domésticas como a las extranjeras. al no poder discriminar. No existirá entonces ventaja competitiva para las empresas nacionales. y aumentará el coste a asumir por el Sector Público.

Por lo tanto, si existe competencia fiscal. las empresas podrán establecerse en el país que deseen, y podrán obtener subsidios por parte de gobiernos extranjeros. De esta manera, la competencia fiscal puede reducir los incentivos para establecer este tipo de subsidios. Ninguno de los países establecerá subsidios y se evitará la situación de dilema del prisionero en la que todos los países den subsidios a las empresas. El bienestar social será mayor. Sin embargo, es importante resaltar que esto solo se cumplirá si no es posible discriminar entre entidades.

\subsubsection{Heterogeneidad recaudatoria: ¿debería preocuparnos?}
\textcolor{blue}{\textbf{OJO -> Origen:} GRUBER (2017)}
\textcolor{yellos}{\textbf{Redistribución entre comunidades}}
El modelo de TIEBOUT proporciona un marco para considerar uno de los problemas más importantes del federalismo fiscal: ¿Debería haber redistribución de fondos públicos entre comunidades? Actualmente existe una enorme desigualdad tanto en la capacidad de las comunidades locales para financiar bienes públicos (el valor de la base del impuesto sobre la propiedad) como en el grado en que lo hacen. Por ejemplo, en el estado de Massachusetts, el municipio de Lakeville recauda solo 13.932 dólares en ingresos fiscales locales por estudiante de escuela pública, mientras que el municipio de Carlisle recauda 22.472 dólares. Parte de esta diferencia procede de decisiones sobre el nivel de imposición local: el impuesto por cada 1.000 dólares de valor de la propiedad es de 14,25 dólares en Lakeville y de 17,68 dólares en Carlisle. Sin embargo, la mayor parte de la diferencia en los ingresos procede de diferencias subyacentes en los valores de las propiedades gravadas: la vivienda unifamiliar mediana vale 367.000 dólares en Lakeville y 851.722 dólares en Carlisle. En el estado de Nueva York, un estudio encontró que los valores de las propiedades por estudiante de escuela pública variaban por un factor de casi seis, teniendo el 10\% más pobre de los distritos valores de propiedad por estudiante inferiores a 313.891 dólares y el 10\% más rico de los distritos valores de propiedad por estudiante superiores a 1.871.956 dólares.

\textbf{¿Debería preocuparnos?}
¿Debería esta desigualdad en las bases de ingresos (tal como se refleja en los valores de las propiedades) o en los ingresos recaudados (el producto de los valores de las propiedades y los tipos del impuesto sobre la propiedad) entre comunidades preocupar a los responsables de las políticas públicas? ¿Deberían los niveles superiores de gobierno imponer la redistribución entre los niveles inferiores de gobierno para compensar estas diferencias? Como se señaló anteriormente, dicha redistribución es una característica importante del federalismo fiscal en algunas naciones, donde el gobierno nacional distribuye subvenciones a las comunidades más pobres que compensan en gran medida las diferencias de ingresos entre comunidades.

La respuesta general a la pregunta ¿Debería preocuparnos? es que depende del grado en que el modelo de TIEBOUT describa la realidad. En un mundo de TIEBOUT perfecto, no redistribuiríamos entre comunidades: las comunidades se habrían formado para la provisión eficiente de bienes públicos, y cualquier redistribución entre ellas obstaculizaría la eficiencia. Si un municipio tiene ingresos bajos o un gasto bajo, es porque los residentes del municipio han elegido proporcionar un nivel bajo de bienes públicos, y este es el resultado eficiente dadas sus preferencias. La redistribución gubernamental en este caso debería centrarse en los individuos, no en las comunidades.

Sin embargo, en la medida en que TIEBOUT no describa perfectamente la realidad, existen dos argumentos para redistribuir desde las comunidades con altos ingresos y alto gasto hacia las comunidades con bajos ingresos y bajo gasto. El primero son los fallos del mecanismo de TIEBOUT. Por ejemplo, supongamos que existen razones por las que las personas no pueden votar eficazmente con los pies, como unas normas de zonificación restrictivas que hacen que las viviendas sean muy grandes y caras en las comunidades con un nivel elevado de bienes públicos (por ejemplo, cada vivienda debe estar situada en una parcela de al menos un acre). En esta situación, puede haber personas que deseen niveles elevados de bienes públicos pero que no puedan permitirse la elevada calidad de vivienda exigida por las normas de zonificación. Estas personas podrían permanecer atrapadas en un municipio con una baja provisión de bienes públicos, el único lugar donde pueden permitirse una vivienda. En este caso, podría ser eficiente redistribuir hacia los municipios con bajos niveles de bienes públicos para ayudar a los individuos atrapados en una situación en la que se ven obligados a consumir una cantidad insuficiente de bienes públicos.

La segunda razón para la redistribución son las externalidades. Si una gran proporción de los ingresos fiscales locales se gasta en bienes públicos locales con efectos de desbordamiento o externalidades para otras comunidades, existe el argumento estándar de las externalidades para que los niveles superiores de gobierno subvencionen el gasto en las comunidades que proporcionan las externalidades. Por ejemplo, supongamos que una educación primaria de alta calidad en un municipio conduce a tasas de delincuencia más bajas tanto en ese municipio como en los municipios vecinos. En este caso, puede ser óptimo que el gobierno estatal grave a los municipios con ingresos elevados y redistribuya hacia los municipios con ingresos bajos para garantizar que los municipios con ingresos bajos puedan proporcionar una educación primaria de alta calidad.




\subsection{Transferencias}

\textcolor{blue}{\textbf{OJO -> Origen:} GRUBER (2017)}
\textcolor{red}{\textbf{Instrumentos de redistribución: subvenciones}}
Si los niveles superiores de gobierno deciden, por una de las dos razones señaladas, redistribuir entre los niveles inferiores de gobierno, lo hacen mediante subvenciones intergubernamentales, que son transferencias de efectivo de un nivel de gobierno a otro. Los niveles superiores de gobierno utilizan varios tipos diferentes de subvenciones. Al definir estos tipos, utilizamos el ejemplo de un estado que redistribuye hacia las comunidades locales (aunque la misma descripción se aplica a otras formas de redistribución de un nivel superior a uno inferior de gobierno, como del nacional al estatal).

Supongamos que Barcelona proporciona solo un bien público a sus residentes: educación. Financia la educación mediante impuestos locales, y cualquier dinero que les quede a las familias después de los impuestos se gasta en bienes privados (como automóviles o ropa). La Figura 10-2 muestra la situación en Lexington antes de que se proporcione ninguna subvención. Los residentes de Barcelona tienen un presupuesto total de 1 millón de dólares para gastar en educación y otros bienes privados, y modelizamos cómo eligen dividir este presupuesto. En el punto A, los residentes de Lexington eligen no gastar nada en educación y gastar todo su presupuesto de 1 millón de dólares en bienes privados. En el punto B, Lexington gasta todo su presupuesto de 1 millón de dólares en educación y nada en bienes privados.

Los votantes de Lexington tienen ciertas preferencias por la educación y los bienes privados que pueden representarse mediante una curva de indiferencia IC1 entre estos dos conjuntos de bienes. Es decir, podemos analizar la elección de Lexington entre educación y bienes privados de la misma manera que podríamos analizar la elección de un individuo entre estos mismos elementos; IC1 representa la agregación de las curvas de indiferencia de los votantes mediante un mecanismo de votación. Antes de que exista ninguna subvención estatal, Lexington elige gastar 500.000 dólares al año en educación y 500.000 dólares al año en bienes privados. Esta combinación de gasto está representada por el punto X, donde la curva de indiferencia del municipio es tangente a su restricción presupuestaria.




\textcolor{red}{\textbf{Origen:} ICEX-CECO}
lll.l.Ill Tram.ferencias

Una de las fuentes de financiación que más se dan en la práctica en los sistemas descentralizados son las transferencias, la transmisión unilateral de recursos entre diferentes niveles de gobierno.

Sin embargo. este método de financiación no está libre de limitaciones. Por ejemplo. como se anticipó. pueden crear tensiones entre gobiernos locales y el gobierno central. y, además, no contribuyen a la corresponsabilidad fiscal. Por esto último, al emplear transferencias se puede obtener un gasto público excesivo financiado por ciudadanos ajenos a dicho territorio.


\textcolor{red}{\textbf{OJO -> Origen:} ROSEN y GARBER}
\textcolor{yellow}{\textbf{Subvenciones intergubernamentales}}

¿Por qué han crecido tanto las transferencias intergubernamentales a largo plazo? Esta cuestión está estrechamente relacionada con la razón por la que ha aumentado el gasto público en general. Una explicación del crecimiento de las subvenciones pone de relieve que, durante las últimas décadas, la demanda de los tipos de servicios proporcionados tradicionalmente por el sector estatal y local —educación, transporte y protección policial— ha estado creciendo rápidamente. Sin embargo, las estructuras de ingresos estatales y locales, que se basan principalmente en impuestos sobre las ventas y sobre la propiedad, no han proporcionado los medios para mantener el ritmo del crecimiento de los gastos deseados. En cambio, los ingresos fiscales federales han crecido automáticamente con el tiempo, en gran medida debido a la naturaleza progresiva del impuesto federal sobre la renta de las personas físicas y, hasta la llegada de la indexación a mediados de la década de 1980, a la inflación. Por tanto, existe un desajuste entre dónde se recauda el dinero de los impuestos y dónde se demanda. Las subvenciones del gobierno central a los estados y localidades proporcionan una forma de corregir este desajuste.

La teoría del desajuste no es satisfactoria porque no logra explicar por qué los estados y las localidades no pueden aumentar sus tipos impositivos para mantenerse al ritmo de los aumentos de la demanda de bienes y servicios públicos locales. Como se señala en la siguiente sección, probablemente tengamos que recurrir a consideraciones políticas para explicar el patrón de las subvenciones intergubernamentales.

\subsubsection{Tipología: ¿qué esquemas de transferencias existen?}



\paragraph{Finalidad: }

\textcolor{blue}{\textbf{OJO -> Origen:} GRUBER (2017)}
\textbf{Subvenciones de contrapartida }
Un tipo de subvención que podría utilizar el gobierno estatal es una subvención de contrapartida, que vincula la cantidad de fondos transferidos a la comunidad local con la cantidad de gasto que esta asigna actualmente a bienes públicos. Por ejemplo, una subvención de contrapartida de uno por uno para educación proporcionaría 1 dólar de financiación del estado por cada 1 dólar de gasto en educación por parte de la comunidad local. Aunque aquí utilizamos como ejemplo una contrapartida de uno por uno, las tasas de contrapartida pueden variar desde 0,01 hasta más de 1.

Esta subvención de contrapartida de uno por uno reduce a la mitad el precio de la educación; cada dólar de gasto en educación cuesta ahora a Lexington solo 0,50 dólares porque el estado de Massachusetts proporciona los otros 0,50 dólares. Este cambio hace pivotar hacia fuera la restricción presupuestaria de AB a AC en la Figura 10-3. Esta subvención aumenta inequívocamente el gasto en educación mediante tanto el efecto renta como el efecto sustitución. En nuestro ejemplo, el gasto total en educación aumenta de 500.000 a 750.000 dólares en el punto Y. Lexington aporta 375.000 dólares a la educación y recibe los otros 375.000 dólares en subvenciones de contrapartida. De su presupuesto original de 1 millón de dólares, Lexington tiene ahora 625.000 dólares para gastar en bienes privados (los 500.000 dólares originales que estaba gastando más los 125.000 dólares que ya no gasta en educación). Como resultado de la subvención de contrapartida, por tanto, ha aumentado el gasto total tanto en educación como en bienes privados.

\textbf{Subvención en bloque }
Otra opción de subvención es una subvención en bloque, mediante la cual el estado simplemente entrega a la comunidad local una determinada cantidad de subvención sin ninguna obligación sobre cómo debe gastarse. Para mantener constante el coste para el gobierno, supongamos que el gobierno concede a Lexington una subvención en bloque de 375.000 dólares. Debido a que la subvención en bloque hace que Lexington sea lo suficientemente rico como para poder permitirse gastar hasta 1,375 millones de dólares tanto en educación como en bienes privados, desplaza hacia fuera la restricción presupuestaria de AB a DE, como ilustra la Figura 10-4.

Aunque Massachusetts está dando a Lexington la misma cantidad de dinero con la subvención en bloque, esta tiene un efecto muy diferente sobre el comportamiento del municipio. Parte de esta nueva riqueza se utilizará para aumentar el gasto en educación, mientras que otra parte se utilizará para aumentar el consumo de bienes privados. En este ejemplo, el municipio se desplaza al punto Z, aumentando el gasto en educación en solo 75.000 dólares y el gasto en bienes privados en 300.000 dólares (de 500.000 a 800.000 dólares).

El aumento del gasto en educación es menor con la subvención en bloque (75.000 dólares) de lo que fue con la subvención de contrapartida (250.000 dólares) porque ahora solo existe un efecto renta sobre el gasto en educación para Lexington, mientras que la subvención de contrapartida tenía tanto un efecto sustitución como un efecto renta. El efecto renta eleva el gasto en educación de 500.000 a 575.000 dólares, desplazando al municipio del punto X al punto Z. El efecto sustitución que se añade con la subvención de contrapartida eleva entonces el gasto en educación en 175.000 dólares adicionales, hasta 750.000 dólares, como refleja el desplazamiento del punto Z al punto Y.

Por otro lado, Lexington ha quedado en una situación mejor con la subvención en bloque que con la subvención de contrapartida. Esto puede observarse gráficamente por el hecho de que, con la nueva restricción presupuestaria bajo la subvención en bloque (DE), el municipio podría haberse permitido su elección en el punto Y, con el gasto en educación aumentando hasta 750.000 dólares y el gasto en bienes privados aumentando hasta 625.000 dólares, pero eligió una combinación diferente. Dado que el municipio eligió en su lugar el punto Z, debe encontrarse en una curva de indiferencia más alta. Es decir, dada la libertad de gastar el dinero de su subvención como desee, sin la restricción de una condición de contrapartida, el municipio preferiría gastar la mayor parte del dinero en bienes privados y relativamente poco en educación. La subvención de contrapartida conduce a un mayor gasto en educación del que el municipio elegiría de otro modo dada esa cantidad de dinero, por lo que deja al municipio en una curva de indiferencia más baja.


\paragraph{Condicionalidad: ¿cómo escoger la forma más adecuada?}
Se pueden dividir a las transferencias inteegubernamentales en dos grandes categorías: las \textbf{transferencias incondicionales} - que no tienen condiciones sobre su destino o empleo; o las \textbf{condicionales} - qué sí deben dedicarse a un uso particular.

Para analizar el efecto que tienen estas transferencias en el bienestar, vamos a suponer una jurisdicción que provee un único bien público llamado G, no sujeto a congestión. El coste unitario de producción es constante. e igual a P<,, y se costea entre el número total de residentes. N.

En cuanto a los residentes de esta hacienda, tienen una renta agregada total Yo. El consumo ele bien privado será igual a la renta disponible, llamada YD, que queda tras el pago de impuestos necesarios para financiar el gasto público. Se supone, además, que tanto el bien privado como el bien público son bienes normales.

Por otro lado, la utilidad del agente representativo dependerá de la renta per cápita disponible de la _jurisdicción, y del nivel de provisión del bien público. De esta forma, tendremos el siguiente problema de optimización:

max U (�� , G)
s. a. YD + PcG = Y0

Resolviendo este problema se obtiene el nivel óptimo de provisión, Go. Vendrá dado por el punto de tangencia de la restricción presupuestaria y la curva de indiferencia más alejada del origen.

\textbf{INCONDICIONAL}
Si a continuaci·ón suponemos que se establece una transferencia incondicional T del gobierno central a la jurisdicción, la restricción presupuestaria pasará a ser:

Gráficamente, se produce un desplazamiento paralelo de la restrieción presupuestaria. en una magnin1d igual a T. Únicamente existe. por tanto, un efecto renta. Dado que los bienes considerados son nomrnles, tendremos que tanto el consumo público G como el privado Yn. aumentan.

\textbf{GRÁFICO 6.- TRANSFERENCIA NACIONAL}

El consumo privado aumenta con la transferencia incondicional. ya que con los recursos adicionales la jurisdicción podrá aumentar las transferencias a los hogares o reducir los impuestos en esa cuantía. Sin embargo, al haber supuesto bienes non11ales, el aumento del gasto público s��rá inferior a la tran sferencia.

Se obtienen resultados diferentes si se considera que la transferencia es condicional. En este caso, además. existen dos tipos de tran sferencias : las condicionales compensatorias, y las no compen satorias.

Las primeras son aquellas que se conceden condicionadas a un uso particular. pero su cuantía es independiente del nivel de gasto que la jurisdicción dedique a ese concepto. Por su parte, las compensatorias sí tienen en cuenta ese nivel de gast o previo por parte de la hacienda territorial.


\textbf{CONDICIONAL}
Si en el ejemplo anterior considerarnos una transferencia condicional no compensatoria T, la restricción se desplaza de nuevo a la derecha, existiendo únicamente un efecto renta, como antes. Pero ahora, los puntos en los que el nuevo gasto público sea inferior a la transferencia no serán alcanzables.

\textbf{GRÁFICO 7.- TRANSFERENCIA CONDICIONAL NO COMPANESATORIA}

Es decir. no son alcanzables aquellos puntos en los que el nuevo gasto público sea inferior a la transferencia. ya que ésta se concede para dedicarse en su integridad a dicho fin.

Si las preferencias de los agentes son muy fuertes por el consumo privado en relación al consumo público, el agente representativo desearia situarse en un punto como el 2. Esto no es posible ya que la condicionalidad opera, y el equilibrio será el punto 2', en una curva de indiferencia con un nivel de bienestar infe1ior, tal y como se observa en el Gráfico 6.

Pero en caso de que las preferencias entre el bien público y el privado sean más similares, la condicionalidad podría no ser operativa, obteniendo, por ��jemplo, un equilibrio en el punto l . En ambos casos tendremos que el gasto público aumenta, aunque en una cuantía inferior a l a de la transferencia, ya que el gobierno local reducirá su aportación a la financiación del bien público. En caso de la transferencia condicional compensatoria se obtienen resultados distintos. En este caso, el gobierno central establece una transferencia igual a t unidades monetarias. por cada unidad monetaria prevista en el gasto para la provisión del bien público por parle de la jurisdicción.

La restricción presupuestaria es entonces:
Y0 + (l + t) G = Y0

La diferencia principal con las transferencias vistas hasta ahora, es que aquí se modifican los incentivos de la jurisdicción de invertir en ese gasto objetivo. al hacerse comparativamente menos costoso que el resto de gastos.

Ahora no solo e��iste un efecto riqueza, sino también un efecto sustitución que refuerza la expansión del gasto. Así, se modifican los precios relativos y gráficamente pivota la restricción presupuestaria.

\textbf{GRÁFICO 8.- TRANSFERENCIA CONDICIONAL COMPENSATORIA}

El equilibrio se producirá en un punto como el punto 1. según la elasticidad-precio de la demanda del bien público.

Tendremos ahora que el consumo del bien público será mayor que en los casos anteriores, ya que, en los casos anteriores, las haciendas territoriales perciben la transferencia como un ingreso extraordinario y lo gastan como tal: dedicando parte del gasto al bien público objetivo. y el resto a consumo privado.

No obstante, en las transferencias condicionales compensatotias surge un efecto distorsionante. ya que con una transferencia incondicional se podría obtener el mismo nivel de bienestar final. y a menor coste.

En definitiva, a la hora de elegir entre los distintos tipos de transferencia, el gobierno central debe tener en cuenta cuál es el objetivo princ ipal. Si desea estimular un gasto público concreto. lo más adecuado sería establecer una transferencia condicional y compensatoria. Sin embargo, si desea únicamente aumentar el bienestar social de las haciendas, las transferencias incondicionales serán preferibles.

\begin{specialblock}{\textbf{Extensión:} Efecto \textit{Flypaper}}
    Sin embargo. existen estudios em píricos que muestran que la propensión marginal de dedicar la transferencia intergubernamental a gasto público es siempre mayor que la propensión marginal de dedicar otros ingresos a gasto público. Hines y Thaler en 1995 llevaron a cabo un análisis intenso de esta literatura empírica. Los estudios incluidos concluían que las transferencias intergubernamentales suponían un incremento del gasto en el progra ma objetivo de un 25% del total de la transferencia. Este resultado es mayor a la propensión del gobierno que recibe la transferencia a gastar en un programa públicos sus ingresos regulares (sin transferencias). Hines y l11aler·estimaron que solo entre el 5 y el 1 0°/4, de los nuevos ingresos que no provinieran de transferencias se destinaría a gasto para programas públicos.

    Arthur Okun denominó a este fenómeno ''.flypaper e/Ject'". o efecto pegamento. Así. el gobierno central deberá tener en cuenta este fenómeno a la hora de diseñar sus objetivos principales en relación al gasto público.
    
    Existen muchas ap011aciones empíricas que analizan la relación entre las transferencias y el gasto público local.
    
    K.night en 2002 mostró que las jurisdicciones que mayor preferencia tenían por los bienes públicos. ejercían una mayor presión ante el gobierno central para obtener transferencias. Con esto. existiría una correlación positiva entre transferencias y el gasto público, pero no debido al efecto pegamento.
    
    Otras aportaciones que obtienen evidencia contra el efecto pegamento son las de Gordon de 2004. que muestra que los gobiernos locales que reciben transferencias ajustan otrns fuentes de financiación como respuesta: Baicker y Staiger en 2005 muestran qúe. a mayor capacidad institucional de las jurisdicciones, menor es el efecto pegamento; o Strumpt: que en 1 998 mostró que el efecto pegamento será mayor cuanto menor sea el poder de los burócratas de los gobiernos locales.
\end{specialblock}

\textcolor{blue}{\textbf{Origen:} GRUBER}
\textbf{Subvención en bloque condicionada}
Supongamos que a Massachusetts le gusta el hecho de haber mejorado la situación de Lexington con una subvención en bloque más que con una subvención de contrapartida, pero no le gusta el hecho de que el gasto en educación no haya aumentado tanto. Una forma en que el estado podría tratar de remediar esto es mediante una subvención en bloque condicionada, una cantidad fija de dinero distribuida al municipio con el mandato de que el dinero se gaste únicamente en educación. En este caso, el estado podría proporcionar a Lexington una subvención en bloque de 375.000 dólares y exigir que gastara la totalidad de la subvención en educación.

El efecto de esta subvención en bloque condicionada se ilustra en la Figura 10-5. Lexington puede ahora gastar hasta 375.000 dólares (la cantidad de la subvención) en educación, mientras continúa gastando su presupuesto original de 1 millón de dólares en bienes privados. Así, el primer segmento de la restricción presupuestaria es ahora AF. Sin embargo, una vez que Lexington gasta más de 375.000 dólares en educación, se enfrenta a la misma disyuntiva entre gastar en educación y gastar en bienes privados a la que se enfrentaba cuando recibió la subvención incondicionada: la condición impuesta a esta subvención no importa si el municipio ya está gastando más de 375.000 dólares en educación. La nueva restricción presupuestaria es, por tanto, AFE. Más allá del punto de 375.000 dólares en el eje horizontal, esta nueva restricción presupuestaria es la misma que la restricción presupuestaria de la subvención en bloque incondicionada.

Como queda claro en la Figura 10-5, añadir esta condición no tiene ningún efecto sobre el comportamiento de Lexington: el municipio sigue eligiendo gastar los mismos 575.000 dólares en educación que gastaba con la subvención en bloque incondicionada (en el punto Z). Dado que Lexington ya estaba gastando más de 375.000 dólares en educación, esta subvención no está, en la práctica, condicionada para el municipio: tiene el mismo efecto que si el estado simplemente le hubiera dado 375.000 dólares para gastar en cualquier cosa.

Así, el municipio ha deshecho el mandato de gastar el dinero en educación reasignando el gasto existente para cumplir el mandato. Este es un ejemplo del tipo de desplazamiento (crowd-out) que analizamos en el Capítulo 7. El gobierno dio al municipio 375.000 dólares para gastar en educación, pero el municipio gastó solo 75.000 dólares netos de ese dinero en educación; gastó los 300.000 dólares restantes en bienes privados. Por tanto, el 80 % (300.000 dólares/375.000 dólares) del gasto del estado fue desplazado por la reacción del municipio. A pesar de una gran subvención estatal, el gasto local en educación aumentó solo en una pequeña cantidad.

El efecto de una subvención en bloque condicionada diferirá del de una subvención en bloque incondicionada únicamente si el municipio que recibe la subvención hubiera gastado menos que la cantidad de la subvención sin que se impusiera la condición. Es decir, añadir la condición a la subvención en bloque afectaría al comportamiento de Lexington únicamente si este hubiera elegido gastar menos de 375.000 dólares en educación con la subvención en bloque incondicionada. En ese caso, hacer que la subvención en bloque fuera condicionada aumentaría el gasto educativo de Lexington en más de solo 75.000 dólares. Si municipios como Lexington gastaran más de 375.000 dólares en educación independientemente de esta restricción, entonces imponer la restricción no tendría ningún efecto.



\subsubsection{¿Cuál es el esquema óptimo de transferencias?}
Por tanto, la elección óptima de los mecanismos de subvención para los niveles superiores de gobierno (como los estados) depende del objetivo del programa de subvenciones. 

Si el objetivo es maximizar el bienestar del nivel inferior de gobierno, las subvenciones en bloque serán las más eficaces. Si el objetivo es fomentar el gasto en bienes públicos como la educación, las subvenciones de contrapartida serán las más eficaces porque pondrán en funcionamiento tanto el efecto renta como el efecto sustitución para aumentar el gasto del municipio.



\textcolor{blue}{\textbf{Origen:} GRUBER}
\tetxcolor{red}{La redistribución en acción: igualación de la financiación escolar}

Los estados pueden intentar compensar estas desigualdades utilizando los tipos de subvenciones que acabamos de analizar. Al recaudar ingresos fiscales de todas las comunidades y después redistribuir los ingresos mediante subvenciones en bloque o de contrapartida a determinadas comunidades con bajos valores de las propiedades o bajo gasto educativo, el estado puede intentar igualar el gasto educativo entre distritos. Desde 1970, todos los estados han realizado al menos un intento de igualación de la financiación escolar, algunos impulsados por los tribunales estatales y otros por el público votante.

La estructura de los sistemas de igualación 
Los sistemas de igualación pueden adoptar formas muy diferentes. Algunos estados tienen sistemas que intentan igualar completa o casi completamente el gasto entre distritos escolares. California, por ejemplo, proporciona un nivel básico de financiación educativa a sus distritos escolares y prohíbe diferencias entre distritos escolares superiores a 350 dólares en el gasto por alumno. Una vez que un distrito está gastando 350 dólares más que el distrito con el gasto más bajo, todos los impuestos adicionales sobre la propiedad recaudados por el municipio se entregan al estado para su distribución a otros distritos. Así, bajo este sistema, un municipio no recibe ningún beneficio por aumentar sus propios impuestos locales sobre la propiedad porque los ingresos adicionales se dividen entre los distritos de todo el estado.

Menos extremos son los estados que han establecido un impuesto sobre la propiedad a escala estatal que se redistribuye de una manera que garantiza un determinado nivel básico de financiación por alumno para cada municipio. Por ejemplo, en el estado de Nueva Jersey, los municipios con valores de las propiedades por encima del percentil 85 de los valores de las propiedades del estado simplemente reciben una pequeña subvención básica del estado y tienen que recaudar localmente los demás ingresos destinados a la educación. Los municipios con valores de las propiedades por debajo del percentil 85 de los valores de las propiedades del estado reciben una subvención de contrapartida que es un múltiplo de su propio gasto educativo, lo que proporciona así a los municipios un incentivo para aumentar su gasto.²⁵

Los efectos de la igualación 
Diversos estudios económicos han evaluado los efectos de la igualación de la financiación escolar. Estos estudios generalmente coinciden en que las leyes de igualación han tenido el efecto pretendido de igualar el gasto escolar entre comunidades, y la igualación del gasto parece haber conducido también a una igualación de los resultados de los estudiantes. Murray, Evans y Schwab (1998), por ejemplo, concluyeron que las igualaciones ordenadas por los tribunales redujeron la desigualdad del gasto dentro de los estados entre un 19 y un 34 %. Card y Payne (2002) encontraron que las igualaciones redujeron en 8 puntos la brecha en las puntuaciones medias del SAT entre los hijos de padres con un alto nivel educativo y los hijos de padres con un bajo nivel educativo, o aproximadamente un 5 % de la brecha.

Existe menos acuerdo acerca de si esta igualación se ha producido aumentando el gasto entre los distritos con bajo gasto, reduciendo el gasto entre los distritos con alto gasto, o ambas cosas. Hoxby (2001) proporciona un estudio detallado de esta cuestión, en el que calculó el precio fiscal de los sistemas de igualación escolar, la cantidad de ingresos que un distrito local tendría que recaudar para obtener 1 dólar más de gasto. Los distritos de California se enfrentan a un precio fiscal infinito: independientemente de cuántos ingresos recauden mediante impuestos locales, no pueden elevar su gasto educativo local a más de 350 dólares por alumno por encima del distrito con el gasto más bajo. Los distritos de Nueva Jersey tienen en su mayoría precios fiscales inferiores a 1: un distrito podría recaudar 0,60 dólares de sus propios ingresos para recibir 0,40 dólares de ayuda estatal, obteniendo un total de 1 dólar de aumento del gasto. Este distrito tendría, por tanto, un precio fiscal de 0,6.

Hoxby encontró que los sistemas de igualación extremos con precios fiscales muy elevados, como el de California, conducen a una reducción general del gasto por alumno. Dado que cualquier impuesto que los municipios recauden por encima del nivel mínimo (más 350 dólares) simplemente es tomado por el estado y redistribuido a otros distritos, existe un incentivo para reducir los impuestos y disminuir el gasto. La igualación de California provocó una caída del 15 % en el gasto por alumno; el gasto de Nuevo México cayó un 13 %; y el gasto de Oklahoma, Utah y Arizona cayó un 10 % en cada caso. Estados como California igualaron el gasto por alumno, pero únicamente «igualando hacia abajo», es decir, reduciendo el gasto educativo general en todos los distritos. El resultado ha sido un deterioro general de la calidad de las escuelas públicas y una huida hacia las escuelas privadas por parte de los estudiantes que pueden permitírselo. Los sistemas de igualación con precios fiscales bajos, como los de Nueva Jersey, Nueva York y Pensilvania, en realidad aumentaron el gasto por alumno entre un 7 y un 8 %; estos estados, por tanto, lograron «igualar hacia arriba». Así, la igualación de la financiación escolar puede lograr sus efectos pretendidos de mejorar el gasto educativo de los distritos con poca riqueza únicamente si el sistema está diseñado de una manera que proporcione a esos distritos incentivos para aumentar su gasto sin penalizar excesivamente a los distritos con mayor riqueza.

Ramónson et al. (2014) confirmaron los hallazgos de Hoxby sobre los diferentes efectos de las reformas de la financiación escolar y posteriormente examinaron las consecuencias en etapas posteriores de la vida para los niños sujetos a estas reformas. Encontraron que había pocos efectos sobre las familias no pobres, pero grandes efectos sobre las más pobres: un aumento del 20 % en el gasto condujo a casi un año completo más de escolarización y a unos ingresos un 20 % mayores, suficiente para eliminar la mayor parte de la brecha entre los niños criados en familias pobres y los criados en familias no pobres.



\textbf{Tipos de subvenciones}

La estructura de una subvención influye en su impacto económico. Existen básicamente dos tipos, condicionadas e incondicionadas, que analizamos por separado.

Subvenciones condicionadas Estas se denominan a veces subvenciones categóricas. El donante especifica, hasta cierto punto, los fines para los cuales el receptor puede utilizar los fondos. La gran mayoría de las subvenciones federales están destinadas a fines específicos, y las reglas para gastar el dinero a menudo se detallan minuciosamente. Por ejemplo, el gobierno federal concede subvenciones a los estados para establecer programas contra la conducción bajo los efectos del alcohol. Los términos de la ley especifican todo, desde el porcentaje de concentración de alcohol en sangre que constituye intoxicación hasta cuánto tiempo después de que un infractor sea condenado debe retirársele el permiso de conducir. Tales restricciones no son atípicas.

\textcolor{red}{Existen varios tipos de subvenciones condicionadas.}

\textcolor{red}{\textbf{OJO -> Origen:} ROSEN y GARBER}
\textbf{Subvenciones de contrapartida }
Por cada dólar entregado por el donante para apoyar una actividad determinada, el receptor debe gastar una determinada suma. Por ejemplo, una subvención podría estipular que, siempre que una comunidad gaste un dólar en educación, el gobierno federal aportará también un dólar.

La teoría estándar de la elección racional puede ayudarnos a analizar las subvenciones de contrapartida. En la Figura 22.4, el eje horizontal mide la cantidad de producción del gobierno local, G, consumida por los residentes del municipio de Smallville. El eje vertical mide el consumo total de Smallville, c. Supongamos, para simplificar, que las unidades de G y c se definen de modo que el precio de una unidad de cada una sea 1 dólar. Por tanto, suponiendo que no hay ahorro, c es igual a la renta después de impuestos. Con estos supuestos, la restricción presupuestaria de Smallville entre c y G es una línea recta (AB en la Figura 22.4) cuya pendiente en valor absoluto es uno.¹⁵ La pendiente unitaria indica que por cada dólar que Smallville está dispuesto a gastar puede obtener una unidad de bien público.

Supongamos que las preferencias de Smallville por G y c pueden representarse mediante un conjunto de curvas de indiferencia de forma convencional.¹⁶ Entonces, si el municipio trata de maximizar su utilidad sujeto a la restricción presupuestaria, elige el punto E₁, donde el consumo de bienes públicos es G₁ y la renta de la comunidad después de impuestos es c₁.

Supongamos ahora que se establece un régimen de subvención de contrapartida de uno por uno. Cuando Smallville renuncia a 1 dólar de renta, puede obtener 2 dólares de G: uno de sus propios dólares y otro del gobierno federal. La pendiente (en valor absoluto) de la línea presupuestaria de Smallville pasa, por tanto, a ser un medio. En efecto, la subvención de contrapartida reduce a la mitad el precio de G. Es una subvención ad valorem al consumo del bien público. La nueva línea presupuestaria aparece representada en la Figura 22.4 como AR.

Smallville consume ahora G₂ bienes públicos y tiene c₂ disponible para consumo privado. Obsérvese que no solo G₂ es mayor que G₁, sino que c₂ también es mayor que c₁. Smallville utiliza parte de la subvención para comprar una mayor cantidad del bien público y parte para reducir su carga fiscal. Sería posible, por supuesto, dibujar las curvas de indiferencia de modo que c₂ fuera igual a c₁, o incluso de modo que c₂ fuera menor que c₁. No obstante, existe una clara posibilidad de que parte de la subvención destinada a estimular el consumo público se utilice no para comprar más G, sino para obtener una reducción de impuestos. En un caso extremo, las curvas de indiferencia de la comunidad podrían ser tales que G₂ = G₁: la comunidad consume la misma cantidad del bien público y utiliza toda la subvención para reducir los impuestos. Así, la teoría por sí sola no puede indicar cómo afecta una subvención de contrapartida al gasto de una comunidad en un bien público. Depende de la capacidad de respuesta de la demanda a los cambios en el precio.

Por tanto, los economistas han realizado estudios estadísticos sobre cómo varían con sus precios las demandas de diversos bienes públicos. Según la literatura examinada por Fisher y Papke [2000], la elasticidad-precio de la demanda de educación se sitúa entre 0,15 y 0,50.

Una subvención de contrapartida es una forma sensata de corregir la presencia de una externalidad positiva. Como se explicó en el Capítulo 5, cuando un individuo o una empresa genera una externalidad positiva en el margen, una subvención apropiada puede aumentar la eficiencia. La misma lógica se aplica a una comunidad. Por supuesto, todos los problemas que surgen al aplicar el sistema de subvenciones siguen estando presentes. En particular, el gobierno central tiene que ser capaz de medir el tamaño real de la externalidad. En este contexto, resulta interesante señalar que muchos programas federales de subvenciones son muy difíciles de racionalizar utilizando criterios de eficiencia. Las elevadas tasas de contrapartida (a menudo del 80 al 90 %) son mucho mayores que las estimaciones razonables de las externalidades generadas por las actividades estatales y locales subvencionadas [Oates, 1999, p. 1129]. De hecho, la literatura examinada por Borck y Owings [2003] sugiere que en la distribución de las subvenciones gubernamentales predominan las consideraciones políticas, más que las de eficiencia. Por ejemplo, tiende a ir más dinero a los estados que tienen representantes en comités importantes del Congreso.

\textcolor{red}{\textbf{OJO -> Origen:} ROSEN y GARBER}
\textbf{Subvención de contrapartida con límite }
El coste para el donante de una subvención de contrapartida depende, en última instancia, del comportamiento del receptor. Si Smallville aumenta sustancialmente su consumo de G, las contribuciones del gobierno central serán bastante grandes, y viceversa. Para poner un límite al coste, el donante puede especificar alguna cantidad máxima que aportará. Una subvención de contrapartida con límite de este tipo se ilustra en la Figura 22.5. Como antes, antes de la subvención, la línea presupuestaria de Smallville es AB, y el equilibrio se encuentra en el punto E₁. Con la subvención de contrapartida con límite, la restricción presupuestaria es la línea quebrada ADF. La pendiente del segmento AD es menos un medio, reflejando la disposición de contrapartida de uno por uno. Pero después de cierto punto D, el donante deja de aportar un dólar por cada dólar. El coste de oportunidad para Smallville de una unidad de gasto público vuelve a ser 1 dólar, lo cual se refleja en la pendiente del segmento DF.

El nuevo equilibrio en E₃ implica un mayor consumo de G que bajo el statu quo, pero menor que bajo la subvención de contrapartida sin límite. El hecho de que la subvención se agote limita su capacidad para estimular el gasto en el bien público. Sin embargo, en algunos casos, el hecho de que tenga un límite puede ser irrelevante. Si el consumo deseado de G por parte de la comunidad implica un gasto inferior al límite máximo, la presencia del límite es irrelevante. En términos gráficos, si la nueva tangencia hubiera estado a lo largo del segmento AD de la Figura 22.5, no importaría que los puntos a lo largo de DR no estuvieran disponibles. Baker et al. [1999] realizaron un interesante estudio del impacto de pasar de un sistema de subvenciones de contrapartida sin límite a uno con límite en Canadá. Antes de la década de 1990, por cada dólar que una provincia canadiense gastaba en programas de asistencia social, el gobierno central aportaba otro dólar. Para contener los costes, en 1990 el gobierno central convirtió el programa en un sistema con límite en tres de las diez provincias. De forma coherente con lo expuesto en la Figura 22.5, el gasto en las tres provincias afectadas disminuyó en relación con las demás.

Subvención sin contrapartida Aquí el donante entrega una suma fija de dinero con la estipulación de que se gaste en el bien público. La Figura 22.6 representa una subvención sin contrapartida para comprar AH unidades de G. En cada nivel de renta de la comunidad, Smallville puede comprar ahora AH unidades más del bien público de las que podía comprar antes. Así, la nueva restricción presupuestaria se obtiene añadiendo una distancia horizontal AH a la restricción presupuestaria original AB. El resultado es la línea quebrada AHM.

Smallville maximiza la utilidad en el punto E₄. Obsérvese que, aunque el consumo del bien público aumenta de G₁ a G₄, la diferencia entre ambos es menor que la cantidad de la subvención, AH. Smallville ha cumplido la estipulación de gastar la totalidad de la subvención en G, pero, al mismo tiempo, ha reducido sus propios gastos en el bien público. Si el donante esperaba que los gastos aumentaran exactamente en AH, entonces la reacción de Smallville frustra estas expectativas. Resulta que la situación representada en la Figura 22.6 es una buena descripción de la realidad. Las comunidades utilizan a menudo una parte del dinero de las subvenciones condicionadas sin contrapartida para reducir sus propios impuestos. Según una estimación, por ejemplo, por cada dólar de ayuda educativa recibido por una comunidad, los impuestos locales disminuyen entre 30 y 70 centavos [Fisher y Papke, 2000, p. 157].

\textcolor{red}{\textbf{OJO -> Origen:} ROSEN y GARBER}
Subvenciones incondicionadas 
Obsérvese en la Figura 22.6 que la línea presupuestaria AHM parece casi como si hubiera sido creada dando a la comunidad una subvención de suma fija sin restricciones de AH dólares. Tales subvenciones incondicionadas se denominan a veces reparto de ingresos. Una subvención incondicionada habría conducido a una línea presupuestaria JM, que es simplemente el segmento MH prolongado hasta el eje vertical. Smallville resulta comportarse exactamente de la misma manera frente a la restricción AHM que como se habría comportado si se hubiera enfrentado a JM. En este caso concreto, por tanto, la subvención condicionada podría igualmente haber sido una subvención de suma fija sin restricciones. Intuitivamente, mientras la comunidad quiera consumir al menos una cantidad del bien público igual a la subvención, el hecho de que la subvención sea condicionada es irrelevante. En cambio, si la comunidad quisiera consumir una cantidad del bien público inferior a AH (si las curvas de indiferencia más altas fueran tangentes en algún punto a lo largo de JM a la izquierda de H), entonces la naturaleza condicionada de la subvención afectaría realmente al comportamiento.


¿Por qué debería el gobierno central dedicarse a conceder subvenciones incondicionadas a los estados y las localidades? La respuesta habitual es que tales subvenciones pueden igualar la distribución de la renta. La validez de este argumento no está clara. Incluso si un objetivo de la política pública es ayudar a las personas pobres, de ello no se deduce que la mejor manera de hacerlo sea ayudar a las comunidades pobres. Después de todo, lo más probable es que una comunidad con una renta media baja tenga algunos miembros relativamente ricos, y viceversa. Si el objetivo es ayudar a los pobres, ¿por qué no darles el dinero directamente?

Una posible explicación es que al gobierno central le preocupa especialmente que los pobres consuman una mayor cantidad del bien proporcionado públicamente. Un ejemplo importante es la educación. Este es un tipo de igualitarismo de bienes (Capítulo 12) aplicado a la producción del sector público. Sin embargo, como acabamos de demostrar, con las subvenciones incondicionadas no podemos saber con seguridad que todo el dinero terminará gastándose en el bien favorecido. (De hecho, lo mismo ocurre también con las subvenciones condicionadas).

\textbf{Medición de la necesidad} 
En cualquier caso, un programa redistributivo de subvenciones exige que el donante determine qué comunidades «necesitan» dinero y en qué cantidades. Las asignaciones federales se basan en fórmulas complicadas establecidas por el Congreso. La cantidad de dinero de subvenciones recibida por un estado depende de factores tales como la renta per cápita, el tamaño de su población urbana y la cantidad de sus recaudaciones del impuesto estatal sobre la renta. Las asignaciones a las localidades son funciones de estos factores económicos convencionales y también pueden depender de elementos tales como la etnia de la población.

Un factor importante para determinar cuánto recibe una comunidad del gobierno federal es su esfuerzo fiscal, normalmente definido como la ratio entre la recaudación tributaria y la capacidad fiscal. La idea es que las comunidades que se esfuerzan mucho por recaudar impuestos pero aun así no pueden financiar un nivel muy elevado de servicios públicos son merecedoras de recibir una subvención. Por desgracia, esta y otras medidas relacionadas pueden proporcionar poca o ninguna información sobre el verdadero esfuerzo de una comunidad. Supongamos que Smallville está en condiciones de exportar su carga fiscal en el sentido de que la incidencia de cualquier impuesto que establezca recae sobre personas de fuera. Entonces, un tipo impositivo elevado no nos dice nada sobre cuánto están sacrificando los miembros de la comunidad.

Más fundamentalmente, el enfoque del esfuerzo fiscal puede quedar totalmente desprovisto de significado debido al fenómeno de la capitalización. Consideremos dos municipios, Sodoma y Gomorra. Son idénticos excepto por el hecho de que Sodoma tiene un arroyo que proporciona agua prácticamente a coste cero. En Gomorra, por otro lado, es necesario cavar un pozo y bombear el agua.

Gomorra establece un impuesto sobre la propiedad para financiar la bomba de agua. Si hay un impuesto en Gomorra y ninguno en Sodoma, y las comunidades son por lo demás idénticas, ¿por qué debería vivir alguien en Gomorra? A medida que las personas migran a Sodoma, los valores de las propiedades aumentan allí (y disminuyen en Gomorra) hasta que no existe ninguna ventaja neta de vivir en una u otra comunidad. En resumen, los valores de las propiedades son más elevados en Sodoma para reflejar la presencia del arroyo.

Por las razones analizadas anteriormente, no esperamos que la ventaja esté necesariamente capitalizada al 100 % en los valores de las propiedades de Sodoma. No obstante, la capitalización compensa al menos parcialmente las diferencias entre los municipios. El mero hecho de que Gomorra establezca un impuesto no significa que se esté «esforzando más» que Sodoma, porque los habitantes de Sodoma ya han pagado por su agua mediante un precio más elevado por vivir allí. Concluimos que las medidas convencionales del esfuerzo fiscal pueden no ser muy significativas.

El efecto papel matamoscas

Nuestro análisis de las curvas de indiferencia de la comunidad deja abierta una cuestión fundamental: ¿de quién son esas curvas de indiferencia? Según la teoría del votante mediano (Capítulo 6), las preferencias son las del votante mediano de la comunidad. Los burócratas y los cargos electos desempeñan un papel pasivo en la aplicación de los deseos del votante mediano.

Una implicación directa de la regla del votante mediano es que un aumento de 1 dólar en la renta de la comunidad tiene exactamente el mismo impacto sobre el gasto público que la recepción de una subvención incondicionada de 1 dólar. En términos de la Figura 22.6, ambos acontecimientos generan desplazamientos paralelos idénticos hacia fuera de la línea presupuestaria inicial. Si los cambios de la línea presupuestaria son idénticos, los cambios en el gasto público también deben ser idénticos.

Se ha realizado una cantidad considerable de trabajo econométrico sobre los determinantes del gasto público local. (Véase Inman [2008] para una revisión). Contrariamente a lo que cabría esperar, prácticamente todos los estudios concluyen que un dólar recibido por la comunidad en forma de subvención da lugar a un mayor gasto público que un aumento de un dólar en la renta de la comunidad. En términos aproximados, las estimaciones sugieren que un dólar recibido como subvención genera 40 centavos de gasto público, mientras que un dólar adicional de renta privada aumenta el gasto público en solo 10 centavos. Este fenómeno ha sido denominado efecto papel matamoscas, porque el dinero parece quedarse pegado en el sector donde cae inicialmente.

Algunas explicaciones del efecto papel matamoscas se centran en el papel de los burócratas. Recuérdese del Capítulo 6 que algunos sostienen que los burócratas buscan maximizar el tamaño de sus presupuestos. Como maximizadores del presupuesto, los burócratas no tienen ningún incentivo para informar a los ciudadanos sobre el verdadero nivel de financiación mediante subvenciones de la comunidad. Al ocultar esta información, los burócratas pueden engañar a los ciudadanos para que voten a favor de un nivel de financiación mayor del que habría existido de otro modo. Según esta visión, el efecto papel matamoscas se produce porque los ciudadanos desconocen la verdadera restricción presupuestaria.






\subsection{Déficit}

\textcolor{red}{\textbf{Origen:} ICEX-CECO}
Otra fuente de financiación de las haciendas territoriales que se observa en muchos casos. es el endeudamiento.

Sin embargo. en un contexto de perfecta corresponsabiliclad fiscal (que se refiere al hecho de que las jurisdicciones financien una parte suslantiva de su gaslo mediante sus impuestos, y se analizará más adelante), sólo constituiría un medio para aplazar mayores impuestos. y existen razones por las que se debe limitar el uso de este instrumento a nivel descentralizado. En primer lugar, la interferencia con objetivos macroeconómicos. La emisión de deuda por parte de los gobiernos territoriales puede interferir en el diseño de la polilica de eslabilización del gobierno central.

Además, dado que los gobiernos terTitoriales no pueden monetizar la deuda. la política de gestión de su deuda exigiría que las emisiones es1én limitadas por las previsiones futuras de ingresos. De esta manera. podría haber incentivos a no emplear estos ingresos a satisfacer las preferencias de los ciudadanos.

Y. por último, surge el problema de riesgo moral. Una jurisdicción tendrá incentivos para llevar a cabo políticas fiscales expansivas. ya que se beneficiará del aumento de la demanda agregada por la poi ítica. pero, ��n muchas ocasiones, sin tener que hacer frente a los costes asociados a ésta. Es decir, no se traduce en un aumento del riesgo región, pero sí en un aumento del riesgo nacional.



\subsubsection{Restricción blanda:}

\textcolor{red}{\textbf{Origen:} ICEX-CECO}
JJJ.lf.JV. Restricción presupuestaria blanda 

Otra línea de investigación impo11ante dentro de las aportaciones de la segunda generación al estudio de la descentralización de los ingresos. se refiere a los incentivos perversos que pueden surgir en cuanto al sistema de financiación.

En un sistema descentralizado, el gobierno central puede imponer una restricción presupuestaria a los gobiernos locales. Ésta puede instrumentalizarse a través de la fi nanciación que provee el gobierno central a las haciendas tetTitoriales.

Se define así el término de restricción presupuestaria dura, que es aquella que debe cumplirse. El caso contrario se refiere a la restricción presupuestaria blanda, que KORNAI empleaba en 1979 para describir el comportamiento de las empresas públicas de economías socialistas que tenían en cuenta la posibilidad de ser rescatadas en el caso de caer en bancarrota.

El concepto de restricción presupuestaria blanda ha sido ampliado a través de varias aportaciones, empleándose dentro del marco del federalismo fiscal para el caso en el que los gobiernos locales pueden manipular la cantidad de fi nanciación que obtienen.  

Como se señaló anteriormente con la mutualización de riesgos. de un sistema de transferencias y del endeudamiento pueden surgir problemas de riesgo moral que dan lugar a soluciones ineficientes. Las haciendas territoriales tienen comportamientos estratégicos y tratarán de aumentar su gasto a través de una mayor financiación por parte del gobierno central. Las teorías de la segunda generación destacan dos elementos fundamentales de cara a evitar restricciones presupuestarias blandas: la credibilidad del gobierno central, y las i nstituciones.

Para analizar la credibilidad, podemos basarnos en las aportaciones de lnman de 2003. Ese autor considera un juego donde participan el gobierno central y los locales. En una primera etapa, el gobierno central indica que no habrá rescates.

Posteriormente. las haciendas evalúan la viabilidad de esta afirmación y toman sus decisiones presupuestarias según consideren que la afirmación es creíble o no. Si consideran que no es creíble. decidirán llevar a cabo un gasto por encima de su capacidad de financiación.

En la etapa final, el gobierno central debe decidir si cumple con su compromiso, o si decide incumplirlo y rescatar a los gobiernos locales. De esta forma, bajo determinadas condiciones, el gobierno central puede considerar que es más bcnelicioso rescatar a las jurisdicciones e incumplir su promesa.

Es importante comprender que. en este marco. el comportamiento perverso es endógeno al sistema. De esta manera. limitar de manera externa la discrecionalidad de los gobiernos (por ejemplo a tra,és de un marco institucional que defina unas reglas presupuestarias) puede reducir estos incentivos endógenos. Por esto. las instituciones políticas )' fiscales son el otro elemento fundamental a la hora de solucionar estos incentivos perversos endógenos al sistema. y evitar la aparición de restricciones presupuestaria'> blandas.

Las aportaciones dentro de esta línea de investigación tratan de comprender cuál es la motivación para aquellos comportamiento�� e��tratégicos que dan lugar a resultados ineficientes. ) por qué en ocasiones la promesa de no rescate del gobierno central nunca es creíble.

En cuanto a la motivación, Goodspeed señala que. si el gobierno central no rescata a una jurisdicción en banca1i-ota. surgen dos problemas. En primer lugar. se reducirá el bienestar social de esa jurisdicción. al verse reducida la provisión de los bienes y servicios públicos locales. Pero, ademá . . podrá haber consecuencias políticas en las elecciones para el gobierno central.

Por tanto, el gobierno central acudirá al rescate por motivos jurídicos o electorales. y también para asegurar que la provisión del bien público local no se vea reducida. 

Wildasin en su modelo de 1 997 analiza la cuestión de por qué en ocasiones la promesa de no rescate no será creíble. Considera un modelo en el que el comportamiento de un gobierno local tiene extcrnalidades sohre el resto de jurisdicciones.  

Señala que el hecho de no rescatar a una hacienda territorial, tendría efectos pe1judiciales para el resto de las jurisdicciones. Por esta razón. en el caso de algunas haciendas, el compromiso de no rescate por parte del gobierno central nunca será creíble.

Esto ocurrirá especialmente en el caso de las jurisdicciones de mayor tamaño. ya que en este caso las externalidades hacia otras haciendas serán mayores. Es en efecto una aplicación del fenómeno ·'too big /O fail" al caso de las Finanzas Públicas.

Con todo esto, surge la pregunta de qué tipo de instituciones económicas y políticas son las más adecuadas para minimizar los incentivos perversos de los gobiernos locales que dan lugar a restricciones presupuestarias blandas. Una respuesta lógica es establecer una mayor centralización. Sin embargo. en línea con el argumento del Leviatán, esto aumenti-i el poder de monopolio, aumentando las ineficiencias.

Sin embargo. como señala Oates en su ai1ículo de 2005. la nueva literatura señala algunas características que pueden ayudar a resolver este problema.

En primer lugar. es importante que haya acceso de las jurisdicciones a los mercados de crédito eficientes para garantizar la disciplina fiscal de los gobiernos locales. Un mal desempeño financiero de un gobierno local en este contexto, puede resultar en un acceso reducido a los mercados de crédito y en tipos de interés mayores para su financiación.

Además. los mercados de suelo eficientes también son deseables cuando hay movilidad de factores. De esta fónna. una deuda excesiva de una jurisdicción puede reducir los incentivos de adquirir propiedades en esa región. o de abandonarla en busca de una mejor gestión.

En esta línea Qian y Roland seí'ialaron en 1 998 que la competencia entre distintas jurisdicciones por atTaer factores, y los movimientos de capitales incentivan una mayor disciplina fiscal de los gobiernos locales.

En tercer lugar, Oates destaca un sistema de financiación impositiva eficiente que garantice ingresos suficientes a las haciendas territoriales, y que les pem1ita tomar decisiones valorando los costes y los beneficios de sus acciones. Así. el sistema de transferencias no debe ser la fuente de financiación principal de las haciendas, sino centrarse únicamente en sus funciones asignativas y redistributivas sin poder ser manipulado. 

Por último. Oates destaca la importancia de un sistema descentralizado en el que la legislación no suponga un poder excesivo para las jurisdicciones. El problema en este caso, es que, si los gobiernos locales tienen mayor poder que el gobierno central, puede producirse una ·'captura local'' por parte de éstas.

De hecho, Blanchard y Shleifer en el año 2000 compararon la descentralización de China y la de Rusia. Estos autores consideran que la de China ha sido más exitosa debido a que su gobierno central ha sido lo suficientemente fuerte como para evitar la captura local y frenar acciones a nivel local que podían debilitar el sistema.

Debemos señalar que, en el caso de países en desarrollo, cumplir con estas características se complica debido a que sus mercados están menos desan-ollados, y a que las instituciones y la legislación están más limitadas.

Para éstos. la literatura económica propone imponer de forma legal la obligación de equilibrio presupuestario para los gobiernos locales, limitando su deuda, y establecer una legislación centrada en las quiebras de entidades públicas que especifique el procedimiento a seguir ante crisis financieras.

Y es que, de acuerdo a Oates. a la hora de establecer las instituciones y de diseñar el sistema descentralizado. se debe tener en cuenta el contexto político, económico, cultural, y la tradición histórica.







\clearpage
\section{Conclusión} 


\subsection{Recapitulación}


\subsection{Extensiones y relación con otras partes del Temario}



\subsection{Opinión}

\subsubsection{Idea final (Salida o cierra)}

\clearpage
\pagenumbering{gobble}
\MyBibliography



\MyBibItem{DAWID y GATTI}{Chapter 2 - Agent-Based Macroeconomics}{2018}{https://doi.org/10.1016/bs.hescom.2018.02.006}


% ANEXOS
\clearpage










\end{document}